\documentclass[11pt]{article}
\usepackage[utf8]{inputenc}
\usepackage[a4paper,margin=1in]{geometry}
\usepackage{amsmath}
\usepackage{amssymb}
\usepackage{amsthm}
\usepackage{graphicx}
\usepackage{authblk}
\usepackage{abstract}
\usepackage{longtable}
\usepackage{booktabs}
\usepackage[font=small,labelfont=bf]{caption}
\usepackage{tabularx}
\usepackage{adjustbox}
\usepackage{xcolor}
\usepackage{listings}
\usepackage{enumitem}
\usepackage{float}
\usepackage{newfloat}
\renewcommand{\arraystretch}{1.22}
% keep floats close to the text that discusses them
\renewcommand{\topfraction}{0.92}
\renewcommand{\bottomfraction}{0.75}
\renewcommand{\textfraction}{0.06}
\renewcommand{\floatpagefraction}{0.75}
\setcounter{topnumber}{3}
\setcounter{bottomnumber}{2}
\setcounter{totalnumber}{5}

% ---- code & algorithm styling ------------------------------------------- %
\definecolor{codebg}{rgb}{0.972,0.972,0.972}
\definecolor{codekw}{rgb}{0.10,0.22,0.55}
\definecolor{codestr}{rgb}{0.62,0.13,0.13}
\definecolor{codecom}{rgb}{0.28,0.48,0.30}
\definecolor{codenum}{rgb}{0.55,0.55,0.55}
\definecolor{codeframe}{rgb}{0.78,0.80,0.83}

\lstdefinestyle{py}{
  language=Python,
  backgroundcolor=\color{codebg},
  basicstyle=\ttfamily\footnotesize,
  keywordstyle=\color{codekw}\bfseries,
  stringstyle=\color{codestr},
  commentstyle=\color{codecom}\itshape,
  numberstyle=\tiny\color{codenum},
  numbers=left, numbersep=8pt,
  showstringspaces=false,
  breaklines=true,
  frame=single, framerule=0.6pt, rulecolor=\color{codeframe},
  columns=fullflexible, keepspaces=true,
  aboveskip=8pt, belowskip=4pt,
  xleftmargin=14pt, xrightmargin=4pt,
}

\lstdefinestyle{pseudo}{
  mathescape=true, escapeinside={(*}{*)},
  basicstyle=\small\sffamily,
  keywordstyle=\color{codekw}\bfseries,
  morekeywords={for,while,do,if,else,then,return,repeat,until,Input,Output,
    foreach,to,with,probability,each,in,and,or,not,initialize,sample,select},
  numbers=left, numberstyle=\tiny\color{codenum}, numbersep=8pt,
  columns=fullflexible, keepspaces=true, xleftmargin=16pt,
  aboveskip=6pt, belowskip=2pt,
}

\lstdefinestyle{console}{
  backgroundcolor=\color{black!88},
  basicstyle=\ttfamily\scriptsize\color{white},
  numbers=none, frame=single, framerule=0pt,
  columns=fullflexible, keepspaces=true,
  breaklines=true, xleftmargin=10pt, aboveskip=8pt, belowskip=4pt,
}

% a properly numbered, caption-able "Algorithm" float
\DeclareFloatingEnvironment[fileext=loa,placement={!ht},name=Algorithm]{algorithm}

\usepackage{lmodern}
\usepackage[T1]{fontenc}
\usepackage{microtype}
\usepackage{mathtools}
\usepackage[round,authoryear]{natbib}
\bibpunct{(}{)}{;}{a}{,}{,}
\usepackage[colorlinks=true,linkcolor={blue!55!black},citecolor={green!45!black},urlcolor={blue!60!black}]{hyperref}
\DeclareMathOperator*{\argmax}{arg\,max}

% DOI helper (the bibliography uses \doi{...}); render as a clickable link.
\providecommand{\doi}[1]{\href{https://doi.org/#1}{\nolinkurl{doi:#1}}}

% lightweight theorem-like environments for the convergence results
\newtheorem{theorem}{Theorem}
\newtheorem{proposition}{Proposition}
\newtheorem{corollary}{Corollary}
\newtheorem{definition}{Definition}
\newtheorem{assumption}{Assumption}
\newtheorem{lemma}{Lemma}
\newtheorem{remark}{Remark}

% jednolite symbole:
\newcommand{\E}{\mathbb{E}}
\newcommand{\evidence}{\mathcal{E}} % evidence - żeby nie mylić z \E
\newcommand{\filt}{\mathcal{F}}     % filtration
% upright, well-spaced operator symbol: fixes "OpDefine" rendering
\newcommand{\Op}[1]{\mathrm{Op}_{\mathrm{#1}}}

\title{\textbf{When Does Design Thinking Work?\\A Convergence Theorem for the Iterate--Test Core}}
\author[1]{Paweł Kuraś}
\author[2]{Adrian Michalski}
\author[3]{Alicja Gerka}

\affil[1]{Institute of Telecommunications, AGH University of Science and Technology, Kraków, Poland}
\affil[2]{Faculty of Mathematics and Applied Sciences, Rzeszów University of Technology, Rzeszów, Poland}
\affil[3]{Department of Complex Systems, Rzeszów University of Technology, Rzeszów, Poland}

\begin{document}

\maketitle

\begin{abstract}
\noindent Design Thinking (DT) has resisted formal analysis because of its descriptive, intuitive nature. This paper gives a mathematical account of DT's convergent decision core and the conditions under which it provably works. Rather than formalizing wicked problems, we analyze the \emph{solution-space} dynamics \emph{within a single, well-framed epoch} (fixed feature space and utility), treating reframing as a nested wicked process our result locates but does not solve. We model DT as a stochastic walk on a state-space graph with divergent (entropy-increasing) and convergent (utility-maximizing) operators, casting selection and testing as a Bayesian decision process over a multi-criteria utility. Our central result is not the convergence theorem itself (a few lines from Doob's submartingale theorem) but the identification of \emph{sufficient, empirically testable and operationally interpretable conditions} --- bounded exploration, retention, creative generativity, approximately unbiased testing --- under which the Ideate--Prototype--Test cycle converges almost surely to a local optimum of expected utility. Generativity is required only in a weakened, decaying (non-summable) form, and we set out how it is measured; bounded test bias, slowly drifting criteria and incremental growth of the feature space degrade the guarantee continuously rather than destroying it; and under stakeholder disagreement the limit becomes set-valued, with a Pareto archive that provably stabilizes on a locally non-dominated set. We are careful about evidence: the mathematics \emph{proves} conditional convergence; a reproducible Monte-Carlo study \emph{demonstrates} internal consistency \emph{in silico} --- including two findings that cut against the model's own presentation, that the process is an order of magnitude slower than its stopping rule suggests, and that the value of the reframing criterion lies in robustness to reframing cost rather than in beating a well-tuned schedule, and that its one subjective input is safe to over- but not under-estimate; and whether real teams satisfy the assumptions \emph{remains untested}, for which we supply a powered measurement protocol, a \emph{pre-registration} of the validation study stating in advance what would count against the model, and a proof-of-concept run whose fragility to coder disagreement we quantify rather than gloss. We connect the model to formal design theories (FBS, C--K, co-evolution) and to Bayesian optimization, stating which of its rate results transfer to the discrete setting and which do not. The contribution is a falsifiable vocabulary for when and why disciplined iteration improves design outcomes --- and where it cannot.
\end{abstract}

\vspace{0.5em}
\noindent\textbf{Keywords:} Design Thinking; wicked problems; design process modelling; Bayesian decision theory; convergence; computational design support

\section{Introduction: The Paradox of Rigor and Ambiguity in Design Thinking}

Design Thinking (DT) is used across corporate innovation \citep{simon1996, brown2008}, public policy \citep{rittel1973} and healthcare, on the strength of a humanistic approach that integrates human needs, technological possibility and business viability \citep{plattner2011}. It is taught as a sequence of phases --- Empathize, Define, Ideate, Prototype, Test --- while being described in the same breath as ``non-linear'', ``iterative'' and ``fluid'' \citep{brown2008, uebernickel2020}. That duality poses a paradox: how can a method that is both structured and deliberately flexible be rigorously analyzed, predicted, or systematically improved?

Existing descriptions of DT are pedagogically invaluable but too informal to serve as a theoretical foundation \citep{liedtka2018, uebernickel2020}, which obstructs both its development as a discipline and the building of computational design support \citep{broberg2023}. Digital tools today mostly facilitate collaboration (virtual whiteboards) rather than navigating the problem and solution space, because such systems require formal, algorithmic descriptions of what they support \citep{fillingim2022, hermsen2015, majka2024}.

Faced with these challenges, this paper poses a focused research question: \emph{under what sufficient, empirically meaningful conditions does the iterative core of Design Thinking provably improve design outcomes, and where must it fail?} We answer it for the convergent decision core of DT within a single problem-framing epoch, and we are explicit throughout about the boundary between what the mathematics guarantees and what real design practice contributes.

\paragraph{Scope of the contribution.}
We stress at the outset what this paper is and is not. It is a \emph{normative, idealized} model of the decision-making scaffolding of DT --- not a descriptive account of the full richness of design practice, and not a claim that designers literally compute posteriors. Two boundaries are worth drawing sharply, because both bear on how the results should be read.

\emph{A well-structured subproblem, not wickedness in general.} Three levels are worth distinguishing. A \textbf{wicked problem} (Rittel \& Webber) has no definitive formulation and no stopping rule; an \textbf{ill-structured problem} lacks clarity about its structure but is worked within some frame; a \textbf{well-structured subproblem} has a fixed feature space $F$ and utility $U$. \emph{Our theorem addresses only the third} --- solution-space convergence given a fixed framing, a well-structured subproblem nested inside the genuinely wicked activity of \emph{problem framing}, in which problem and solution co-evolve \citep{rittel1973, dorstcross2001}. We do not claim to have formalized wicked problems; Section~\ref{sec:framing} states in full what is and is not claimed at the framing level, and what it costs to draw the epoch boundary wrongly.

\phantomsection\label{sec:scope-normative}\emph{A normative model, stated plainly.} Before anything else: \textbf{this is a normative model, not a descriptive one.} It says what a design team \emph{would} have to do for iteration to be guaranteed to improve its outcome. It does \emph{not} say that design teams do this, and we are not claiming that they do. Designers do not compute posteriors, do not represent their options as subsets of a feature space, and do not evaluate a stopping rule; nothing in this paper asserts otherwise. Three consequences follow, and we state them here rather than leaving them to be assembled from remarks scattered through the text.

\emph{(i) What the theory is about.} It is a theory of a \emph{decision procedure} --- one that a team can approximate, that a support tool can execute on a team's behalf (Section~\ref{sec:tools}), and against which actual practice can be compared. Its object is the procedure, not the practitioner.

\emph{(ii) What would count as evidence for it.} Not that teams behave as modelled --- they do not --- but that the \emph{conditional} holds: that variation in how closely a team approximates the four conditions predicts variation in its outcomes, in the direction and roughly the magnitude the model implies. That is a comparative claim, and it is testable; Section~\ref{sec:prereg} pre-registers the test.

\emph{(iii) What would falsify it.} A team measured to satisfy the assumptions whose leading concept nonetheless fails to improve over honest test--iterate cycles; or, comparatively, no association between measured assumption-satisfaction and outcome quality across a powered sample of epochs.

Two places where normative language can slide into descriptive language, and where we have tried not to let it: ``the process has non-negative drift'' is a statement about the modelled process, never an observed team, and the simulation's outputs describe the model's behaviour, never a team's.

\emph{Three kinds of claim, kept distinct.} We are careful throughout to separate (i) what the \emph{mathematics proves} --- conditional almost-sure convergence under idealized assumptions; (ii) what the \emph{simulation demonstrates} --- internal consistency of the model and the direction and size of each assumption's effect \emph{in silico}; and (iii) what \emph{remains to be tested empirically} --- whether real teams satisfy the assumptions and whether the predicted failure modes occur in practice, which we address only with a validation \emph{protocol} (Section~\ref{sec:roadmap}), not with data. Where we write ``violating an assumption does $X$,'' $X$ is always a simulation result unless explicitly attributed otherwise.

\noindent The model's purpose is threefold: to give design researchers a precise vocabulary for a process usually described only in prose, to isolate the assumptions that make disciplined iteration work, and to serve as a blueprint for computational design support. The main contributions are:
\begin{enumerate}
    \item A compact formal model of one framing epoch of DT that couples a graph-based state space with a Bayesian decision rule over a multi-criteria utility, positioned explicitly against existing formal design theories (function--behaviour--structure, C--K theory, co-evolution) and against Bayesian optimization (Section~\ref{sec:related}).
    \item A shift of emphasis from the (routine) convergence \emph{theorem} to its \emph{assumptions}: four sufficient, operationally interpretable behavioural conditions --- bounded exploration, retention, creative generativity, approximately unbiased testing --- stated as empirically testable, frequently violated propositions, with the consequences of each violation analysed theoretically and in simulation. Generativity is stated in a weakened, decaying form (non-summable rather than uniformly bounded below), and we set out how it is measured.
    \item An operational definition of convergence for a \emph{framed subproblem}, distinguishing the analyst-side optimum from an implementable, cost-calibrated stopping rule, and clarifying its relation to Simon's satisficing.
    \item Extensions covering the cases the single-utility theorem does not: bounded test bias, slowly drifting criteria, incremental growth of the feature space, and stakeholder disagreement, for which we prove set-valued convergence of a Pareto archive (Proposition~\ref{prop:pareto}).
    \item A pre-registered validation study (Section~\ref{sec:prereg}) stating the hypotheses, sampling and stopping rules, exclusions, analysis plan and --- in advance --- the results that would count against the model.
    \item A reproducible Monte-Carlo study that verifies internal consistency, characterizes landscape ruggedness, quantifies how bias and impoverished testing erode the guarantee, separates two notions of ``converged'' that differ by an order of magnitude, tests whether the reframing criterion actually outperforms over-exploiting and over-exploring, replicates the qualitative results across an independently generated \emph{family} of landscapes (Section~\ref{sec:landscapes}), and benchmarks the loop against random, greedy, $\varepsilon$-greedy and UCB baselines under a matched test budget (Section~\ref{sec:benchmarks}).
\end{enumerate}

We emphasize that the mathematics is in the service of design theory: its value is not rigour for its own sake but a falsifiable explanation of \emph{why} and \emph{when} iteration helps, and a map of the method's limits on genuinely wicked problems.

\section{Conceptual Foundations: From Wicked Problems to Design Process Models}

\subsection{The Nature of the Artificial and Wicked Problems}

Simon's ``sciences of the artificial'' distinguish natural science, which describes the world, from design, which asks \textit{how things ought to be} \citep{simon1996}. Against that backdrop, Rittel and Webber's \textbf{wicked problems} \citep{rittel1973, thomas2021, chen2022} are ``resistant'' rather than ``evil'', and four of their ten properties map onto properties of a search space: \emph{no definitive formulation} means problem and solution are co-defined, so the state-space graph is \emph{dynamic}; \emph{no stopping rule} and \emph{good-or-bad rather than true-or-false} evaluation mean search in a multimodal landscape $U(p)$ for a good-enough maximum \citep{wikipediaNDb, smith2025}; implementation as a costly \emph{one-shot operation} means evaluation is expensive and noisy. ``Wickedness'', so read, is search in a dynamic, rugged, high-dimensional space under a costly evaluation function --- the translation of a philosophical concept into the constraints our model rests on.

\subsection{A Critical Review of Dominant Process Models}

Two descriptive models dominate. The \textbf{Stanford d.school / HPI model} presents five or six phases that, while drawn linearly, are emphatically iterative: a team may return to any earlier phase when testing reveals a misread of user needs \citep{plattner2011, brenner2016, abrell2016}. The conditions and logic of those ``jumps'' are left ambiguous --- exactly what we aim to make precise. The \textbf{Double Diamond} \citep{designcouncilNDa} foregrounds instead the alternation of \emph{divergent} and \emph{convergent} thinking, which is useful here because both have direct mathematical analogues --- increasing versus decreasing the cardinality or entropy of the working sets --- that our operators (Section~\ref{sec:utility}) make precise \citep{hermsen2015, menning2018, designcouncilNDb}. Both models describe \textit{what} happens; neither states \textit{under what conditions} it works.

\subsection{Positioning Relative to Existing Formal Models of Design}
\label{sec:related}

Design is not a formalism-free field, and it would be a mistake to present this model as the first to bring rigour to it. \textbf{Shape grammars} formalize generation as rule-based rewriting over spatial vocabularies \citep{stiny1972}. The \textbf{function--behaviour--structure} ontology \citep{gero2004} decomposes designing into transformations among what an artefact is for, how it behaves, and how it is structured. \textbf{Axiomatic design} \citep{suh1990} prescribes independence and information axioms. \textbf{C--K theory} \citep{hatchuel2009} models designing as the coupled expansion of a concept space $C$ and a knowledge space $K$; its expansion/restriction maps almost one-to-one onto our divergent/convergent operators, and our Bayesian convergent operator can be read as one decision-theoretic instantiation of C--K restriction. The \textbf{co-evolution} model \citep{dorstcross2001, dorst1995} and \textbf{design-space exploration} \citep{woodbury2006} emphasize that problem and solution spaces are refined together. Table~\ref{tab:related} summarizes the comparison; Table~\ref{tab:frameworks} contrasts DT with TRIZ and mathematical programming. Those two are the informative extremes for our purposes. \textbf{TRIZ} \citep{altshuller1999} attacks the same difficulty --- a rugged space in which incremental search stalls --- with the opposite instrument: instead of iterating and testing, it resolves a stated \emph{contradiction} by analogy to a catalogue of inventive principles, which in our vocabulary is a way of \emph{jumping between basins without paying for the jump in tests}. It therefore complements rather than competes with the present model: TRIZ is a proposal for raising the generativity factor $\pi^{\text{gen}}$ and for making reframes cheaper, both of which our analysis treats as exogenous. It offers no convergence result and needs none, because it does not model iteration. \textbf{Mathematical programming} sits at the other extreme and is a genuine limiting case of our model: as evaluation becomes cheap and exact ($\sigma\to0$, $B\to0$) and the neighbourhood becomes the whole space, the stopping rule collapses to optimality and the local guarantee becomes global.

Two things distinguish the present contribution. First, none of these theories states a \emph{convergence result}: they describe the moves of designing, not the conditions under which iterating those moves provably improves an outcome. Second, once concept selection is modelled as choosing which noisy option to prototype next, the Ideate--Prototype--Test loop is formally a \textbf{Bayesian-optimization} / best-arm-identification problem --- sequential selection under a surrogate belief with costly, noisy evaluation \citep{shahriari2016, bubeck2012}. That literature offers far sharper \emph{rate} results than almost-sure convergence, and we borrow its lens: to say precisely what the DT loop is (a bandit over prototypes), to import its winner's-curse and posterior-consistency cautions, and to explain (Remark~\ref{rem:diversity}) why keeping several live prototypes has option value. Which of its rates do and do not carry over to our discrete, independent-belief setting is a question we defer to Section~\ref{sec:rates}, where we answer it rather than gesture at it.

\begin{table}[h!]
\centering
\caption{Positioning of the present model relative to existing formal design theories. ``Convergence result'' asks whether the theory states conditions under which iteration provably improves the outcome.}
\label{tab:related}
\footnotesize
\setlength{\tabcolsep}{5pt}
\begin{tabularx}{\textwidth}{@{}>{\raggedright\arraybackslash}p{0.155\textwidth} >{\raggedright\arraybackslash}X >{\centering\arraybackslash}p{0.10\textwidth} >{\raggedright\arraybackslash}X@{}}
\toprule
\textbf{Approach} & \textbf{What it formalizes} & \textbf{Conv.\ result?} & \textbf{Relation to this work} \\
\midrule
Shape grammars \citep{stiny1972} & Rule-based generation of form & No & A generative model for the divergent operator \\
FBS ontology \citep{gero2004} & Function/behaviour/structure transformations & No & Complementary process ontology; orthogonal to our decision core \\
Axiomatic design \citep{suh1990} & Design axioms (independence, information) & No & Prescriptive criteria; could shape the utility $U$ \\
C--K theory \citep{hatchuel2009} & Coupled concept/knowledge expansion & No & Our operators instantiate C--K expansion/restriction \\
Co-evolution \citep{dorstcross2001} & Joint refinement of problem \& solution & No & Motivates our framing-epoch scope and reframing resets \\
Bayesian optimization \citep{shahriari2016} & Sequential choice under a surrogate belief & Yes (rates) & The I--P--T loop \emph{is} such a process; we borrow its lens \\
\textbf{This work} & Convergent DT core within a framing epoch & \textbf{Yes (a.s.)} & Minimal conditions for provable improvement \\
\bottomrule
\end{tabularx}
\end{table}



\begin{table}[!tb]
\centering
\caption{Design Thinking contrasted with TRIZ and classical mathematical programming.}
\label{tab:frameworks}
\footnotesize
\setlength{\tabcolsep}{5pt}
\begin{tabularx}{\textwidth}{@{}l >{\raggedright\arraybackslash}X >{\raggedright\arraybackslash}X >{\raggedright\arraybackslash}X@{}}
\toprule
\textbf{Dimension} & \textbf{Design Thinking} & \textbf{TRIZ / SIT} & \textbf{Mathematical programming} \\
\midrule
Search space & Dynamic, co-defined with the problem & Fixed inventive principles / contradiction matrix & Fixed, given feasible set \\
Objective function & Multi-criteria, subjective, re-weighted & Implicit (resolve contradiction) & Explicit, fixed \\
Evaluation & Costly, noisy, human-in-the-loop tests & Analogical, largely deductive & Cheap, exact \\
Guarantee sought & ``Good-enough'' satisficing & Inventive resolution & Global/local optimum \\
Our model captures & Convergent core within a framing epoch & --- (comparison target) & Limiting special case (cheap, unbiased eval) \\
\bottomrule
\end{tabularx}
\end{table}

\section{A Mathematical Toolkit for Creative and Inferential Processes}
\label{sec:toolkit}

We need a language for both the design artifacts and the moves that produce them: set theory, graph theory, and abductive reasoning with Bayesian decision theory. Table~\ref{tab:notation} collects every symbol in one place for the interdisciplinary reader.

\begin{table}[h!]
\centering
\caption{Notation used throughout the paper. \emph{Subscript convention:} an unsubscripted symbol (e.g.\ $S_i$, $P_r$) denotes the space or set in general; the same symbol with a step subscript (e.g.\ $S_{i,k}$, $P_{r,k}$) denotes its concrete instance at iteration $k$.}
\label{tab:notation}
\footnotesize
\setlength{\tabcolsep}{6pt}
\begin{tabularx}{\textwidth}{@{}>{\raggedright\arraybackslash}p{0.34\textwidth} >{\raggedright\arraybackslash}X@{}}
\toprule
\textbf{Symbol} & \textbf{Meaning} \\
\midrule
$F$; $p\subseteq F$ & Feature space (atomic functionalities); a prototype is a subset of features \\
$P_r=2^F$ & Prototype space (all feature subsets); $P_{r,k}$ = materialized prototypes at step $k$ \\
$I,\ N,\ \Pi,\ S_i$ & Insight, need, point-of-view, and idea spaces; $\widetilde{\Pi}$ = abductive POV candidates \\
$v_k=(I_k,N_k,\Pi_{k},S_{i,k},P_{r,k})$ & Design state at step $k$; $G=(V,A)$ the state-space graph \\
$U(p)\in[0,1]$ & Multi-criteria utility, Eq.~\eqref{eq:utility}; $U_D,U_V,U_F$ its desirability/viability/feasibility parts \\
$\evidence$ & Test evidence (user feedback, feasibility/market assessments) \\
$\filt_k$ & Filtration: the $\sigma$-algebra of all idea generations and test outcomes up to step $k$ \\
$M_k^p=\E[U(p)\mid\filt_k]$ & Posterior-mean (expected) utility of $p$ given the evidence so far \\
$X_k=\max_{p\in P_{r,k}} M_k^p$ & Maximum expected utility (the submartingale of Theorem~\ref{thm:converge}) \\
$Y_k=\max_{j\le k}U(p_j^*)$ & Retained-best \emph{true} utility; $p_k^*$ = leading prototype at step $k$ \\
$\mathcal{N}(p)$ & Operator-reachable neighbourhood (Hamming ball of radius $r_{\max}$) \\
$b_k,\ B$ & Test bias and its bound; $\sigma$ = observation noise s.d. \\
$\Phi_k=U_{\max}-X_k$ & Lyapunov-type optimality gap \\
$\mu:S_i\to P_r$ & Materialization map (idea $\mapsto$ testable prototype), Section~\ref{sec:materialization} \\
$\delta_k$ & Generativity: $\Pr$(an improving neighbour is generated \emph{and} retained at step $k$) \\
$\mathrm{EI}_k$; $\epsilon,\tau$ & Expected improvement; stopping tolerance and belief-width threshold \\
$G_{\text{ref}},c_{\text{ref}},c_{\text{iter}}$ & Expected gain from a reframe; cost of a reframe; cost of one more cycle \\
$U^{(s)}$; $A_k$; $\succ$ & Stakeholder $s$'s utility; Pareto archive; the domination order \\
$L$; $\Lambda$ & Longest possible, and realized (random), strictly improving chain in the $\mathcal{N}$-graph \\
\bottomrule
\end{tabularx}
\end{table}

\subsection{Set Theory for Representing Design Artifacts}

The artifacts a design process accumulates are naturally sets \citep{bagaria2021, wikipediaNDa}.
The \textbf{insight space} $I$ collects raw observations, quotes and field data; the \textbf{need space} $N$ their synthesized interpretations; the \textbf{point-of-view space} $\Pi$ the problem formulations, each $\pi$ specifying for whom, what and why we design; the \textbf{idea space} $S_i$ the candidate solution concepts responding to a $\pi$; the \textbf{feature space} $F$ the atomic functionalities (a ``folding handle'', a ``one-touch button''); and the \textbf{prototype space} $P_r=2^F$ the materialized artefacts, each a subset $p\subseteq F$. Ideas become prototypes by a \emph{materialization} map $\mu:S_i\to P_r$. The distinction between $S_i$ and $P_{r,k}$ carries weight later: $S_i$ holds \emph{untested} proposals, which carry no posterior and are ranked in our implementation by a cheap proxy, whereas $P_{r,k}$ holds the \emph{tested} prototypes that do carry a posterior --- and it is over that tested set alone that $X_k$ and $\Op{Select}$ operate.

\paragraph{The materialization map is not innocuous.}
\label{sec:materialization}
A prototype is a lossy realization of an idea --- a paper mock-up of a service, a wizard-of-oz stand-in for an algorithm --- and it matters for the theory exactly where the loss lands. We distinguish three regimes, because they have different formal consequences.

\emph{(i) Faithful materialization.} $\mu$ is injective on the candidates the team entertains and $U(\mu(s))$ is the quantity the team cares about. This is the implicit assumption of Section~\ref{sec:proof} and the case our simulation implements: an idea \emph{is} its feature set.

\emph{(ii) Noisy materialization.} The build introduces a zero-mean fidelity error: the realized artefact is $\mu(s)$ but what the test measures is $U(\mu(s))+\eta$ with $\E[\eta\mid\filt_k]=0$. Nothing changes. The fidelity error is statistically indistinguishable from observation noise, so it is absorbed into $\varepsilon$ in Assumption~\ref{as:unbiased} and inflates $\sigma$ only; Theorem~\ref{thm:converge} applies verbatim, at the price of the slower, noisier convergence the $\sigma$-sweep quantifies (Table~\ref{tab:sens}).

\emph{(iii) Biased (systematically lossy) materialization.} Low-fidelity prototypes typically \emph{under}-realize their ideas: $\E[\eta\mid\filt_k]=-\beta_k$ with $\beta_k\in[0,B_\mu]$. This is a bounded, signed bias, so Proposition~\ref{prop:bias} covers it directly with $B_\mu$ in the role of $B$: the limit is $2B_\mu$-locally optimal. The interesting consequence is not the bound but its interaction with generativity. A lossy $\mu$ can convert an improving idea into a non-improving prototype, so it acts on $\pi^{\mu}$ in the factorization~\eqref{eq:delta-factor}: fidelity does not merely add noise, it \emph{lowers the generativity rate itself}. Two testable predictions follow --- that raising prototype fidelity should raise measured $\delta_k$ (not just reduce variance), and that a team whose ideas consistently outperform its prototypes is fidelity-limited rather than idea-limited, and should invest in building, not ideating. Finally, teams do form beliefs about \emph{untested} ideas by analogy; restricting posteriors to tested prototypes is a deliberate conservatism, and admitting reasoned priors over untested ideas is a natural extension.

\subsection{Graph Theory for Modeling the Process Structure}

The process is a directed, weighted graph $G=(V,A)$ \citep{mateosND, smith2024}, which lets the non-linearity be stated rather than asserted (Figure~\ref{fig:state_diagram}). A node $v\in V$ is a \emph{design state} $v_k=(I_k,N_k,\Pi_k,S_{i,k},P_{r,k})$, a snapshot of the project at step $k$; an edge $a\in A$ is an \emph{operator} (a DT phase) transforming one state into another, weighted by the probability of choosing it or by its cost in time and resources. A project is then a walk on $G$.

\begin{figure}[h!]
    \centering
    \includegraphics[width=0.48\linewidth]{Fig2.png}
    \caption{State-space representation of the Design Thinking process as a directed graph}
    \label{fig:state_diagram}
\end{figure}

\subsection{Abductive Reasoning and Bayesian Inference for Modeling Decisions}

\begin{itemize}
    \item \textbf{Abductive Reasoning:} The Empathize $\rightarrow$ Define transition is not a summary of data but an act of \textbf{abduction}, ``inference to the best explanation'' \citep{mallinson2019, menning2018}: from the insights $I_k$ the team generates plausible hypotheses about the underlying problem (Points of View $\Pi$), the form of inference Peirce associated with genuinely new explanatory hypotheses \citep{norman2013, wiltse2016}. Consistent with treating $\Op{Ideate}$ as a black box, we do not model \emph{how} framings are conceived; the prior $P(\pi)$ is read as the team's initial plausibility weighting (e.g.\ maximum-entropy over the candidate set $\widetilde\Pi$), refined by how well each framing accounts for $I_k$ via the likelihood $P(\pi\mid I_k,N_k)$, so no framing is privileged except through fit to the data.

    \item \textbf{Bayesian Inference:} Convergence on the ``best'' hypothesis is a \textbf{Bayesian decision process} \citep{theodoridisND, enkmann2025}: a prior $P(H)$, a likelihood $P(\evidence\mid H)$ from test data --- the evidence $\evidence$ being any operationalized signal (satisfaction or adoption for $U_D$, cost for $U_F$, ROI for $U_V$) --- and the usual posterior \citep{johnson2019, laskey2023}. Each hypothesis is ``prototype $p$ has utility $u$'', so the posterior over $u$ yields the expected utility $\E[U(p)\mid\filt_k]$ that drives selection (Figure~\ref{fig:bayesian_update}) \citep{elgamalND, howard1966}.
\end{itemize}

\begin{figure}[h!]
    \centering
    \includegraphics[width=0.78\linewidth]{Fig3.png}
    \caption{Simulation of the Test operator as a Bayesian update. \textbf{(a)} As unbiased test evidence accumulates, the posterior belief over a prototype's utility concentrates on the true value $u^*$. \textbf{(b)} The posterior mean is a martingale whose $\pm\sigma$ band contracts around $u^*$.}
    \label{fig:bayesian_update}
\end{figure}

\section{A Formal Model of the Design Thinking Process}
\label{sec:model}

\subsection{The Design Process as a Stochastic Walk on a State-Space Graph}

A project is a sequence of states $v_0,v_1,\dots,v_T$ on the graph $G=(V,A)$ of Section~\ref{sec:toolkit}. The transition from state $v_k$ to $v_{k+1}$ is the result of applying an operator $Op$ belonging to the set of operators corresponding to the DT phases,
\[
\{\Op{Empathize},\; \Op{Define},\; \Op{Ideate},\; \Op{Prototype},\; \Op{Test}\}.
\]
The choice of which operator to apply next is stochastic, governed by a transition kernel $T$ with $T_{ij}=\Pr(v_{k+1}=v_j\mid v_k=v_i)$. Two remarks keep this honest. First, $T$ describes the \emph{meta-level} choice of which operator to apply, whereas the convergence analysis of Section~\ref{sec:proof} concerns the finer dynamics \emph{within} the Ideate--Prototype--Test operators once selected; the theorem therefore does not describe the whole walk on $G$ but \emph{phases of concentrated iteration}, and $T$ enters only through one condition: \emph{within a framing epoch the I--P--T operators are applied infinitely often} (the epoch does not stall in Empathize/Define). If a team instead spends essentially all its time reframing, the premise fails and the result simply does not apply --- it is a statement about \emph{when} a team is iterating, not a claim that teams mostly iterate. Second, the backtracking entries of $T$ (Test back to Empathize/Define) implement the reframing resets of Section~\ref{sec:proof}; a full Markov analysis over framing epochs is left to future work. Figure~\ref{fig:dt_loop} shows a simplified view.



\begin{figure}[h!]
    \centering
    \includegraphics[width=0.62\linewidth]{Fig1.png}
    \caption{Process Flow of Design Thinking with Backtracking and Iteration Loops}
    \label{fig:dt_loop}
\end{figure}

\subsection{Formalizing Divergence and Convergence}

The Double Diamond's alternation of divergent and convergent thinking is formalized by operators that modify the design state.

\begin{itemize}
    \item \textbf{Divergent Operators ($\Op{Div}$):} These operators, corresponding to phases like \textit{Discover} (Empathize) and \textit{Develop} (Ideate), are modeled as functions that increase the cardinality and complexity of the relevant sets in the state tuple.
    \begin{itemize}
        \item $\Op{Empathize}: v_k \rightarrow v_{k+1}$, where the insight set grows by \emph{genuinely new} information: $I_{k+1}=I_k\cup\Delta I$ with $\Delta I\neq\varnothing$ and $\Delta I\cap I_k=\varnothing$. Non-redundancy matters: re-recording known facts is not divergence. We model only novelty, not the \emph{depth} of an insight --- that a single deep insight can outweigh many shallow ones is a limitation we return to in Section~\ref{sec:limitations}.
        \item $\Op{Ideate}: v_k \rightarrow v_{k+1}$, where the cardinality of the idea set increases: $|S_{i,k+1}| > |S_{i,k}|$. This represents a brainstorming session and the generation of new concepts.
        \item The effect of divergence is tracked by the Shannon entropy of the state (Figure~\ref{fig:entropy_increase}): pooling the features of all artifacts in a set and letting $\rho_k(j)$ be the normalized frequency of feature $j$ in the pool, $H(v_k)=-\sum_j\rho_k(j)\log_2\rho_k(j)$. That divergence raises $H$ is a \emph{modelling choice, not a theorem}: adding near-duplicate ideas that repeat existing features can \emph{lower} normalized entropy, so $|S_i|\!\uparrow$ does not imply $H\!\uparrow$ in general. The implication holds under a diversity condition --- each new idea contributes features not already saturated in the pool --- which our generative operator satisfies by construction, and Figure~\ref{fig:entropy_increase} should be read accordingly. Feature-weighted entropy or a KL divergence from a baseline would serve equally well; we use plain Shannon entropy because it is parameter-free, and no result depends on the choice.
    \end{itemize}





    \item \textbf{Convergent Operators ($\Op{Conv}$):} These operators, corresponding to phases like \textit{Define} and \textit{Deliver} (Prototype/Test), are modeled as functions that reduce the cardinality of sets by selecting a high-utility subset.
    \begin{itemize}
        \item $\Op{Define}: v_k \to v_{k+1}$. Abductive reasoning first generates a candidate set of Points of View $\widetilde{\Pi}$ from the insights $I_k$ and needs $N_k$; the operator then selects the single most promising one and contracts the POV component to that singleton:
\[
\pi^{*} = \argmax_{\pi \in \widetilde{\Pi}} P\!\big(\pi \mid I_k,N_k\big)\,U(\pi),
\qquad
v_{k+1} = \big(I_k,\, N_k,\, \{\pi^{*}\},\, S_{i,k},\, P_{r,k}\big),
\]
        so that $|\Pi_{k+1}| = 1 \le |\widetilde{\Pi}|$. This is an abductive and decision-making step: $\pi^{*}$ maximizes the product of the posterior probability and a POV-level utility $U(\pi)$ (the natural extension of the utility of Section~\ref{sec:utility} to problem framings).
        \item $\Op{Select}: v_k \to v_{k+1}$, where
\[
p^* \in \argmax_{p \in S_{i,k}} \,\E\!\big[U(p)\mid \filt_k\big],\qquad
v_{k+1} = \big(I_k,\, N_k,\, \Pi_{k},\, S_{i,k},\, P_{r,k} \cup \{p^*\}\big).
\]
This is a selection step based on decision theory. From the set of ideas $S_{i,k}$, the team selects the concept with the highest \textit{expected utility} $\E[U(p)\mid\filt_k]$ given the information $\filt_k$ available at step $k$, and materializes it as a prototype.

\emph{A worked instance, to fix the two roles $\Op{Select}$ plays.} Suppose $F$ contains \{folding handle, one-touch button, spill lid, temperature display\} and the team holds two untested ideas, $s_1=\{$handle, lid$\}$ and $s_2=\{$button, display$\}$, alongside one tested prototype $p_0=\{$lid$\}$ with posterior mean $0.55$. $\Op{Select}$ is invoked in two distinct capacities. \emph{As an acquisition rule} it decides \emph{what to build next} and must rank $s_1$ against $s_2$ --- objects with no posterior at all --- so in any implementation it necessarily falls back on a cheap prior proxy $\widehat U$ (an expert guess, an analogy, a heuristic score). \emph{As the definition of the incumbent} it ranges only over the tested set $\{p_0,\dots\}$, where genuine posterior means exist, and it is this second use that defines $X_k$ in Eq.~\eqref{eq:Xk}. The distinction matters because only the second enters the convergence argument: a poor proxy changes the \emph{order} in which the neighbourhood is explored (and hence the speed), never the quantity whose convergence we prove.
    \end{itemize}
\end{itemize}

\begin{figure}[!tb]
    \centering
    \includegraphics[width=0.68\linewidth]{Fig8.png}
    \caption{Simulation: the Shannon entropy $H(v_k)$ and cardinality $|S_i|$ of the design state rise under the divergent operator $\Op{Div}$ (Ideate) and fall under the convergent operator $\Op{Conv}$ (Select).}
    \label{fig:entropy_increase}
\end{figure}

\subsection{The Multi-Criteria Utility Function}
\label{sec:utility}

The core of convergent decision-making is the utility function. According to the framework proposed by IDEO, innovation must lie at the intersection of three key values: Desirability, Feasibility, and Viability.\citep{brown2008, uebernickel2020, designcouncilNDb} We therefore define the utility of a given prototype $p$ as a convex combination of these three criteria:

\begin{equation}
\label{eq:utility}
U(p) = w_d \, U_D(p) + w_v \, U_V(p) + w_f \, U_F(p),
\qquad w_d,w_v,w_f \ge 0, \quad w_d+w_v+w_f = 1 .
\end{equation}

Each component is taken to be bounded, $U_D,U_V,U_F \in [0,1]$; since the weights form a convex combination, the aggregate utility is likewise bounded, $U(p)\in[0,1]$. This boundedness is used in the convergence proof of Section~\ref{sec:proof}. We flag at once, to forestall a common misreading, that the convergence result does \emph{not} depend on this additive desirability--viability--feasibility form: Theorem~\ref{thm:converge} uses only that $U$ is \emph{bounded} and \emph{time-invariant within the epoch} (Section~\ref{sec:limitations}). The additive scalarization is a concrete, interpretable \emph{instance} used for the simulation and figures, not an ontological claim that design value is compensatory or cardinal; non-compensatory, lexicographic or merely ordinal utilities are all covered so long as they are bounded and stable within the epoch. Because the true criterion values are not known a priori but are estimated from noisy tests, decisions are made with respect to the \emph{expected} utility $\E[U(p)\mid\filt_k]$, the conditional mean given the evidence $\filt_k$ gathered up to step $k$.

For illustration, Figure~\ref{fig:utility_surface} shows a two-criteria cross-section of the 3D utility surface.

\begin{figure}[h!]
    \centering
    \includegraphics[width=0.82\linewidth]{Fig4.png}
    \caption{The utility landscape. \textbf{(a)} Enumerated prototypes in the desirability--feasibility plane coloured by aggregate utility $U=w_dU_D+w_vU_V+w_fU_F$, revealing the characteristic desirability--feasibility trade-off. \textbf{(b)} The distribution of utility \emph{values}. Ruggedness is quantified not by this histogram but by direct enumeration: $4$ local optima under the $1$-flip neighbourhood and a fitness--distance correlation of $-0.22$.}
    \label{fig:utility_surface}
\end{figure}

\begin{itemize}
    \item \textbf{Desirability ($U_D(p)$):} Measured based on user feedback obtained during testing. It can be modeled as a function of satisfaction ratings, adoption rates in a test, or other qualitative data transformed into a numerical scale.\citep{theodoridisND} This is a measure of whether people want the given solution.
    \item \textbf{Viability ($U_V(p)$):} Measured based on business analysis. It can be modeled as a projected return on investment (ROI), market size, alignment with company strategy, etc. This is a measure of whether the solution makes sense from a business perspective.
    \item \textbf{Feasibility ($U_F(p)$):} Measured based on technical assessment. It can be modeled as the estimated development cost, required resources, or an assessment of technical risk. This is a measure of whether the solution can be built using available technologies and resources.
\end{itemize}

In practice the weights shift as priorities change (early emphasis on desirability, later on feasibility and viability), and this matters for the theory: \emph{time-varying} weights make $U_k$ non-stationary, which --- like biased testing --- breaks the martingale structure underlying Theorem~\ref{thm:converge}. The result therefore holds \emph{within} a period of fixed weights; a large reweighting is treated exactly like a reframe, and Corollary~\ref{cor:drift} covers the intermediate regime of \emph{slowly} varying $U_k$. Table~\ref{tab:formal_definitions} is the resulting lexicon.





\begin{table}[h!]
\centering
\caption{Formal definitions of the Design Thinking phases and their operators. Each phase is an operator that transforms the design state $v_k=(I_k,N_k,\Pi_{k},S_{i,k},P_{r,k})$. Divergent phases (Empathize, Ideate) enlarge a knowledge set; convergent phases (Define, Prototype, Test) select or update under the utility/Bayesian criterion.}
\label{tab:formal_definitions}
\setlength{\tabcolsep}{5pt}
\footnotesize
\begin{tabularx}{\textwidth}{@{}l c c >{\raggedright\arraybackslash}X >{\raggedright\arraybackslash}X@{}}
\toprule
\textbf{Phase} & \textbf{Input} & \textbf{Output} & \textbf{Mathematical representation} & \textbf{Design goal} \\
\midrule
Empathize (divergent) &
$I_k$ &
$I_{k+1}$ &
$I_{k+1} = I_k \cup \{\text{new insights}\}$ &
Broaden \emph{and deepen} understanding of the user (here modelled only by breadth, $\lvert I\rvert\!\uparrow$) \\[4pt]

Define (convergent) &
$I_k, N_k$ &
$\Pi_{k+1}$ &
$\pi^{*} = \displaystyle\argmax_{\pi \in \widetilde{\Pi}} P(\pi \mid I_k,N_k)\,U(\pi)$ &
Select the most useful problem framing (posterior $\times$ utility) \\[4pt]

Ideate (divergent) &
$\Pi_{k}$ &
$S_{i,k+1}$ &
$S_{i,k+1} = S_{i,k} \cup \{\text{new ideas}\}$ &
Generate diverse candidate solutions ($\lvert S_i\rvert\!\uparrow$) \\[4pt]

Prototype (convergent) &
$S_{i,k}$ &
$P_{r,k+1}$ &
$p^* \in \displaystyle\argmax_{p \in S_{i,k}} \E[U(p)\mid\filt_k]$ &
Build the most promising testable artifact \\[4pt]

Test (convergent) &
$P_{r,k}$ &
$(I_{k+1}, U_D(p_k))$ &
$P(p_k \mid \evidence) = \dfrac{P(\evidence \mid p_k)\,P(p_k)}{P(\evidence)}$ &
Update beliefs from evidence $\evidence$ (Bayes' rule) \\
\bottomrule
\end{tabularx}

\smallskip
{\footnotesize\textit{Notation.} $I$ -- insight space (raw observations); $N$ -- need space (synthesized user needs); $\Pi$ -- point-of-view space (problem framings), with $\widetilde{\Pi}$ the abductively generated candidate set; $S_i$ -- idea space (candidate solution concepts); $P_r$ -- prototype space, where each prototype $p\subseteq F$ is a subset of the feature space $F$; $\evidence$ -- test evidence; $\filt_k$ -- the information (filtration) available at step $k$; $U(\cdot)$ -- the multi-criteria utility of Eq.~\eqref{eq:utility}; subscript $k$ -- iteration index.\par}
\end{table}

\section{Proof of Convergence for the Iterative Design Process}

\subsection{Defining Convergence for a Framed Subproblem}

Convergence here cannot mean finding a single objectively correct solution; on a wicked problem no such thing exists \citep{rittel1973, moustakas2023}. It must be defined in terms of optimization and satisfaction.





We distinguish an \emph{analyst-side} definition, used to state the theorem, from a \emph{team-side} rule that a real team (or a support tool) can actually evaluate, since a recurring objection to formal stopping criteria is that they require quantities no one can observe.

\begin{definition}[Convergence: analyst-side and team-side]
\label{def:convergence}
Let $p_k^* \in \argmax_{p\in P_{r,k}} \E[U(p)\mid\filt_k]$ be the leading prototype, $\epsilon>0$ a tolerance and $N\in\mathbb{N}$ a horizon.
\emph{(a, analyst-side)} The process has \textbf{$\epsilon$-converged} at step $k$ if the expected marginal improvement of the leading \emph{true} utility over the next $N$ iterations satisfies
\begin{equation}
\label{eq:eps}
\E\!\big[\,U(p_{k+N}^*) - U(p_k^*)\,\big] \;<\; \epsilon .
\end{equation}
This uses the true utility $U$ and is evaluable only by an omniscient analyst (or, in simulation, by the oracle landscape).
\emph{(b, team-side)} The \textbf{implementable} rule stops when the posterior \emph{expected improvement} of trying one more prototype falls below $\epsilon$ and the incumbent's belief is sufficiently concentrated:
\begin{equation}
\label{eq:ei}
\mathrm{EI}_k := \E\!\big[(U(p)-\hat{U}_k^{\max})^+ \mid \filt_k\big] < \epsilon
\quad\text{and}\quad \operatorname{sd}\!\big(U(p^*_k)\mid\filt_k\big) < \tau,
\end{equation}
where $\hat{U}_k^{\max}=\max_p \E[U(p)\mid\filt_k]$ and $\mathrm{EI}_k$ is the standard expected-improvement acquisition of Bayesian optimization \citep{shahriari2016}. Rule (b) uses only posterior quantities; under Assumptions~\ref{as:unbiased}--\ref{as:retention} it implies (a) in expectation.
\end{definition}

\noindent The contrast between (a) and (b) is not observable versus unobservable --- $\mathrm{EI}_k$ and $\operatorname{sd}(U(p^*_k)\mid\filt_k)$ are no more available to a team's unaided intuition than $U$ is. It is that (b) is \emph{computable from the team's own record} and (a) is not computable at all. Given a log of which concept was tested and what the test returned --- exactly the data the coding protocol of Section~\ref{sec:protocol} is designed to produce --- rule (b) is a deterministic function of that log, whereas rule (a) requires the true utility of prototypes that were never built. Three readings of (b) follow, and they are worth distinguishing: as a \emph{benchmark}, it says what an ideal observer with the team's own evidence would conclude; as \emph{tool logic}, it is what a support system computes on the team's behalf (Section~\ref{sec:tools}); and as an \emph{audit}, it can be evaluated retrospectively from a coded protocol to ask whether a team stopped too early or too late. None of the three asserts that anyone computes it in the moment.

\begin{remark}[Why expected improvement, and how to set $\epsilon$]
\label{rem:ei}
Two questions about rule~(b) deserve answers rather than convention.

\emph{Why EI and not a posterior probability of superiority?} Because in this setting EI \emph{is} the value of information. If the team's terminal payoff on shipping is the incumbent's utility, $\max_p M^p_k$, then the one-step expected gain from running one more test on $p$ is exactly $\E[(U(p)-\hat U_k^{\max})^{+}\mid\filt_k]$ --- the knowledge-gradient policy of \citet{frazier2008} reduces to EI whenever the tested option is the only one whose belief the test revises, which is the case for independent per-prototype beliefs. A posterior probability of superiority, $\Pr(U(p)>\hat U^{\max}_k\mid\filt_k)$, discards magnitude: it would keep a team iterating for a gain that is very likely but negligible, which is precisely the failure mode (terminal over-polishing) the rule is meant to prevent. The companion condition $\operatorname{sd}(U(p^*_k)\mid\filt_k)<\tau$ carries the ``confident enough'' leg that EI alone does not. When beliefs about prototypes are \emph{correlated} --- a realistic refinement, since testing one variant informs its neighbours --- EI is no longer the exact one-step value and the knowledge gradient is the proper generalization; we note this as the natural extension rather than adopting it here.

\emph{How should $\epsilon$ be chosen?} Not arbitrarily. Testing is costly, and if $c_{\text{test}}$ is the marginal cost of one more test--iterate cycle expressed in utility units, the economically correct threshold is
\begin{equation}
\label{eq:eps-cost}
\epsilon \;:=\; c_{\text{test}},
\end{equation}
i.e.\ stop when the expected gain from another cycle no longer covers what the cycle costs. This makes explicit something the bare tolerance hides: $\epsilon$ is not a statistical nuisance parameter but a \emph{price}, and the same $c_{\text{test}}$ reappears as $c_{\text{iter}}$ in the reframing criterion of Eq.~\eqref{eq:reframe}, so stopping and reframing are calibrated on one scale. Two consequences are worth stating. First, per-test cost does \emph{not} change Theorem~\ref{thm:converge}: the theorem is horizon-free, and cost enters only the stopping rule, so expensive testing changes \emph{when} a team stops, never \emph{whether} the process would converge if it continued. Second, it sharpens the satisficing reading: Simon's aspiration level is recovered as the cost-adjusted stopping threshold, and cheap testing (software A/B tests) rationally justifies a smaller $\epsilon$ --- more iteration --- than expensive testing (clinical trials, tooled manufacture).
\end{remark}

This plateauing behaviour is illustrated in Figure~\ref{fig:threshold} (rule~(a) with $N=5$, $\epsilon=0.01$). Both rules formalize the intuition that design stops not when perfection is reached but when the marginal gain no longer justifies the cost --- the ``good-enough'' criterion of wicked problems, which is precisely Simon's \emph{satisficing} under bounded rationality \citep{simon1956}. We stress that a local optimum and a satisficing solution are related but not identical: a local optimum satisfies rule~(b) with $\epsilon\to0$, whereas satisficing halts earlier, at the first state clearing an aspiration level; our theorem targets the former, and Eq.~\eqref{eq:ei} interpolates to the latter as $\epsilon$ and $\tau$ grow.

\begin{figure}[!tb]
    \centering
    \includegraphics[width=0.66\linewidth]{Fig6.png}
    \caption{Simulation of the stopping rule (Definition~\ref{def:convergence}). The expected $N$-step marginal utility improvement decays monotonically and crosses the threshold $\epsilon$ at iteration $k^*$, after which the process is $\epsilon$-converged.}
    \label{fig:threshold}
\end{figure}

\subsection{Expected Utility as a Lyapunov-Type Potential}
\label{sec:lyapunov}

We consider the iterative loop (Ideate $\rightarrow$ Prototype $\rightarrow$ Test $\rightarrow$ Ideate $\rightarrow \cdots$) as a discrete-time stochastic system adapted to the filtration $\{\filt_k\}_{k\ge0}$, where $\filt_k$ is the $\sigma$-algebra generated by all idea generations and test outcomes up to step $k$. The central quantity of the proof is the \textbf{maximum expected utility} over the materialized prototypes,
\begin{equation}
\label{eq:Xk}
X_k \;:=\; \max_{p \in P_{r,k}} \E\!\big[U(p)\mid\filt_k\big].
\end{equation}
In deterministic dynamical-systems theory a Lyapunov function is a non-negative scalar that \emph{decreases} along trajectories toward an equilibrium. The role of a Lyapunov function is played here, in stochastic form, by the non-negative \emph{potential} (or \emph{gap}) $\Phi_k := U_{\max} - X_k \ge 0$, where $U_{\max} := \max_{p\in 2^F} U(p)$. We will show that $X_k$ is a bounded submartingale; equivalently, $\Phi_k$ is a non-negative supermartingale that decreases in expectation, $\E[\Phi_{k+1}\mid\filt_k]\le\Phi_k$, and hence drives the system toward a stable set --- exactly the behavior a Lyapunov argument is designed to capture.

\subsection{Formal Proof of Convergence to a Local Optimum}
\label{sec:proof}

We work on a probability space with the filtration $\{\filt_k\}$ above. The proof rests on five assumptions, each stated not as a self-evident truth but as an \emph{empirically testable behavioural condition}. The scientific content of the theorem is that these conditions are \emph{sufficient} --- and operationally interpretable --- for disciplined iteration to provably improve the outcome; we do \emph{not} claim they are necessary, nor minimal in the class of all sufficient conditions, and a weaker sufficient set may well exist. Section~\ref{sec:sweep} measures in simulation what happens as each is relaxed.

\begin{assumption}[Finiteness within a project]
\label{as:finite}
Within a single design project only finitely many features $F$ are ever entertained, so the prototype space $P_r = 2^F$ is finite and, by Eq.~\eqref{eq:utility}, the utility is bounded, $U(p)\in[0,1]$, with $U_{\max}:=\max_{p\in 2^F} U(p)$. \emph{Testability / limits:} this idealizes away the open-ended generativity of design (truly novel features can emerge); it is the price of a convergence statement, and Section~\ref{sec:limitations} discusses generative-grammar alternatives that relax it.
\end{assumption}

\begin{assumption}[Bounded exploration]
\label{as:bounded}
$\Op{Ideate}$ proposes new prototypes only within a bounded neighbourhood $\mathcal{N}(p)$ of already-materialized prototypes (in our model, the Hamming ball of radius $r_{\max}$: prototypes differing from an explored one in at most $r_{\max}$ features). ``Local optimum'' means a state with no strictly better prototype in $\mathcal{N}(\cdot)$. \emph{Testability / limits:} teams overwhelmingly explore variations of what they already have; unbounded ``restarts'' that can reach the entire space $2^F$ are rare. This assumption is what makes \emph{local} optimality the genuine limit --- if $\mathcal{N}$ were the whole space, the reachable neighbourhood of every state would be $2^F$ and the argument below would instead force convergence to the \emph{global} optimum. We verify this dichotomy empirically in Section~\ref{sec:sweep}.
\end{assumption}

\begin{assumption}[Creative generativity and faithful selection]
\label{as:creative}
Let $A_k$ be the event that round $k$ generates an improving neighbour in $\mathcal{N}$ of the current incumbent \emph{and} the Prototype/Test operators rank it into incumbency, and let
\[
\delta_k \;:=\; \Pr\!\big(A_k \mid \filt_{k-1}\big)
\quad\text{on the event that the incumbent is not $\mathcal{N}$-optimal.}
\]
We assume the generativity sequence is \textbf{non-summable}:
\begin{equation}
\label{eq:nonsummable}
\sum_{k\ge1}\delta_k \;=\; \infty \qquad\text{a.s.\ on the event }\{\text{incumbent not $\mathcal{N}$-optimal}\}.
\end{equation}
\emph{Two readings, and why the weaker one matters.} The \emph{uniform} version, $\delta_k\ge\delta>0$, is the special case usually assumed; it is what one gets from Assumption~\ref{as:finite} as soon as every individual improving move has positive probability, because the incumbent then ranges over a finite set. But uniformity is a strong and unjustified claim about design teams: it says the chance of finding an improvement does \emph{not} decay as the team approaches a local peak, whereas diminishing returns are exactly what practitioners report. Condition~\eqref{eq:nonsummable} is what the proof actually needs, and it accommodates decay: $\delta_k\asymp c/k$ (or any $\delta_k$ with divergent sum) is admissible, while $\delta_k\asymp c/k^{2}$ is not. Substantively, \eqref{eq:nonsummable} says only that a team's generativity does not die out \emph{fast enough to be summable} --- a much weaker, and much more defensible, requirement than a constant $\delta$.
\emph{Testability / limits:} the condition couples generation with correct selection --- a genuinely better idea helps only if it is generated, prototyped, \emph{and} correctly identified; that coupling is exactly what biased testing (Assumption~\ref{as:unbiased}) can break. Section~\ref{sec:delta} discusses what $\delta_k$ means operationally and how it can be estimated.
\end{assumption}

\begin{assumption}[Approximately unbiased, consistent testing]
\label{as:unbiased}
Each test returns $y = U(p) + \varepsilon$ with $\E[\varepsilon\mid\filt_k]=b_k$ and $|b_k|\le B$. In the ideal case $B=0$ the evidence is conditionally unbiased, so the posterior-mean utility $M_k^{p}:=\E[U(p)\mid\filt_k]$ is a martingale (tower property). We further assume a retest schedule under which every \emph{retained} prototype is tested infinitely often (Algorithm~\ref{alg:ipt} realizes this by an explicit round-robin retest, Section~\ref{sec:code}), so each retained belief is \emph{consistent}: $M_k^{p}\to U(p)$ a.s.\ (standard for a well-specified conjugate update), which rules out a persistent winner's-curse gap between belief and truth. This is a property of the \emph{idealized} process; the finite-horizon Algorithm~\ref{alg:ipt} only approximates it, and Remark~\ref{rem:idealvsalgo} keeps the two apart. \emph{Testability / limits:} real user tests exhibit framing, confirmation and demand effects ($b_k\neq0$); Proposition~\ref{prop:bias} and Section~\ref{sec:sweep} quantify the damage.
\end{assumption}

\begin{assumption}[Retention]
\label{as:retention}
The materialized set never discards a prototype, $P_{r,k}\subseteq P_{r,k+1}$, so the best known solution is never lost. \emph{Testability / limits:} teams do abandon prototypes under time pressure; relaxing this weakens $X_k$ from a submartingale to a process that can drop, and is a natural direction for future work.
\end{assumption}

\subsubsection*{The status and the operational content of $\delta_k$}
\label{sec:delta}

\paragraph{Definition or empirical claim?} Both, and separating the two roles is what dissolves the appearance of circularity (``if a team does not converge, say its $\delta$ was too low''). \emph{Within the model}, \eqref{eq:nonsummable} is a \textbf{definition}: a practice is \emph{creatively generative} exactly when its improvement-generation rate is non-summable. So read, it explains nothing and is not meant to; it names the property the theorem needs. \emph{Outside the model}, $\delta_k$ is a \textbf{measurable quantity} attached to a particular team-and-practice, and pairing the definition with the measurement is what removes the circularity: ``this team converged because it was generative'' is empty, but ``this team's $\hat\delta_k$, estimated from its protocol before and independently of the outcome, was non-summable, and it converged'' is not. The theorem is then falsifiable in the usual conditional way --- a team with measurably non-summable generativity, bounded exploration, retention and near-unbiased testing that nonetheless fails to improve would refute it.

\paragraph{What $\delta_k$ is, and is not, a proxy for.} $\delta_k$ is deliberately \emph{not} a measure of creativity, novelty, or ideational fluency. It is the rate of \emph{retained improvement}: the joint probability that (i) ideation proposes a better neighbour, (ii) materialization realizes it faithfully (Section~\ref{sec:materialization}), and (iii) testing and selection recognise it. It therefore factorizes,
\begin{equation}
\label{eq:delta-factor}
\delta_k \;=\; \underbrace{\pi^{\text{gen}}_k}_{\text{ideation proposes it}}\cdot
              \underbrace{\pi^{\mu}_k}_{\text{materialization preserves it}}\cdot
              \underbrace{\pi^{\text{sel}}_k}_{\text{testing selects it}},
\end{equation}
and each factor is loaded by different practice variables: $\pi^{\text{gen}}$ by team diversity, facilitation quality, time pressure and psychological safety; $\pi^{\mu}$ by prototype fidelity; $\pi^{\text{sel}}$ by testing quality and bias. Reporting $\delta_k$ as one number would hide which of the three is the binding constraint, which is precisely the diagnostic value of the decomposition: a team with high fluency but low fidelity and a team with honest testing but a homogeneous group both show a low $\delta_k$ for opposite reasons and need opposite interventions.



\paragraph{Three estimators, in increasing order of demand.} (E1) From published corpora, the rate at which an evaluation is followed by renewed generation --- $39/52\approx0.75$ in \citet{sung2019} --- bounds $\pi^{\text{gen}}$, and hence $\delta$, \emph{from above} at no collection cost. (E2) From coded sessions, the ratio of rounds in which the retained best improved to rounds in which an improvement was still available --- a \emph{retrospective estimator of the realized improvement-opportunity rate}, not a measure of creativity, since the denominator is identified only \emph{post hoc} from the epoch's own final value; this is the estimator we apply to the simulator in Section~\ref{sec:rates-sim}, where it returns $\hat\delta=0.104$. (E3) From instrumented studies, the three factors separately, which requires blind post hoc ratings and is the only route to knowing \emph{which} factor binds. All three are written out as estimating equations in Section~\ref{sec:protocol}, with the data each needs and the direction in which each is biased; the point to carry forward is that whether $\hat\delta_k$ decays, and whether the decay is summable, is a measurement rather than a matter of opinion.

\paragraph{The division of labour we are keeping.} Assumption~\ref{as:creative} still does not \emph{explain} creativity, and we do not claim it should: our results concern the \emph{structural} conditions under which iteration provably helps and take the \emph{substantive} mechanism of idea generation as a black box. The honest consequence is that \textbf{the theorem's applicability to any given team rests on an empirical condition this paper does not validate}. A team whose generativity is summable --- homogeneous, rushed, unfacilitated --- falls outside the premise, and no convergence is claimed for it. What the formalism buys is the reclassification of ``keep generativity alive'' from a slogan into a measurement.

\paragraph{Are Assumptions~\ref{as:bounded} and~\ref{as:creative} compatible?} One might worry that bounded exploration (small $\mathcal{N}$) and generativity from every non-optimal state pull against each other on a rugged landscape. There is no tension, because both are defined \emph{relative to the same neighbourhood}: ``locally optimal'' in Assumption~\ref{as:creative} means $\mathcal{N}$-optimal, and \eqref{eq:nonsummable} is required only while the incumbent is not. When it is, no improving neighbour exists by definition --- and that is the state the theorem converges to. So $r_{\max}$ need not be large enough to escape every basin; it sets \emph{which} optima count as local, which makes it a tunable quality knob: exhaustive enumeration of our landscape gives $4$ local optima under the $1$-flip neighbourhood, $2$ at $r_{\max}=2$ and $1$ at $r_{\max}=3$ (Table~\ref{tab:rates}), while an unbounded radius collapses the count to one at the cost of leaving the regime that describes real teams. In short, the \emph{quality} of the guaranteed optimum is a function of the neighbourhood definition, $\text{quality guarantee}=f(\mathcal N)$: ``local optimum'' throughout means optimal within the operator-reachable $\mathcal N$, so it is in large part a \emph{parameter} of the model (the choice of $r_{\max}$) rather than an emergent property of Design Thinking, and this qualification should be read into every later occurrence of ``local optimum''.

\begin{theorem}[Almost-sure convergence of the design process]
\label{thm:converge}
Under Assumptions~\textup{\ref{as:finite}}, \textup{\ref{as:unbiased}} with $B=0$, and \textup{\ref{as:retention}}, the maximum expected utility $X_k$ defined in Eq.~\eqref{eq:Xk} is a bounded submartingale and therefore converges almost surely (and in $L^1$) to a finite limit $X_\infty$. If, in addition, Assumptions~\textup{\ref{as:bounded}} and \textup{\ref{as:creative}} hold, the limit is almost surely a \emph{local} optimum of expected utility with respect to the neighbourhood $\mathcal{N}$.
\end{theorem}

\begin{proof}
For each fixed prototype $p$, Assumption~\ref{as:unbiased} states that $M_k^{p}=\E[U(p)\mid\filt_k]$ is a martingale. Using retention ($P_{r,k}\subseteq P_{r,k+1}$), the monotonicity of the maximum, and the conditional Jensen inequality applied to the (convex) maximum functional,
\begin{align}
\E\!\big[X_{k+1}\mid\filt_k\big]
&= \E\!\Big[\max_{p\in P_{r,k+1}} M_{k+1}^{p}\;\Big|\;\filt_k\Big] \nonumber\\
&\ge \E\!\Big[\max_{p\in P_{r,k}} M_{k+1}^{p}\;\Big|\;\filt_k\Big]
   &&\text{(larger index set)}\nonumber\\
&\ge \max_{p\in P_{r,k}} \E\!\big[M_{k+1}^{p}\mid\filt_k\big]
   &&\text{(Jensen: } \max\E \le \E\max)\nonumber\\
&= \max_{p\in P_{r,k}} M_{k}^{p} \;=\; X_k
   &&\text{(martingale property).}
\end{align}
Hence $\{X_k\}$ is a submartingale. For boundedness we do not appeal to the observation model at all: $M_k^{p}=\E[U(p)\mid\filt_k]$ is the conditional expectation of the utility $U(p)$, which by Assumption~\ref{as:finite} is a (deterministic-valued) random variable taking values in $[0,1]$. A conditional expectation of a $[0,1]$-valued random variable lies in $[0,1]$ almost surely, so $M_k^{p}\in[0,1]$ and hence $X_k=\max_{p}M_k^{p}\in[0,1]$ --- regardless of how the belief is parameterized. In particular $X_k$ is $L^1$-bounded, $\sup_k\E[X_k^{+}]\le 1<\infty$. By Doob's submartingale convergence theorem a submartingale bounded in $L^1$ converges almost surely to an integrable limit $X_\infty$, and boundedness additionally yields convergence in $L^1$.
\end{proof}

\begin{remark}[The belief model, and why boundedness is automatic]
\label{rem:bounded}
The boundedness step is a statement about the \emph{ideal} posterior mean $M_k^{p}=\E[U(p)\mid\filt_k]$, and it holds for \emph{any} belief model in which $M_k^{p}$ is a genuine conditional expectation of the bounded utility: a $\mathrm{Beta}$/$\mathrm{Bernoulli}$ likelihood for $[0,1]$ outcomes (Section~\ref{sec:protocol}) delivers it exactly. The conjugate normal--normal update used in the simulation and in Algorithm~\ref{alg:ipt} is instead an \emph{unbounded Gaussian approximation} whose mean can in principle leave $[0,1]$ --- a Gaussian observation model $y=U(p)+\varepsilon$ does not by itself keep a precision-weighted posterior mean in the unit interval. We therefore treat the Gaussian belief as a computational device and take the \emph{metric projection} of its mean onto $[0,1]$, which enforces $M_k^{p}\in[0,1]$ as the ideal posterior requires; because observed utilities cluster well inside $[0,1]$ relative to the noise, the projection is almost never active and changes no reported number. The reading we adopt is that Theorem~\ref{thm:converge} governs the idealized bounded-belief process and the projected-Gaussian implementation approximates it, the two coinciding up to the (empirically negligible) projection.
\end{remark}

\begin{remark}[Where exploration lives, and a practice recommendation]
\label{rem:diversity}
The slack in the Jensen step, $\E[\max_p M_{k+1}^p\mid\filt_k]-\max_p\E[M_{k+1}^p\mid\filt_k]\ge0$, is not a technicality: it is exactly the \emph{value of diversification}. It is strictly positive when beliefs about competing prototypes are uncertain and imperfectly correlated, and vanishes when one prototype dominates with certainty --- so the exploration--exploitation trade-off of Bayesian optimization lives in this gap. The practical corollary is testable: \emph{carry more than one promising prototype for as long as beliefs about the leaders remain uncertain and weakly correlated, and collapse to a single incumbent only once the leader's advantage is statistically clear.} The number worth keeping live rises with posterior variance and falls with belief correlation; committing early forfeits the Jensen term, diversifying forever forfeits exploitation. The mathematics thus recovers the design-studio heuristic of ``keep two or three horses in the race,'' and locates the divergent operator's role as what makes the gap positive.
\end{remark}

This submartingale behavior --- a non-negative expected drift with locally fluctuating sample paths --- is illustrated \emph{in silico} (not confirmed empirically) in Figure~\ref{fig:submartingale}.

\begin{figure}[h!]
    \centering
    \includegraphics[width=0.95\linewidth]{Fig5.png}
    \caption{Monte-Carlo verification of theorem-consistent dynamics (ensemble of $N=250$ runs). This checks that the \emph{implemented model} behaves as the theorem predicts --- not that Design Thinking does (cf.\ Section~\ref{sec:simulation}). \textbf{(a)} The running maximum expected utility $X_k$ rises and saturates below the estimated global optimum $U^*$; the team's noisy belief estimate tracks the retained-best value. \textbf{(b)} The mean one-step increment $\E[\Delta X_k]$ is non-negative at every iteration, the defining condition of a submartingale. \textbf{(c)} The distribution of per-run limits $X_\infty$ shows that most runs settle at distinct \emph{local} optima strictly below $U^*$.}
    \label{fig:submartingale}
\end{figure}

\paragraph{From convergence to a local optimum.}
Theorem~\ref{thm:converge} guarantees that the submartingale $X_k$ stabilizes almost surely; we now characterize the limit. The posterior mean $X_k$ is itself only a submartingale (its sample paths may fluctuate as beliefs are updated), so to localize the limit we track the \emph{retained-best true utility}
\[
Y_k \;:=\; \max_{0\le j\le k}\, U\big(p_j^*\big),
\]
i.e.\ the highest \emph{true} utility actually achieved by an incumbent up to step $k$. Under retention (Assumption~\ref{as:retention}) $Y_k$ is non-decreasing, and by Assumption~\ref{as:finite} it takes values in the finite set $\{U(p):p\in 2^F\}$; a non-decreasing sequence with finitely many possible values is eventually constant, so there is a (random) step $K$ with $Y_k = Y_\infty$ for all $k\ge K$. It remains to show $Y_\infty$ is a local optimum with respect to $\mathcal{N}$. Suppose not, i.e.\ suppose the event $E=\{Y_k=Y_\infty \text{ for all }k\ge K,\ Y_\infty \text{ not } \mathcal{N}\text{-optimal}\}$ has positive probability. On $E$ the incumbent is non-optimal at every $k\ge K$, so Assumption~\ref{as:creative} applies at every such round and $\sum_k \Pr(A_k\mid\filt_{k-1})=\infty$. By L\'evy's conditional extension of the second Borel--Cantelli lemma \citep[Thm.~4.3.4]{durrett2019}, for an adapted sequence of events the divergence of $\sum_k\Pr(A_k\mid\filt_{k-1})$ implies $\Pr(A_k \text{ i.o.})=1$ on that event. Hence an improving neighbour is generated \emph{and} selected infinitely often on $E$, contradicting $Y_k\equiv Y_\infty$; so $\Pr(E)=0$ and $Y_\infty$ is almost surely $\mathcal{N}$-locally optimal. (Under the uniform reading $\delta_k\ge\delta>0$ the same conclusion follows from the elementary bound $(1-\delta)^n\to0$; the conditional Borel--Cantelli route is what buys the weaker, decaying-generativity condition~\eqref{eq:nonsummable}.) Finally, consistency of the incumbent's belief (Assumption~\ref{as:unbiased}) gives $X_\infty = U(p^*_\infty) = Y_\infty$ almost surely: in the limit the subjective maximum posterior mean and the true utility of the shipped prototype coincide, so the winner's curse does not survive. From step $K$ onward the expected marginal improvement in Eq.~\eqref{eq:eps} vanishes and the process is $\epsilon$-converged (Definition~\ref{def:convergence}).

\paragraph{Finite horizons vs.\ the limit.}
Under bounded exploration the almost-sure limit is a \emph{local} optimum, and that is also what one sees at practical horizons (Figure~\ref{fig:submartingale}c). Permitting unbounded global restarts makes Assumption~\ref{as:bounded} fail, $\mathcal{N}=2^F$, and the same contradiction argument then forces $Y_\infty$ to be the \emph{global} optimum --- but only as $k\to\infty$; at any finite horizon the team still observes a local optimum, because a random restart almost never lands in the global basin (Section~\ref{sec:sweep} confirms both). The headline is therefore a statement about the bounded-exploration regime that describes real teams, not an artifact of truncation.

\begin{remark}[The idealized process vs.\ the implemented algorithm]
\label{rem:idealvsalgo}
Theorem~\ref{thm:converge} concerns the \emph{idealized} process satisfying Assumptions~\ref{as:finite}--\ref{as:retention}, in which (i) beliefs are genuine bounded posteriors and (ii) every retained prototype is tested infinitely often. Algorithm~\ref{alg:ipt} is a \emph{finite-horizon heuristic} approximating that process, and two gaps are worth naming rather than hiding. First, its top-$r$ retest rule alone could let an older prototype fall out of the leading set and stop being tested, so we add an explicit round-robin retest (Section~\ref{sec:code}) that visits every materialized prototype in turn; over an infinite horizon this makes the implemented loop genuinely satisfy the infinite-retest clause of Assumption~\ref{as:unbiased} rather than merely approximate it. Second, the algorithm uses a projected-Gaussian belief (Remark~\ref{rem:bounded}) in place of the exact bounded posterior. Neither gap affects the theorem, which is a statement about the idealized limit; both are why the simulation is reported as a check of \emph{internal consistency} (Section~\ref{sec:simulation}), not as the theorem's proof.
\end{remark}

\subsection{Robustness to Biased Testing}
\label{sec:bias}
Assumption~\ref{as:unbiased} with $B=0$ is an idealization; the most important robustness question is what happens when it fails. A systematic bias $b_k\neq0$ breaks the martingale property, so $X_k$ is no longer a submartingale. The following statement bounds the damage.

\begin{proposition}[Graceful degradation under bounded bias]
\label{prop:bias}
Suppose Assumptions~\textup{\ref{as:finite}}, \textup{\ref{as:bounded}}, \textup{\ref{as:retention}} hold and the test bias is bounded, $|b_k|\le B$ (Assumption~\textup{\ref{as:unbiased}}). Write $\widehat M_k^{p}$ for the biased posterior mean and $M_k^{p}=\E[U(p)\mid\filt_k]$ for the unbiased one, and assume the standard conjugate-Gaussian update, under which $|\widehat M_k^{p}-M_k^{p}|\le B$ for every $p,k$ (the posterior mean is a precision-weighted average of observations, so a bias of at most $B$ in each observation shifts it by at most $B$). Let $\widehat X_k=\max_{p\in P_{r,k}}\widehat M_k^{p}$. Then:
\begin{enumerate}[label=(\roman*)]
\item $Z_k:=\widehat X_k + 2Bk$ is a submartingale, so $\widehat X_k$ has expected drift at least $-2B$ per step;
\item the retained-best true utility $Y_k$ converges almost surely, and its limit is \emph{$2B$-locally optimal}: no neighbour in $\mathcal{N}$ has true utility exceeding $Y_\infty + 2B$;
\item as $B\to0$ the guarantee of Theorem~\ref{thm:converge} is recovered.
\end{enumerate}
\end{proposition}

\begin{proof}
\emph{(i)} Fix $p$. The unbiased posterior mean $M_k^{p}$ is a martingale (Assumption~\ref{as:unbiased} with the tower property). The biased mean satisfies $\widehat M_k^{p}=M_k^{p}+\beta_k^{p}$ with $|\beta_k^{p}|\le B$ by hypothesis, so
\[
\E[\widehat M_{k+1}^{p}\mid\filt_k]=\E[M_{k+1}^{p}\mid\filt_k]+\E[\beta_{k+1}^{p}\mid\filt_k]\ge M_k^{p}-B=\widehat M_k^{p}-\beta_k^{p}-B\ge \widehat M_k^{p}-2B.
\]
Taking the maximum over the (retained, hence growing) set $P_{r,k}$ and using the same monotonicity--Jensen steps as in Theorem~\ref{thm:converge} gives $\E[\widehat X_{k+1}\mid\filt_k]\ge \widehat X_k-2B$. Now define $Z_k:=\widehat X_k+2Bk$. Then
\[
\E[Z_{k+1}\mid\filt_k]=\E[\widehat X_{k+1}\mid\filt_k]+2B(k+1)\ge(\widehat X_k-2B)+2B(k+1)=\widehat X_k+2Bk=Z_k,
\]
so $Z_k$ is a submartingale (and $\widehat X_k\ge0$ bounds it below). This is statement~(i); no constant is ``absorbed'' into $B$, which keeps its fixed meaning as the bias bound throughout.
\emph{(ii)} $Y_k=\max_{j\le k}U(p_j^*)$ is non-decreasing (retention) and lies in the finite set $\{U(p)\}$ (Assumption~\ref{as:finite}), hence is eventually constant, $Y_k=Y_\infty$ for $k\ge K$. Suppose a neighbour $q\in\mathcal{N}(p^*_\infty)$ had $U(q)>Y_\infty+2B$. Once the incumbent's belief is consistent (Assumption~\ref{as:unbiased}, up to the $\le B$ bias) we have $\widehat M^{p^*_\infty}\le U(p^*_\infty)+B\le Y_\infty+B$ and, after $q$ is generated and tested (which happens with positive probability by Assumption~\ref{as:creative} and infinitely-often testing), $\widehat M^{q}\ge U(q)-B>Y_\infty+B\ge \widehat M^{p^*_\infty}$, so a myopic team would switch its incumbent to $q$, contradicting $Y_k=Y_\infty$. Hence no such $q$ exists: the limit is $2B$-locally optimal.
\emph{(iii)} Immediate from (ii): the optimality gap $2B\to0$.
\end{proof}

\noindent The bound is deliberately loose; its role is qualitative --- bias causes \emph{bounded}, $O(B)$, continuous degradation rather than divergence (we do \emph{not} claim the expected final utility is a \emph{linear} function of $B$; the simulation shows a monotone, roughly linear trend, but that is an empirical regularity, not a theorem), and Section~\ref{sec:sweep} measures the actual effect (a confirmation bias of $b=0.2$ costs about $0.03$ of final utility). Its \emph{scope} is worth stating: the proposition covers \textbf{bounded additive bias in the observations}, $|b_k|\le B$, which includes confirmation bias as we model it (systematically over-rating the incumbent). It does \emph{not} cover distorted framing --- a mis-specified $U$ or $F$, treated in Section~\ref{sec:limitations} --- or selective attention that grows without bound.

\begin{corollary}[Tracking bound under slowly drifting utility]
\label{cor:drift}
Suppose the utility is not fixed but drifts, $U_k$, with bounded per-step change $\sup_p |U_{k+1}(p)-U_k(p)|\le D$ (e.g.\ criterion weights that are re-tuned gradually rather than reset). Then, running the argument of Proposition~\ref{prop:bias} with $D$ in the role of $B$, the retained-best utility obeys a \emph{tracking bound}: from some step on it remains within a $2D$-neighbourhood of a local optimum of the \emph{current} $U_k$. We deliberately do \emph{not} call this ``convergence'': when the target $U_k$ moves at every step there need be no fixed limit point, only a moving neighbourhood the process tracks, with tracking error proportional to the drift rate $D$ and not to elapsed time. Convergence to a fixed limit is recovered only once the drift ceases ($D\to0$, or $U_k$ eventually constant).
\end{corollary}

\begin{corollary}[Gradual co-evolution: feature accretion]
\label{cor:accretion}
Suppose the feature space grows monotonically, $F_0\subseteq F_1\subseteq\cdots$ (new features are admitted, none retired), with $U_k$ defined on $2^{F_k}$ and agreeing with $U_{k-1}$ on $2^{F_{k-1}}$ up to a per-step change of at most $D$. Then:
\begin{enumerate}[label=(\roman*)]
\item $X_k$ still converges almost surely. The ``larger index set'' step in the proof of Theorem~\ref{thm:converge} requires only $P_{r,k}\subseteq P_{r,k+1}$, and feature accretion \emph{enlarges} the ambient space without removing anything already materialized, so retention and boundedness are untouched; with $D>0$ the drift term of Proposition~\ref{prop:bias} applies with $D$ in place of $B$.
\item If the feature space stabilizes ($F_k=F_\infty$ for all $k\ge K_F$), the limit is almost surely a local optimum of $U_\infty$ with respect to $\mathcal{N}$ in $F_\infty$.
\item If it never stabilizes, $X_k$ converges but ``local optimum'' has no fixed referent: the neighbourhood is a moving target, and the process is best described as tracking rather than converging to a characterized limit.
\end{enumerate}
Feature \emph{retirement} ($F_{k+1}\not\supseteq F_k$) is a different matter: it can invalidate a retained prototype and so breaks Assumption~\ref{as:retention}.
\end{corollary}

\noindent Corollaries~\ref{cor:drift} and~\ref{cor:accretion} together answer a natural objection --- that real reframing is gradual (problem and solution co-evolve continuously \citep{dorstcross2001}) whereas our reframe is a discrete reset. The answer is that \emph{gradual} co-evolution decomposes into two moves the machinery already absorbs: slow re-weighting of criteria (Corollary~\ref{cor:drift}) and incremental admission of new features (Corollary~\ref{cor:accretion}). Neither ends the epoch; both merely add a tracking error proportional to the drift rate. What ends an epoch is the \emph{discontinuous} move --- retiring features, or resetting the criteria wholesale --- which is what Definition~\ref{def:reframe} isolates. So the epoch boundary is not an artefact of assuming framings are static; it is the point at which continuity of $(F,U)$ actually fails.

\noindent This matters for the ``how long is an epoch?'' worry. Real teams rarely hold $w_d,w_v,w_f$ exactly fixed; they nudge them as they learn. Corollary~\ref{cor:drift} says such \emph{gradual} re-weighting does not shatter the epoch into uselessly short pieces --- convergence degrades continuously in the drift rate --- whereas only a \emph{discrete, large} change of $F$ or $U$, a genuine reframe (Definition~\ref{def:reframe}), resets the process. So ``epoch'' means ``a period without a discontinuous reframe,'' not ``a period of literally constant $U$'' --- a far weaker and more realistic requirement, and whether real weight changes are gradual or abrupt is itself one of the empirical questions the protocol of Section~\ref{sec:roadmap} can settle. Consequently the theorem is \emph{not} a claim to ``solve'' wicked problems: within a fixed framing epoch, under bounded exploration, retention, and approximately unbiased testing, iterative improvement converges almost surely to a \emph{local} optimum of expected utility, degrading gracefully under bias --- matching the simulation, where $96.8\%$ of runs settle below the global optimum (Section~\ref{sec:checks}) --- and escaping a poor optimum requires reframing, not further iteration.

\subsection{Stakeholder Pluralism: Set-Valued Convergence}
\label{sec:pareto}

Design is rarely optimization against one agreed objective. If stakeholders hold a \emph{family} of utilities $\{U^{(s)}\}_{s\in\mathcal S}$ over a finite $\mathcal S$, the natural target is not a maximum but the Pareto set, and it is worth proving rather than asserting that the process converges to it. Write $U(p)=\big(U^{(1)}(p),\dots,U^{(|\mathcal S|)}(p)\big)\in[0,1]^{|\mathcal S|}$ and say $q$ \emph{dominates} $p$ ($q\succ p$) if $U^{(s)}(q)\ge U^{(s)}(p)$ for all $s$ with strict inequality for some $s$.

\paragraph{What replaces the argmax.} Two model elements change, and both are weaker than scalarization. The convergent operator becomes a \textbf{Pareto filter}: $\Op{Select}$ prototypes some candidate that is \emph{not dominated} by the current archive, with the choice among non-dominated candidates left free. And retention becomes \emph{Pareto retention}.

\begin{assumption}[Pareto retention and Pareto generativity]
\label{as:pareto}
The team maintains an \textbf{archive} $A_k\subseteq P_{r,k}$ of tested prototypes, and (a) a prototype leaves $A_k$ only when a newly archived prototype dominates it, never otherwise; (b) whenever some $p\in A_k$ has a neighbour $q\in\mathcal{N}(p)$ with $q\succ p$, such a neighbour is generated, materialized and archived at round $k$ with conditional probability $\delta_k$ satisfying the non-summability condition~\eqref{eq:nonsummable}.
\end{assumption}

\begin{proposition}[Convergence of the archive to a locally non-dominated set]
\label{prop:pareto}
Under Assumptions~\textup{\ref{as:finite}}, \textup{\ref{as:bounded}}, \textup{\ref{as:unbiased}} with $B=0$ holding coordinate-wise, and \textup{\ref{as:pareto}}:
\begin{enumerate}[label=(\roman*)]
\item \emph{(Coordinate-wise convergence.)} For each stakeholder $s$, $X_k^{(s)}:=\max_{p\in P_{r,k}}\E[U^{(s)}(p)\mid\filt_k]$ is a bounded submartingale and converges almost surely and in $L^1$.
\item \emph{(The archive stabilizes.)} The archive's value set $V_k:=\{U(p):p\in A_k\}\subseteq[0,1]^{|\mathcal S|}$ is eventually constant almost surely: there is a random $K$ with $V_k=V_\infty$ for all $k\ge K$.
\item \emph{(Local Pareto optimality.)} Almost surely, no $p\in A_\infty$ has a neighbour $q\in\mathcal{N}(p)$ with $q\succ p$; that is, $A_\infty$ is \emph{locally non-dominated} with respect to $\mathcal{N}$.
\item \emph{(Consistency with the scalar case.)} For $|\mathcal S|=1$, Assumption~\ref{as:pareto} reduces to Assumptions~\ref{as:creative}--\ref{as:retention} and (ii)--(iii) reduce to Theorem~\ref{thm:converge}. Fixed-weight scalarization, Eq.~\eqref{eq:utility}, is exactly this case.
\end{enumerate}
\end{proposition}

\begin{proof}
\emph{(i)} Fix $s$ and repeat the argument of Theorem~\ref{thm:converge} verbatim in that coordinate: $\E[U^{(s)}(p)\mid\filt_k]$ is a martingale for each $p$ by unbiasedness, $P_{r,k}$ grows (the \emph{materialized} set never shrinks even though the \emph{archive} may), and the monotonicity-plus-Jensen step gives $\E[X^{(s)}_{k+1}\mid\filt_k]\ge X^{(s)}_k$. Boundedness in $[0,1]$ and Doob's theorem finish it.

\emph{(ii)} Define the \emph{attainment set} $D_k:=\{u\in[0,1]^{|\mathcal S|}:\exists\,p\in A_k \text{ with } U(p)\succeq u\}$. By Assumption~\ref{as:pareto}(a) any $p$ removed at step $k$ is dominated by some $q\in A_{k+1}$, and nothing else is removed, so every $u\in D_k$ is still attained: $D_k\subseteq D_{k+1}$. By Assumption~\ref{as:finite} all values lie in the finite set $U(2^F)$, so $\{D_k\cap U(2^F)\}$ is a non-decreasing sequence of subsets of a finite set and is eventually constant. Since $V_k$ is precisely the set of maximal elements of $D_k\cap U(2^F)$, $V_k$ is eventually constant.

\emph{(iii)} Suppose the event $E$ that $V_k=V_\infty$ for $k\ge K$ while some $p\in A_\infty$ has a dominating neighbour has positive probability. On $E$, Assumption~\ref{as:pareto}(b) applies at every round $k\ge K$, so $\sum_k\Pr(\text{a dominating neighbour is archived}\mid\filt_{k-1})=\infty$ and, by L\'evy's conditional Borel--Cantelli lemma \citep[Thm.~4.3.4]{durrett2019}, such an event occurs almost surely on $E$. Archiving a dominator strictly enlarges $D_k$, contradicting the eventual constancy of $V_k$. Hence $\Pr(E)=0$.

\emph{(iv)} Immediate: with one coordinate, domination is strict improvement, the archive is the singleton incumbent, and (ii)--(iii) are the statements about $Y_k$ in Section~\ref{sec:proof}.
\end{proof}

\begin{remark}[Archive pruning does not conflict with retention]
\label{rem:pareto-monotone}
Removing dominated points from $A_k$ may look as if it violates the retention that Theorem~\ref{thm:converge} needs. It does not: the object required to grow is the \emph{attainment set} $D_k$ (the value vectors dominated by some archive member), and part~(ii) of the proof shows $D_k$ is monotone even while the archive itself sheds dominated points --- a removed point's value vector is still attained by the dominator that replaced it. Retention is thus preserved at the level of \emph{attained value}, which is all the submartingale argument uses; Proposition~\ref{prop:pareto} is built on this potential rather than on the raw archive.
\end{remark}

\paragraph{What this does and does not license.} Three clarifications, because ``convergence to the Pareto front'' is easy to over-claim.

\emph{(a) The result is invariant to the tie-break.} Proposition~\ref{prop:pareto} constrains only which moves are \emph{allowed} (Pareto-monotone replacement), not which non-dominated candidate is chosen next. A team may break ties by a scalarization, by rotating whose criterion leads, or by unstructured negotiation; the archive converges regardless. This is exactly why the model needs no aggregation rule: \textbf{convergence does not require the stakeholders to agree}. What the tie-break changes is \emph{which} locally non-dominated set is reached and how fast, not \emph{whether} one is reached.

\emph{(b) Locally non-dominated $\ne$ the whole local front}, and the gap is closable. Proposition~\ref{prop:pareto} says no archive member can be improved on by a neighbour; it does not say $A_\infty$ \emph{covers} the local Pareto front. The missing ingredient is diversity, and it can be supplied explicitly. Say generativity is \textbf{front-covering} if, whenever a locally non-dominated prototype $q$ is $\mathcal{N}$-reachable from the archive and no archive member attains its value vector, $q$ is generated and archived at round $k$ with conditional probability $\delta_k$ satisfying~\eqref{eq:nonsummable} --- that is, if generativity is spread \emph{across} the front rather than merely present somewhere on it. Under that strengthening the same conditional Borel--Cantelli argument applies to each missing front point in turn, and since the local front is finite (Assumption~\ref{as:finite}), $A_\infty$ almost surely contains a prototype attaining every value vector on the $\mathcal{N}$-local Pareto front. In practice front-covering generativity is demanding, and $\epsilon$-dominance archiving \citep{laumanns2002} is the standard relaxation: retaining one representative per $\epsilon$-box bounds the archive's size while guaranteeing coverage to within $\epsilon$, at the cost of converging to an $\epsilon$-approximate rather than exact front. The failure this analysis names is real and diagnosable: an archive that is locally non-dominated but clustered means the team's ``options'' all sit in one corner of the trade-off space, which a support tool can detect by measuring the spread of $A_k$ and which no amount of further iteration will repair.

\emph{(c) Choosing a point on the front remains political, and that is not a gap in the theory.} The theorem says the \emph{menu} stabilizes; it is silent on the selection from that menu, and deliberately so, because a decision-theoretic model has no standing to adjudicate whose utility should win. What the model does contribute is to make the boundary explicit: everything up to the front is evidential and can be delegated to a tool, and everything after it is a value judgement that cannot. Multi-attribute utility theory \citep{keeney1976} supplies elicitation methods for teams that \emph{do} want to aggregate --- trade-off weights, value functions, partial specification via dominance cones --- while reminding us that elicitation is costly and noisy, hardest exactly where stakeholders disagree.

\subsection{How Fast? What Transfers from Bayesian Optimization, and What Does Not}
\label{sec:rates}

Almost-sure convergence is a weak guarantee: it says the process gets there, not when. Section~\ref{sec:related} invoked Bayesian-optimization regret bounds as a sharper companion, and a fair objection is that those bounds were not connected to this model. We connect them here, and are explicit about which do not transfer.

\paragraph{What does not transfer.} The GP-UCB bound $R_T=O^{*}(\sqrt{T\gamma_T})$ \citep{srinivas2012}, and the simple-regret rate $\widetilde O(\sqrt{\gamma_T/T})$ we quoted, presuppose a \emph{continuous} domain, a Gaussian-process surrogate whose kernel is well specified, and beliefs that are \emph{correlated} across the domain --- correlation is what lets a finite number of evaluations say anything about an uncountable space, and $\gamma_T$ measures exactly that. Our model has none of these: the domain is finite, beliefs are independent across prototypes, and no kernel is posited. So the ``roughly $1/\sqrt{T}$'' reading applies to the \emph{continuous analogue} of the model discussed in Section~\ref{sec:limitations}, not to the discrete process the theorem is about, and we withdraw it as a description of the latter.

\paragraph{What does transfer.} In the finite, independent-belief setting the relevant literature is best-arm identification. Within one plateau the team faces a bandit over the reachable neighbourhood $\mathcal{N}$, and identifying an $\epsilon$-optimal arm with confidence $1-\eta$ requires on the order of
\begin{equation}
\label{eq:bai}
\sum_{q\in\mathcal N}\frac{\sigma^2}{\max(\Delta_q,\epsilon)^2}\,\log\!\frac{|\mathcal N|}{\eta}
\qquad\text{tests}, \qquad \Delta_q := U(p^*)-U(q),
\end{equation}
by standard action-elimination arguments \citep{evendar2006, bubeck2012}. Two design-relevant readings follow. First, the cost is governed by the \emph{inverse squared utility gaps}: separating two near-equally-good concepts is quadratically expensive, which is the formal content of the practitioner's experience that the last comparison is the hardest --- and a justification for the $\epsilon$ in Eq.~\eqref{eq:eps-cost}, since one should not pay to resolve differences smaller than the cost of resolving them. Second, the dependence on breadth is only \emph{logarithmic} in $|\mathcal N|$: generating more alternatives is cheap to adjudicate, generating \emph{closer} alternatives is not.

\paragraph{A finite-time bound the model does prove.} Under the uniform reading of Assumption~\ref{as:creative} an elementary bound is available.

\begin{proposition}[Finite-time bound under uniform generativity]
\label{prop:rate}
Suppose $\delta_k\ge\delta>0$ and let $L$ be the length of the longest strictly improving path in the neighbourhood graph $\big(2^F,\{(p,q):q\in\mathcal N(p),\,U(q)>U(p)\}\big)$, which is a DAG. Let $T_{\mathrm{loc}}$ be the first round at which the retained best is $\mathcal{N}$-locally optimal. Then $T_{\mathrm{loc}}$ is stochastically dominated by a sum of at most $L$ independent $\mathrm{Geometric}(\delta)$ variables, so
\[
\E[T_{\mathrm{loc}}]\;\le\;\frac{L}{\delta},
\qquad
\Pr\big(T_{\mathrm{loc}}>n\big)\;\le\;\exp\!\Big(-\tfrac{(n\delta-L)^2}{2n\delta}\Big)\ \ \text{for } n>L/\delta .
\]
\end{proposition}
\begin{proof}
While the incumbent is not $\mathcal{N}$-optimal, each round produces a strict improvement with probability at least $\delta$, so the waiting time for the next improvement is dominated by $\mathrm{Geometric}(\delta)$. Retention makes the sequence of incumbent utilities strictly increasing, so the incumbents trace a strictly improving path in the DAG, which has at most $L$ edges. Summing $L$ geometric waiting times gives the mean; the tail follows from a Chernoff bound on the number of successes in $n$ Bernoulli($\delta$) trials.
\end{proof}

\noindent The bound is \emph{valid but loose}: with $L=532$ at $r_{\max}=2$ and a measured $\hat\delta=0.104$ it gives $\approx5{,}100$ rounds against an observed $46$ (Section~\ref{sec:rates-sim}). The looseness is entirely attributable to $L$, a worst-case chain length over the whole space, and it is worth replacing it with the sharper statement --- which is a theorem, not a fitted regularity.

\begin{proposition}[Rate decomposition]
\label{prop:decomp}
Let $\Lambda$ be the number of strict improvements the incumbent makes before reaching an $\mathcal{N}$-local optimum, and $T_{\mathrm{loc}}=\sum_{i=1}^{\Lambda}G_i$ where $G_i$ is the number of rounds between the $(i-1)$st and $i$th improvement. If $\delta_k\ge\delta>0$ then $\E[G_i\mid\filt]\le 1/\delta$ for every $i$, $\Lambda$ is a stopping time for the improvement-indexed filtration with $\Lambda\le L$ almost surely, and Wald's identity gives
\begin{equation}
\label{eq:rate-decomp}
\E[T_{\mathrm{loc}}]\;\le\;\frac{\E[\Lambda]}{\delta},
\end{equation}
with equality when $\delta_k\equiv\delta$. Proposition~\ref{prop:rate} is the corollary obtained by bounding $\E[\Lambda]$ by $L$.
\end{proposition}
\begin{proof}
Conditionally on the past, each round improves with probability at least $\delta$ while the incumbent is not $\mathcal{N}$-optimal, so $G_i$ is stochastically dominated by $\mathrm{Geometric}(\delta)$ and $\E[G_i\mid\filt_{i-1}]\le1/\delta$. Retention makes incumbent utilities strictly increasing, so the incumbents trace a strictly improving path in the DAG and $\Lambda\le L$; in particular $\E[\Lambda]<\infty$. Wald's identity for a stopped sum of conditionally dominated increments then gives $\E[\sum_{i\le\Lambda}G_i]\le\E[\Lambda]\sup_i\E[G_i\mid\filt_{i-1}]\le\E[\Lambda]/\delta$.
\end{proof}

\noindent Eq.~\eqref{eq:rate-decomp} is the operative statement, and it earns its keep by \emph{separating} the two factors: $\E[\Lambda]$ is a property of the \emph{landscape and the starting point}, $\delta$ a property of the \emph{team}. Halving a team's generativity costs exactly as much time as doubling the depth of the improvement chain it must climb --- and, unlike $L$, both factors are estimable from a real protocol, $\Lambda$ by counting coded improvements and $\delta$ by the estimator of Section~\ref{sec:delta}. In our landscape $\E[\Lambda]=4.8$ against a worst case of $L=532$, which is why the bound of Proposition~\ref{prop:rate} is two orders of magnitude out.

\paragraph{Scaling.} The relevant size is not $|2^F|$ but $|\mathcal N|=\sum_{j\le r_{\max}}\binom{|F|}{j}$, which grows \emph{polynomially} in $|F|$ with degree $r_{\max}$ ($|\mathcal N|=105$ for $|F|=14$, $r_{\max}=2$, against $|2^F|=16{,}384$). Combined with the logarithmic dependence in Eq.~\eqref{eq:bai}, this predicts that enlarging a design problem degrades primarily the \emph{quality} of the local optimum --- more features means more basins --- and only mildly the \emph{time} to reach one. That is a falsifiable prediction: bigger projects should not iterate proportionally longer, they should settle proportionally further from the best available.

\subsection{When to Reframe: An Optimal-Stopping View}
\label{sec:reframe}

Before proceeding we fix a distinction the rest of the paper depends on, because two very different moves are easily conflated.
\begin{definition}[Within-frame move vs.\ reframe]
\label{def:reframe}
A \textbf{within-frame move} keeps both the feature space $F$ and the utility $U$ (equivalently, the Point of View $\pi$) fixed and changes only the prototype under consideration. This includes bounded mutation (Assumption~\ref{as:bounded}) \emph{and} an unbounded ``global restart'' that samples a fresh prototype from all of $2^F$: a restart enlarges the reachable neighbourhood $\mathcal{N}$ but does not touch $F$ or $U$, so it is still within-frame. A \textbf{reframe} changes $F$, $U$, or both --- a move to a different Point of View $\pi'$ that redefines what counts as a feature and what counts as good.
\end{definition}
\noindent The two are not degrees of the same thing. Within-frame convergence (Theorem~\ref{thm:converge}) concerns the former; the ``global restart'' ablation of Section~\ref{sec:sweep} is a within-frame stress test of Assumption~\ref{as:bounded} (it shows the a.s.\ limit becomes global while the finite-horizon outcome stays local), \emph{not} a model of reframing. Reframing, by contrast, resets the process to a new $(F,U)$ and is what the criterion below governs. With that fixed:

Theorem~\ref{thm:converge} says iteration converges within a frame; it is silent on the practically decisive question --- \emph{when should a team stop optimizing the current frame and reframe (change $F$ or $U$)?} This is the design counterpart of the exploration--exploitation dilemma of organizational learning \citep{march1991}, and neglecting it is precisely the ``myopia of learning'' \citep{levinthal1993}: a team that only exploits its current basin adapts efficiently to a local peak of a rugged landscape while missing higher peaks elsewhere \citep{levinthal1997}. We give a simple, decision-theoretic criterion.

Let the team, at step $k$, hold the incumbent $p_k^*$ and its belief. Continuing within the frame yields, in expectation, the within-neighbourhood expected improvement $\mathrm{EI}_k$ of Eq.~\eqref{eq:ei} per round, at a marginal cost $c_{\text{iter}}$ (time, testing). A reframe --- a jump to a new frame $(F',U')$ along a backtracking edge --- offers an expected gain $G_{\text{ref}}$ (the anticipated value of a better basin, net of the value already secured) at a one-off cost $c_{\text{ref}}$ (re-empathizing, rebuilding context, discarding sunk work).

\begin{proposition}[Reframing criterion]
\label{prop:reframe}
A myopically rational team should reframe rather than continue iterating exactly when the expected value of one more within-frame cycle falls below that of a reframe:
\begin{equation}
\label{eq:reframe}
\mathrm{EI}_k \;-\; c_{\text{iter}} \;<\; G_{\text{ref}} \;-\; c_{\text{ref}} .
\end{equation}
Because $\mathrm{EI}_k$ is non-increasing along a converging epoch (it $\to 0$ as $X_k\to X_\infty$), the left-hand side decreases monotonically, so there is a well-defined first step at which reframing becomes preferable; before it, iterate; after it, reframe.
\end{proposition}

\noindent The criterion is a \emph{heuristic}, one-step look-ahead policy --- not a full optimal-stopping solution, and not a corollary of Theorem~\ref{thm:converge}, with which it has a different epistemic status (Theorem~\ref{thm:converge} is a proved limit statement; Eq.~\eqref{eq:reframe} is a normative decision rule we test in simulation, Section~\ref{sec:reframe-sim}). We resolve rather than leave open the tension between calling it ``myopically rational'' and describing it as a meta-bandit. Both are true of different objects. Eq.~\eqref{eq:reframe} is the \emph{one-step index policy} for the meta-bandit whose arms are frames; the meta-bandit is the decision problem, the inequality is a particular (greedy) policy for it, and the optimal policy --- a Gittins-style index or a finite-horizon dynamic program over frames --- is not solved here. What we can do is test whether the greedy policy is good enough to be worth stating, and Section~\ref{sec:reframe-sim} does exactly that against three baselines. With that said, the criterion makes three things precise. First, reframing is worthwhile only once within-frame gains have thinned ($\mathrm{EI}_k$ small), which is why premature reframing (``idea-hopping'') and terminal over-optimization (polishing a doomed concept) are \emph{both} failures of the same inequality. Second, the decision hinges on quantities the team can at least estimate rather than on the unobservable true global optimum. Third, it locates DT's characteristic move, the deliberate reframe, as the method's remedy for myopia.

\paragraph{What is derived and what is subjective.} Which side of Eq.~\eqref{eq:reframe} the model determines matters, since the criterion could otherwise read as ``reframe when you feel like it''. The \emph{left}-hand side is derived: $\mathrm{EI}_k$ is computable from the team's own posteriors --- the acquisition function a support tool already maintains --- and $c_{\text{ref}}$ is estimable from a team's history. The term $G_{\text{ref}}$ is \emph{not} derivable from within the current frame (that is what makes framing wicked), but it is bounded above by the optimality shortfall $U^{\star}_{\text{global}}-Y_k$ and estimable as the posterior mean of an as-yet-unentered frame's ceiling under a hierarchical prior over frame maxima, updated from cross-project experience and from weak signals such as recurring complaints the current frame cannot address. That makes the structure a \textbf{bandit over frames}: each frame is an arm whose payoff is its ceiling, $\mathrm{EI}_k$ is the value of pulling the current arm again, and $G_{\text{ref}}-c_{\text{ref}}$ the cost-adjusted value of switching. The criterion is therefore \emph{informational rather than decision-determining}: it computes the exploitation value exactly and locates the one subjective input precisely, so disagreement about it can be argued in the open. In organizational-learning terms \citep{march1991} within-frame iteration is exploitation and reframing is exploration, and the ``myopia of learning'' \citep{levinthal1993, levinthal1997} is what Eq.~\eqref{eq:reframe} is built to detect --- here in a form specific to designing, in which exploration is not a random jump but the disciplined reframe of the Double Diamond's first diamond.

\subsubsection*{Estimating $G_{\text{ref}}$}
\label{sec:gref}

$G_{\text{ref}}$ carries the whole subjective burden of Eq.~\eqref{eq:reframe}, so leaving it as ``the anticipated value of a better basin'' would leave the criterion unusable. It is estimable, and the estimator is worth writing out.

\paragraph{The estimator.} Let $V_f$ denote the outcome a team of this capability attains in a frame $f$, and place a Normal--Normal hierarchical prior over it:
\begin{equation}
\label{eq:gref-prior}
V_f \mid \theta \sim \mathcal N(\theta,\varsigma^2),
\qquad \theta \sim \mathcal N(m_0,\tau_0^2),
\end{equation}
where $\theta$ is the team-and-domain level from which frames are drawn, $\varsigma^2$ the frame-to-frame spread, and $(m_0,\tau_0^2)$ the cross-project prior. After visiting frames with attained outcomes $v_1,\dots,v_n$ the posterior is the usual precision-weighted update, $\tau_n^{-2}=\tau_0^{-2}+n\varsigma^{-2}$ and $m_n=\tau_n^2(\tau_0^{-2}m_0+\varsigma^{-2}\sum_i v_i)$, and the criterion's gain term is the \emph{frame-level expected improvement}
\begin{equation}
\label{eq:gref}
G_{\text{ref}} \;=\; \E\big[(V_{\text{new}}-Y_k)^{+}\big]
\;=\; (m_n-Y_k)\,\Phi(z) + \varsigma_n\,\phi(z),
\qquad z=\frac{m_n-Y_k}{\varsigma_n},\quad \varsigma_n^2=\tau_n^2+\varsigma^2 .
\end{equation}
The symmetry is not decorative: $G_{\text{ref}}$ is an expected improvement one level up, over frames rather than prototypes, which is what makes the meta-bandit reading operational rather than metaphorical. Three inputs, three sources: $m_0,\tau_0^2$ from the team's own history or a domain benchmark (``how good do our projects usually end up?''); $\varsigma^2$ from the spread of those same outcomes; $Y_k$ from the current epoch's coded record. All three are quantities a team could record without instrumentation it does not already have, and $G_{\text{ref}}$ is in any case bounded above by the optimality shortfall $U^{\star}_{\text{global}}-Y_k$.

\paragraph{A structural warning about which prior to use.} Eq.~\eqref{eq:gref} prices a reframe as \emph{one draw} from the outcome distribution. That is an under-estimate, and predictably so: after reframing the team does not receive $V_{\text{new}}$ and stop, it gets to \emph{iterate} in the new frame. The reframe is an option on a whole epoch, not a lottery ticket. Worse, if $(m_n,\varsigma_n)$ are estimated from \emph{attained} outcomes while $Y_k$ is a maximum over frames already visited, the comparison is between an average and a maximum and $G_{\text{ref}}$ is biased down a second time. The practical correction is to centre the prior on the \emph{good} outcomes of past projects rather than the typical ones --- an extreme-value rather than a mean prior --- and Section~\ref{sec:gref-sim} measures what happens when one does not. Note that this diagnosis also explains why charging the reframe an extra ramp-up cost did not help in Section~\ref{sec:reframe-sim}: the myopia of Eq.~\eqref{eq:reframe} is on the \emph{benefit} side, so a correction on the cost side moves the wrong way.

\subsection{What the Model Says --- and Does Not Say --- About Framing}
\label{sec:framing}

Because the theorem is explicitly about a \emph{framed} subproblem, the natural question is what, if anything, the paper claims about the framing level itself. The answer has three parts, and it is worth stating them together rather than leaving the reader to infer the boundary from disclaimers.

\paragraph{What is claimed.} Three things, all of them modest. \emph{(a)} Framing is where the guarantee ends: iteration within a frame converges to a \emph{local} optimum of that frame's utility, so a poor frame cannot be escaped by more iteration --- only by a reframe (Section~\ref{sec:reframe}). \emph{(b)} The \emph{decision} to reframe, unlike the framing itself, is amenable to decision theory: Eq.~\eqref{eq:reframe} makes the exploitation side exactly computable and isolates the one input that is not (Section~\ref{sec:gref}). \emph{(c)} Framing changes that are \emph{gradual} --- slow re-weighting of criteria, incremental admission of features --- are not epoch-ending at all; Corollaries~\ref{cor:drift} and~\ref{cor:accretion} absorb them with a tracking error proportional to the rate of change. Only a discontinuous change of $(F,U)$ resets the process.

\paragraph{What is not claimed.} We do not model how a framing is \emph{generated}: abduction from insights to a Point of View is treated as a black box exactly as ideation is, and the model has nothing to say about what makes one framing more insightful than another. We do not claim convergence \emph{across} frames --- Eq.~\eqref{eq:reframe} is a greedy policy for the meta-bandit, not a solution to it, and no result here says a team that keeps reframing eventually finds the best available frame. And we do not claim the frame sequence is Markov, exchangeable, or in any way statistically well behaved; the meta-level analysis of Section~\ref{sec:reframe} assumes only that a team can price its own alternatives.

\paragraph{The boundary problem, and what mis-drawing it costs.} A fair objection is that real design is not partitioned into epochs at all, so an analysis conditioned on epochs may be conditioning on an artefact. We accept the premise and answer it in two ways. First, Corollary~\ref{cor:accretion} substantially weakens the need for a sharp boundary: since gradual drift and feature accretion do \emph{not} end an epoch, the partition only has to catch \emph{discontinuities}, which are far rarer and far more visible than the continuous co-evolution the objection has in mind. Second, and more usefully, the consequences of drawing the boundary wrongly are asymmetric and predictable, so the error can be reasoned about rather than merely regretted:

\begin{itemize}[leftmargin=1.4em, itemsep=2pt, topsep=3pt]
\item \emph{Missing a reframe} (coding an epoch-ending move as a within-frame move) merges two epochs. The utility changes mid-sequence, so the retained-best series can \emph{fall} at the join; the drift statistic is biased \emph{downward} and generativity is \emph{under}-estimated. This error therefore works \emph{against} the model's predictions --- it cannot manufacture support for them.
\item \emph{Inventing a reframe} (splitting one epoch in two) discards the information carried across the boundary and shortens the series. It biases estimated generativity \emph{upward} within each fragment and inflates the apparent frequency of ``reframing after a plateau''. This error \emph{can} manufacture support, which is exactly why the coding markers set out in Section~\ref{sec:protocol} are defined observationally and coded before any drift analysis, and why the pre-registration (Section~\ref{sec:prereg}) fixes the segmentation before the outcome variables are computed.
\end{itemize}

\noindent The asymmetry is the point: the error that would flatter the model is the one the protocol is designed to prevent, and the error the protocol cannot fully prevent is the one that penalizes the model. A sensitivity analysis over plausible alternative segmentations is a pre-registered secondary analysis for the same reason.

\section{Internal Consistency and Parameter Sensitivity (In Silico)}
\label{sec:simulation}

We now report a Monte-Carlo study of the model (scripts \texttt{simulation.py}, \texttt{analysis.py}, \texttt{rates.py} and \texttt{reframing.py}, requiring only \texttt{numpy} and \texttt{matplotlib}, with fixed seeds). Its epistemic status needs stating once, plainly: because the simulation implements the model's own assumptions, agreement with the theorem checks \emph{internal consistency and code correctness}, not empirical adequacy. A model of DT is validated against design data, not against itself; Section~\ref{sec:roadmap} sets out how. What the simulation \emph{can} do is (i) verify the analytics are correctly implemented, (ii) characterize the search landscape independently of the process, (iii) measure what happens as each assumption is relaxed, (iv) compare policies the theory leaves open, (v) check that the qualitative findings survive across a \emph{family} of independently generated landscapes rather than the one tuned instance (Section~\ref{sec:landscapes}), and (vi) benchmark the loop against standard search and decision baselines under a matched test budget (Section~\ref{sec:benchmarks}) --- the questions with scientific content. Every ``result'' below is a statement about the model.

\subsection{Experimental Setup}
\label{sec:setup}

A prototype is a binary vector $p\in\{0,1\}^{|F|}$ over $|F|=14$ features, so the prototype space $2^F$ has $16{,}384$ elements --- small enough to enumerate exhaustively, which lets us characterize the landscape exactly rather than assert its shape. The three criteria of Eq.~\eqref{eq:utility} are instantiated as bounded functions of the feature set: desirability rewards covering latent user needs with diminishing returns; viability rewards business value while penalizing scope creep; and feasibility decreases with build cost and with pairwise integration costs (weights $(w_d,w_v,w_f)=(0.45,0.30,0.25)$).

\paragraph{Landscape ruggedness.}
``Convexity'' is undefined on the discrete cube, so instead of calling the landscape non-convex we measure ruggedness by full enumeration. Under the $1$-flip neighbourhood it has \textbf{$4$ local optima} and a fitness--distance correlation to the global optimum of \textbf{$\mathrm{FDC}=-0.22$} (a smooth landscape would have one optimum and $\mathrm{FDC}\approx-1$). This mild ruggedness --- several basins, weak global gradient --- is what makes convergence to a \emph{local} optimum the interesting regime. Note that the multimodal \emph{value histogram} of Figure~\ref{fig:utility_surface}(b) is not itself evidence of a multimodal landscape; the optima count and FDC are the proper measures. These numbers are specific to this instance: stronger interactions or non-additive utilities create more basins, larger $|F|$ raises the count, and a continuous feature space would be characterized instead by a surrogate's length-scale. We use one instance deliberately, to characterize it \emph{exactly}, and claim only the qualitative point the theorem needs --- that a rugged-enough landscape makes local optimality the operative regime; Section~\ref{sec:landscapes} then checks that the qualitative findings replicate across an independently drawn family of $30$ landscapes spanning $1$--$6$ local optima, so nothing below hangs on this single instance.

\paragraph{Process.}
Each iteration applies the operators of Section~\ref{sec:model} in the \emph{bounded-exploration} regime of Assumption~\ref{as:bounded}: $\Op{Ideate}$ mutates a top prototype within Hamming radius $r_{\max}=2$ (no global restarts in the baseline); $\Op{Select}$ chooses the most promising untested candidate; and $\Op{Test}$ returns $y=U(p)+b+\varepsilon$ with $\varepsilon\sim\mathcal{N}(0,\sigma^2)$ and bias $b$ ($b=0$ in the baseline), updating a conjugate Gaussian belief. A retentive team re-tests its leading prototypes, giving the posterior consistency required in Assumption~\ref{as:unbiased}. We average over $250$ runs of $80$ iterations (sensitivity sweeps use $120$ runs of $70$ iterations) and obtain $U^*$ by exhaustive enumeration.





\subsection{Algorithmic and Code Realization}
\label{sec:code}


Algorithm~\ref{alg:ipt} states the iterative core exactly as implemented; it is a literal transcription of the operator semantics of Section~\ref{sec:model}, with the divergent step expanding the candidate set, the convergent step selecting a prototype, and the test step performing a Bayesian update. One implementation detail needs saying, because it is where code and model could silently diverge: an \emph{untested} candidate has no posterior, so it cannot be ranked by posterior mean as $\Op{Select}$ requires. The code therefore orders untested candidates by a cheap proxy $\widehat U$ purely to decide \emph{which to prototype next}, and only \emph{tested} prototypes enter the posterior-mean maximization that defines $X_k$. The proxy affects exploration order, not the quantity whose convergence we study.

\begin{algorithm}
\caption{Ideate--Prototype--Test loop (one run)}
\label{alg:ipt}
\begin{lstlisting}[style=pseudo]
Input: utility landscape $U$, horizon $T$, ideas-per-round $m$, retest count $r$
Output: trajectory $(X_0,\dots,X_{T})$ of max expected utility
initialize beliefs $\gets \varnothing$;  candidates $\gets$ $m$ random prototypes
for $k = 1$ to $T$ do
   // -- Ideate: divergent, BOUNDED to Hamming radius $r_{\max}$ (Asm. 2) --
   base $\gets$ a top prototype by posterior mean
   for $j = 1$ to $m$ do  candidates $\gets$ candidates $\cup\;\{\mathrm{mutate}_{\le r_{\max}}(\text{base})\}$
   // -- Prototype: convergent operator $\Op{Select}$ --
   $p^\star \gets \arg\max_{p\,\in\,\text{candidates}}\ \widehat{U}(p)$   // quick proxy orders *untested* ideas
   move $p^\star$ to beliefs;  initialize neutral prior $\mathcal{N}(0.5,\sigma_0^2)$
   // -- Test: unbiased evidence + Bayesian update --
   beliefs[$p^\star$].update$\big(U(p^\star)+\varepsilon\big)$,  $\varepsilon\sim\mathcal{N}(0,\sigma^2)$
   foreach $p$ in top-$r$ prototypes by posterior mean do
        beliefs[$p$].update$\big(U(p)+\varepsilon\big)$        // re-test the leaders
   // -- Exploratory retest (Asm. 3-4): infinite retesting of every retained $p$ --
   $p_{\mathrm{rr}} \gets$ next materialized prototype in round-robin order
   beliefs[$p_{\mathrm{rr}}$].update$\big(U(p_{\mathrm{rr}})+\varepsilon\big)$
   $X_k \gets \max_{p\,\in\,\text{beliefs}}\ \mathbb{E}[U(p)\mid\mathcal{F}_k]$   // Eq. (*\eqref{eq:Xk}*)
return $(X_0,\dots,X_T)$
\end{lstlisting}
\end{algorithm}

The three components of the algorithm map directly onto the model. (i)~The utility of Eq.~\eqref{eq:utility} combines three bounded criteria $U=w_dU_D+w_vU_V+w_fU_F$; the pairwise integration cost in the feasibility term $U_F$ is what makes the surface rugged (multiple local optima under the mutation neighbourhood; Figure~\ref{fig:utility_surface}). (ii)~Each test is a conditionally unbiased measurement whose conjugate normal--normal update makes the posterior mean a martingale (Assumption~\ref{as:unbiased}), the property exploited in the proof of Theorem~\ref{thm:converge} and verified in Figure~\ref{fig:bayesian_update}. (iii)~The convergence claim is checked by running the loop over an ensemble and confirming that the mean one-step increment of the retained best $X_k$ is non-negative at every iteration (Figure~\ref{fig:submartingale}b). Two implementation choices make the code match the idealized process of Theorem~\ref{thm:converge} rather than merely resemble it (Remark~\ref{rem:idealvsalgo}): the \texttt{UtilityBelief} updater projects its posterior mean onto $[0,1]$ (Remark~\ref{rem:bounded}), and each round performs one additional \emph{round-robin} retest so that every materialized prototype is revisited infinitely often, discharging the infinite-retest clause of Assumption~\ref{as:unbiased}. Both have empirically negligible effect on the trajectories. The full commented implementation of all three --- \texttt{utility()}, the \texttt{UtilityBelief} conjugate updater and the ensemble driver --- is released with the paper (\texttt{simulation.py}).

\subsection{Internal-Consistency Checks (In Silico)}
\label{sec:checks}

The simulation reproduces, \emph{in silico}, each behaviour the theorem predicts (Table~\ref{tab:sim}).

\begin{table}[h!]
\centering
\caption{Theoretical predictions and their \emph{in-silico} checks (250-run ensemble). These verify that the model behaves as proved; they are not empirical evidence about design teams (cf.\ Section~\ref{sec:roadmap}).}
\label{tab:sim}
\begin{adjustbox}{max width=\textwidth}
\begin{tabular}{@{}p{0.30\textwidth}p{0.40\textwidth}p{0.27\textwidth}@{}}
\toprule
\textbf{Theoretical claim} & \textbf{Numerical test} & \textbf{Result} \\
\midrule
Submartingale property (Thm.~\ref{thm:converge}): $\E[\Delta X_k\mid\filt_k]\ge 0$ &
Mean one-step increment of $X_k$ over the ensemble &
$\ge 0$ at $100\%$ of steps (Fig.~\ref{fig:submartingale}b) \\
Boundedness \& a.s.\ convergence &
Ensemble mean of $X_k$ vs.\ estimated $U^*$ &
plateaus at $0.737 < U^*=0.763$ \\
Convergence to a \emph{local} optimum &
Fraction of runs with limit $X_\infty < U^*$ &
$96.8\%$ (Fig.~\ref{fig:submartingale}c) \\
Bayesian update / martingale beliefs (Asm.~\ref{as:unbiased}) &
Posterior mean vs.\ true utility after testing &
$0.589$ vs.\ $u^*=0.579$ (Fig.~\ref{fig:bayesian_update}) \\
Divergence raises, convergence lowers entropy &
$H(v_k)$ across pure operator phases &
$1.0 \to 3.78 \to 2.59$ bits (Fig.~\ref{fig:entropy_increase}) \\
$\epsilon$-stopping rule (Def.~\ref{def:convergence}) &
First crossing of marginal gain below $\epsilon=0.01$ &
$k^*=4$ (Fig.~\ref{fig:threshold}) \\
\bottomrule
\end{tabular}
\end{adjustbox}
\end{table}

The headline check --- a check, not a finding, since the submartingale property is built into the assumptions --- is that the code exhibits the predicted drift: the expected one-step increment $\E[\Delta X_k]$ is non-negative at every iteration (Figure~\ref{fig:submartingale}b) while individual trajectories show the flat-then-jump pattern of a discovery process. The ensemble mean of $X_k$ saturates strictly below $U^*$, and $96.8\%$ of runs settle at local optima (Figure~\ref{fig:submartingale}c). We test the drift with a \emph{single, multiplicity-safe} statistic rather than $80$ per-iteration comparisons: the cumulative drift $\Delta=X_T-X_0$ over a run, for which a submartingale requires $\E[\Delta]\ge0$. Across the ensemble $\E[\Delta]=0.065\pm0.007$ ($95\%$ CI), decisively positive; per-iteration increments are non-negative at $100\%$ of iterations but we rely on the joint statistic for inference. The Bayesian, entropy and stopping-rule experiments corroborate the remaining assumptions.

\paragraph{Convergence diagnostics.}
Figure~\ref{fig:diagnostics} reports three diagnostics that go beyond the binary pass/fail of Table~\ref{tab:sim}. Panel~(a) plots the Lyapunov-type potential $\Phi_k = U^*-X_k$ on a logarithmic scale: it decays monotonically in expectation, as required of the non-negative supermartingale of Section~\ref{sec:lyapunov}, and then flattens at the residual optimality gap that separates the local from the global optimum. Panel~(b) shows the distribution of the per-run \emph{$\epsilon$-stopping} iteration $k^*$ (median $2$, $90$th percentile $6$) --- a number we are careful not to over-read, for reasons Section~\ref{sec:rates-sim} makes precise. Panel~(c) tracks the breadth of exploration: the number of materialized prototypes $|P_r|$ grows linearly while the state entropy $H(v_k)$ rises and saturates, quantifying the divergence pressure that the convergent operators must counterbalance.

\begin{figure}[h!]
    \centering
    \includegraphics[width=\linewidth]{Fig9.png}
    \caption{Convergence diagnostics over the $250$-run ensemble. \textbf{(a)} The Lyapunov potential $\Phi_k=U^*-X_k$ decays (log scale) toward the residual optimality gap. \textbf{(b)} Distribution of the per-run convergence iteration $k^*$ (mean $2.5$, median $2$). \textbf{(c)} Exploration breadth $|P_r|$ and state entropy $H(v_k)$ versus iteration.}
    \label{fig:diagnostics}
\end{figure}

Ensemble statistics expose the monotone rise of $\E[X_k]$ ($0.666\to0.737$), the decay of the potential $\E[\Phi_k]$ ($0.097\to0.026$) and the shrinking $\mathrm{sd}(X_k)$ ($0.041\to0.011$): the ensemble concentrates as it converges. The distribution of the limit is informative too --- $(P_5,P_{25},P_{50},P_{75},P_{95})(X_\infty)=(0.723,0.729,0.734,0.742,0.757)$, a tight, slightly right-skewed spread in which even the luckiest $5\%$ of runs fall short of $U^*=0.763$, exactly what convergence to \emph{local} optima on a rugged landscape predicts.

\subsection{Rates: Two Notions of ``Converged'', and Why They Differ by an Order of Magnitude}
\label{sec:rates-sim}

The $\epsilon$-stopping iteration $k^*$ has a median of $2$, which invites the reading that the process converges almost immediately --- a reading in obvious tension with the rate discussion of Section~\ref{sec:rates}. The tension is real but it is a conflation of two different quantities, and separating them is worth doing because the same conflation is available to anyone measuring a real design team.

The $\epsilon$-stopping rule declares convergence when the \emph{expected marginal gain} over the next $N$ rounds falls below $\epsilon$ (Definition~\ref{def:convergence}). Improvements in this process are rare and lumpy: most rounds produce nothing, so a short window with a loose tolerance is crossed almost at once. The quantity a designer actually cares about --- the round at which the run reaches the value it will finish with --- is much larger. Measuring both on the same $250$-run ensemble (script \texttt{rates.py}):

\begin{table}[h!]
\centering
\caption{Landscape geometry and realized rates. Local optima, neighbourhood sizes and the longest strictly improving chain $L$ are computed by \emph{exhaustive enumeration} of all $2^{14}=16{,}384$ prototypes; $\hat\delta$ and the timing rows are measured over the $250$-run baseline ensemble.}
\label{tab:rates}
\footnotesize
\begin{tabular}{@{}lccc@{}}
\toprule
\textbf{Neighbourhood radius} & $r_{\max}=1$ & $r_{\max}=2^\dagger$ & $r_{\max}=3$ \\
\midrule
$|\mathcal N|$ (reachable neighbours) & $14$ & $105$ & $469$ \\
local optima (exhaustive) & $4$ & $2$ & $1$ \\
longest improving chain $L$ & $66$ & $532$ & $1{,}647$ \\
mean improving depth & $23.6$ & $255.5$ & $826.7$ \\
\midrule
\multicolumn{4}{@{}l@{}}{\textit{Measured over the baseline ensemble} ($r_{\max}=2$):}\\
measured generativity $\hat\delta$ & \multicolumn{3}{c}{$0.104$}\\
realized improving chain $\E[\Lambda]$ & \multicolumn{3}{c}{$4.8$}\\
rounds to the run's own limit & \multicolumn{3}{c}{mean $46.1$, median $46$}\\
rounds to within $1\%$ of that limit & \multicolumn{3}{c}{mean $31.8$, median $30$}\\
$\epsilon$-stopping iteration $k^*$ ($\epsilon=0.01$, $N=5$) & \multicolumn{3}{c}{mean $2.5$, median $2$}\\
\bottomrule
\end{tabular}
\end{table}

\noindent Three things follow. \emph{(i)} The process is \emph{not} fast: reaching the run's limit takes $\approx46$ rounds, an order of magnitude more than $k^*$ suggests, and even reaching within $1\%$ of it takes $\approx32$. The earlier framing of ``median two iterations'' as fast convergence was an over-reading of the stopping rule and we withdraw it. \emph{(ii)} Proposition~\ref{prop:decomp} is met with equality here, as it should be for a roughly constant $\delta$: $\E[\Lambda]/\hat\delta=4.8/0.104=46.1$ rounds against an observed $46.1$. The value of the identity is not the fit but the attribution --- it splits the observed duration into a landscape factor and a team factor that a real study can estimate separately. \emph{(iii)} The worst-case bound of Proposition~\ref{prop:rate}, $L/\delta\approx5{,}100$ rounds, is valid and vacuous --- the realized chain is two orders of magnitude shorter than the longest chain in the space. \emph{(iv)} The fast \emph{apparent} convergence is not an artifact of the small feature space. $|\mathcal N|=105$ is what the process searches per plateau, and it would be $|\mathcal N|=\binom{|F|}{\le2}$ for any $|F|$; the enumerable $|2^F|=16{,}384$ is what lets us characterize the landscape exactly, not what the team searches. Scaling $|F|$ therefore lengthens the improving chain and multiplies the basins --- degrading the \emph{quality} of the limit --- more than it slows the climb, exactly as Section~\ref{sec:rates} predicts.

\subsection{Does the Reframing Criterion Pay? A Policy Comparison}
\label{sec:reframe-sim}

Proposition~\ref{prop:reframe} says a team should reframe when Eq.~\eqref{eq:reframe} flips. Whether following that rule actually produces better designs than the obvious alternatives is a separate question, and one the paper would be weaker for asserting rather than testing. We test it \emph{in silico} (script \texttt{reframing.py}).

\paragraph{Setup.} A run has a fixed budget of $T=150$ rounds. The world contains $M=6$ frames, each an independent landscape with its own feature semantics and hence its own ceiling. Within a frame the team runs the bounded I--P--T loop of Section~\ref{sec:setup}; a reframe costs $c_{\text{ref}}=10$ rounds of the budget (re-empathizing, rebuilding context) and moves the team to a fresh frame with beliefs reset. The payoff is the best \emph{true} utility attained anywhere, $Y_T$ --- a team may ship the best thing it found. $\mathrm{EI}_k$ is computed from the team's own posteriors; $G_{\text{ref}}$ is the frame-level expected improvement under a hierarchical Normal--Normal prior over frame outcomes, updated from the frames already visited. (The symmetry is worth noting: $G_{\text{ref}}$ is itself an expected improvement, one level up --- which is what makes the meta-bandit reading operational rather than metaphorical.) Five policies are compared over $600$ paired worlds: Eq.~\eqref{eq:reframe} itself; a variant that additionally charges the expected \emph{ramp-up} in a fresh frame (the less-myopic, two-level reading); never reframing (over-exploitation); a fixed schedule every $15$ rounds (over-exploration); and a \emph{count-matched random} control, which reframes at random rounds with the number of reframes matched run-by-run to what Eq.~\eqref{eq:reframe} did, isolating the value of reframing at the right \emph{time} from that of reframing the right \emph{number of times}.

\begin{figure}[h!]
    \centering
    \includegraphics[width=\linewidth]{Fig13.png}
    \caption{Does Eq.~\eqref{eq:reframe} pay? ($600$ paired runs, budget $T=150$, $M=6$ frames; error bars are $95\%$ CIs.) \textbf{(a)} Final attained utility by policy. \textbf{(b)} Mean trajectories under a common budget; the reframing cost is visible as the flat segments. \textbf{(c)} Sensitivity to the cost of a reframe: the fixed schedule wins when reframing is cheap and collapses when it is expensive, whereas the criterion tracks the cost.}
    \label{fig:reframe}
\end{figure}

\paragraph{Results, including a negative one.} Following Eq.~\eqref{eq:reframe} beats never reframing decisively: $\E[Y_T]=0.7189\pm0.0020$ against $0.7108\pm0.0023$, a paired difference of $+0.0081\pm0.0024$ ($t=6.5$; the two policies coincide in the $45\%$ of worlds where the criterion never fires, and Eq.~\eqref{eq:reframe} wins $64\%$ of the remainder). It fires sparingly --- a mean of $0.55$ reframes per run. Terminal over-optimization is thus a real and measurable cost, and the criterion detects it. Against the \emph{scheduled} reframer, however, the criterion is \emph{indistinguishable} at the baseline cost ($+0.0003\pm0.0024$, $t=0.3$), and likewise indistinguishable from the count-matched random control ($+0.0003\pm0.0015$, $t=0.4$). We report this negative result rather than suppress it: at $c_{\text{ref}}=10$, in this landscape, \emph{how many} times a team reframes matters and \emph{when} it does so does not.

That is not the end of the story, and panel~(c) shows why. Sweeping the cost of reframing, the fixed schedule is the \emph{best} policy when reframing is cheap ($0.724$ at $c_{\text{ref}}=4$) and the \emph{worst} of the three when it is expensive ($0.706$ at $c_{\text{ref}}=35$, below even never reframing), while Eq.~\eqref{eq:reframe} varies only between $0.7145$ and $0.7187$ across the whole range. The criterion's value is therefore \emph{adaptivity}, not per-configuration optimality: a team that knows its own reframing cost can hand-tune a schedule to match, and a team that does not --- the usual case --- is better served by a rule that reads the cost off its own situation. The less-myopic variant that also charges the expected ramp-up is essentially tied with the plain rule ($0.7183$ vs.\ $0.7189$), which suggests the one-step approximation is not what is limiting performance here, and correspondingly that solving the meta-bandit exactly is unlikely to be where the remaining gains are.

\paragraph{Caveats.} This is an \emph{in-silico} result about the model, subject to the epistemic status of Section~\ref{sec:simulation}: frames are drawn independently, whereas real reframes are informed by what the previous frame taught; the reframing cost is a fixed number of rounds rather than a distribution; and the ceiling prior is well calibrated by construction, where a real team's prior over ``how much better could another framing be'' is exactly the subjective input Section~\ref{sec:reframe} flags. All three plausibly \emph{favour} the scheduled policies, so we read the comparison as a lower bound on the criterion's advantage rather than an estimate of it.

\subsection{How Wrong Can $G_{\text{ref}}$ Be? A Misspecification Study}
\label{sec:gref-sim}

$G_{\text{ref}}$ is the one input Eq.~\eqref{eq:reframe} cannot derive from inside the current frame, so a criterion that failed under a mildly wrong $G_{\text{ref}}$ would not be worth stating. We measure the failure (script \texttt{reframing.py}).

\paragraph{Design.} Running the loop to convergence in each frame of the pool calibrates the distribution of outcomes a team of this capability actually attains, $\mathcal N(0.708,0.028^2)$ --- the correctly specified prior of Eq.~\eqref{eq:gref-prior}. We then displace its mean by $\pm0.05$ and $\pm0.10$ utility (large errors: the entire spread of attainable outcomes is about $0.1$), shrink or inflate its width, and switch off learning entirely, re-running $300$ worlds per condition against the two naive baselines.

\begin{table}[h!]
\centering
\caption{Sensitivity of the Eq.~\eqref{eq:reframe} policy to misspecification of the $G_{\text{ref}}$ prior ($300$ runs per condition, $T=150$, $c_{\text{ref}}=10$; $\pm$ is the $95\%$ CI). The correctly specified prior is $\mathcal N(0.708,0.028^2)$, calibrated from the frame pool.}
\label{tab:gref}
\footnotesize
\begin{tabular}{@{}lccl@{}}
\toprule
\textbf{Prior on the frame outcome} & $\E[Y_T]$ & \textbf{reframes/run} & \textbf{reading} \\
\midrule
pessimistic $-0.10$        & $0.7108\pm0.0032$ & $0.02$ & collapses onto ``never reframe'' \\
pessimistic $-0.05$        & $0.7129\pm0.0030$ & $0.07$ & most of the advantage lost \\
calibrated$^\dagger$       & $0.7154\pm0.0027$ & $0.28$ & \\
optimistic $+0.05$         & $0.7155\pm0.0031$ & $0.68$ & no loss \\
optimistic $+0.10$         & $0.7163\pm0.0028$ & $1.10$ & \emph{better} than calibrated \\
\midrule
diffuse ($\tau=0.15$)      & $0.7170\pm0.0029$ & $0.52$ & width errors are benign \\
over-confident, $+0.05$ off& $0.7216\pm0.0025$ & $0.99$ & \\
no learning, $+0.05$ off   & $0.7254\pm0.0022$ & $1.28$ & best of all conditions \\
\midrule
\textit{baseline:} never reframe   & $0.7101\pm0.0032$ & $0$ & \\
\textit{baseline:} fixed schedule  & $0.7173\pm0.0020$ & --- & \\
\bottomrule
\end{tabular}
\end{table}

\begin{figure}[h!]
    \centering
    \includegraphics[width=0.92\linewidth]{Fig14.png}
    \caption{Robustness of Eq.~\eqref{eq:reframe} to a misspecified $G_{\text{ref}}$ prior ($300$ runs per condition; bands are the $95\%$ CIs of the two baseline policies). \textbf{(a)} Outcome as the prior mean is displaced. Optimism is harmless; pessimism drives the policy onto the never-reframe baseline. \textbf{(b)} The mechanism --- the prior mean controls how often the criterion fires.}
    \label{fig:gref}
\end{figure}

\paragraph{The error is one-sided.} Misspecification does not degrade the criterion symmetrically. \emph{Optimistic} priors cost nothing: at $+0.05$ the outcome is unchanged ($0.7155$ vs.\ $0.7154$) and at $+0.10$ it is slightly better. \emph{Pessimistic} priors are what hurt: at $-0.05$ most of the advantage over never-reframing is gone, and at $-0.10$ the criterion essentially never fires ($0.02$ reframes per run) and its outcome is statistically indistinguishable from never reframing at all. Errors in the prior's \emph{width} are benign in either direction. The practical rule that falls out is blunt and useful: \emph{when uncertain about what another framing could offer, err high.} A team that systematically under-imagines the alternatives does not merely reframe a little less --- it stops reframing.

\paragraph{Why optimism helps, which is not an accident of this landscape.} The most striking rows are the last two: switching off learning, so that an optimistic prior is never revised downward, produces the \emph{best} outcome of any condition ($0.7254$). That looks perverse until one notices what learning does here. The prior is updated with \emph{attained} outcomes, whose mean is $0.708$, while $Y_k$ --- the incumbent the criterion compares against --- is a \emph{maximum} over the frames visited so far. Comparing an average against a maximum drives $G_{\text{ref}}$ toward zero, and the criterion stops firing. Compounding this, Eq.~\eqref{eq:gref} prices a reframe as one draw from the outcome distribution when it is in fact an option on an entire subsequent epoch of iteration. Both biases push the same way, and optimism is simply a crude correction for them.

Three conclusions follow, and the first two are cautionary about our own criterion. \emph{(i)} The right prior for Eq.~\eqref{eq:gref} is over the \emph{good} outcomes of past projects, not the typical ones: an extreme-value prior, not a mean prior. \emph{(ii)} The myopia of the one-step rule lives on the benefit side, which explains the otherwise puzzling result of Section~\ref{sec:reframe-sim} that charging an additional ramp-up cost did not help --- it corrects the wrong side of the inequality. A proper multi-step treatment of the meta-bandit should therefore raise, not lower, the value of switching, and we expect it to close most of the gap between Eq.~\eqref{eq:reframe} and the fixed schedule. \emph{(iii)} Even so, the criterion is never \emph{worse} than never-reframing in any condition tested, including the badly pessimistic ones --- its floor is the naive policy it is meant to improve on, which is the minimum one should ask of a decision rule whose key input is a guess.

\subsection{Sensitivity and Ablation Analysis}
\label{sec:sweep}

To establish that the convergence behaviour is not an artifact of one parameter choice, we sweep the operator parameters --- test-noise $\sigma$, divergence strength $m$, retest budget, confirmation-bias magnitude $b$, neighbourhood radius $r_{\max}$, and prior variance $\sigma_0^2$ --- each over an independent $120$-run ensemble of $70$ iterations on the same landscape ($U^*=0.763$), and ablate one assumption at a time. Tables~\ref{tab:sens} and~\ref{tab:abl} report every configuration, with $95\%$ confidence intervals on $\E[X_\infty]$; Figure~\ref{fig:sensitivity} shows the main trends. The stopping-rule parameters $\epsilon$ and $N$ do not affect the process, only when it is declared converged.



Non-negative drift and below-$U^*$ convergence are observed in \emph{every} configuration: across all twenty parameter settings, all $120/120$ runs in each ensemble converge below $U^*$ with non-negative measured drift. We are careful about what this can show: a simulation cannot establish \emph{almost-sure} convergence, which is an asymptotic property; it can only report that no run violated the predicted behaviour over the observed horizon. The trends are intuitive: noisier tests ($\sigma\uparrow$) and larger confirmation bias ($b\uparrow$) erode the attainable utility monotonically; a wider mutation radius $r_{\max}$, a larger retest budget, and a more diffuse prior ($\sigma_0^2\uparrow$) modestly improve it. The bias sweep answers the robustness question raised about Assumption~\ref{as:unbiased}: even substantial systematic bias does not break convergence, but it lowers the quality of the converged solution in proportion to $b$, exactly as Proposition~\ref{prop:bias} predicts. The ablations isolate which assumption carries the weight: removing belief refinement (``no retest'') costs the most ($0.736\to0.719$ with nearly double the variance), confirmation bias at $b=0.1$ costs $0.736\to0.728$, and unbounded global restarts --- a within-frame move violating Assumption~\ref{as:bounded}, not a reframe --- barely change the finite-horizon outcome ($0.738$), because at practical horizons a random restart almost never lands in the global basin. This resolves the ``restarts reach everything'' versus ``convergence is local'' tension: they live at different horizons.

\begin{figure}[h!]
    \centering
    \includegraphics[width=\linewidth]{Fig10.png}
    \caption{Sensitivity and ablation results (error bars / $\pm$ are $95\%$ CIs). \textbf{(a)} Final utility degrades gracefully as test noise $\sigma$ increases. \textbf{(b)} Systematic confirmation bias $b$ lowers final utility monotonically (Proposition~\ref{prop:bias}). \textbf{(c)} Ablating assumptions: removing belief refinement (``no retest'') causes the largest drop; unbounded ``global restart'' barely changes the finite-horizon outcome, consistent with the theory.}
    \label{fig:sensitivity}
\end{figure}



\begin{table}[h!]
\centering
\caption{Parameter-sensitivity sweep ($120$ runs each, $U^*=0.763$; $\pm$ is the $95\%$ CI on $\E[X_\infty]$). $^\dagger$ marks the baseline value.}
\label{tab:sens}
\footnotesize
\begin{tabular}{@{}lccccc@{}}
\toprule
\textbf{Configuration} & $\E[X_\infty]$ & gap to $U^*$ & \% local & median $k^*$ & drift $\ge0$ \\
\midrule
$\sigma$=0.05 & 0.7429\,$\pm$\,0.0021 & 0.0205 & 88\% & 3 & 100\% \\
$\sigma$=0.10$^\dagger$ & 0.7364\,$\pm$\,0.0019 & 0.0270 & 96\% & 2 & 100\% \\
$\sigma$=0.20 & 0.7313\,$\pm$\,0.0020 & 0.0320 & 98\% & 1 & 100\% \\
$\sigma$=0.30 & 0.7297\,$\pm$\,0.0022 & 0.0337 & 98\% & 1 & 100\% \\
\midrule
ideas=2 & 0.7334\,$\pm$\,0.0020 & 0.0300 & 96\% & 2 & 100\% \\
ideas=5$^\dagger$ & 0.7364\,$\pm$\,0.0019 & 0.0270 & 96\% & 2 & 100\% \\
ideas=8 & 0.7377\,$\pm$\,0.0020 & 0.0256 & 95\% & 2 & 100\% \\
\midrule
retest=1 & 0.7356\,$\pm$\,0.0020 & 0.0278 & 94\% & 2 & 100\% \\
retest=3$^\dagger$ & 0.7364\,$\pm$\,0.0019 & 0.0270 & 96\% & 2 & 100\% \\
retest=6 & 0.7388\,$\pm$\,0.0019 & 0.0246 & 98\% & 2 & 100\% \\
\midrule
$b$=0.00$^\dagger$ & 0.7364\,$\pm$\,0.0019 & 0.0270 & 96\% & 2 & 100\% \\
$b$=0.05 & 0.7341\,$\pm$\,0.0019 & 0.0293 & 97\% & 1 & 100\% \\
$b$=0.10 & 0.7283\,$\pm$\,0.0022 & 0.0350 & 98\% & 1 & 100\% \\
$b$=0.20 & 0.7097\,$\pm$\,0.0039 & 0.0536 & 100\% & 0 & 100\% \\
\midrule
$r_{\max}$=1 & 0.7358\,$\pm$\,0.0019 & 0.0276 & 99\% & 2 & 100\% \\
$r_{\max}$=2$^\dagger$ & 0.7364\,$\pm$\,0.0019 & 0.0270 & 96\% & 2 & 100\% \\
$r_{\max}$=3 & 0.7388\,$\pm$\,0.0019 & 0.0246 & 97\% & 2 & 100\% \\
\midrule
$\sigma_0^2$=0.03 & 0.7366\,$\pm$\,0.0021 & 0.0268 & 98\% & 2 & 100\% \\
$\sigma_0^2$=0.08$^\dagger$ & 0.7364\,$\pm$\,0.0019 & 0.0270 & 96\% & 2 & 100\% \\
$\sigma_0^2$=0.20 & 0.7375\,$\pm$\,0.0020 & 0.0259 & 97\% & 2 & 100\% \\
\bottomrule
\end{tabular}
\end{table}

\begin{table}[h!]
\centering
\caption{Ablation of modelling assumptions ($120$ runs each, $U^*=0.763$; $\pm$ is the $95\%$ CI).}
\label{tab:abl}
\footnotesize
\begin{tabular}{@{}lcccc@{}}
\toprule
\textbf{Configuration} & $\E[X_\infty]$ & gap to $U^*$ & \% local & median $k^*$ \\
\midrule
baseline & 0.7364\,$\pm$\,0.0019 & 0.0270 & 96\% & 2 \\
confirmation bias ($b{=}0.1$) & 0.7283\,$\pm$\,0.0022 & 0.0350 & 98\% & 1 \\
no retest (greedy belief) & 0.7192\,$\pm$\,0.0033 & 0.0442 & 100\% & 1 \\
global restart (within-frame, unbounded $\mathcal{N}$) & 0.7376\,$\pm$\,0.0019 & 0.0258 & 98\% & 2 \\
\bottomrule
\end{tabular}
\end{table}

\subsection{Robustness Across a Family of Landscapes}
\label{sec:landscapes}

The sweeps above vary the operator parameters on \emph{one} landscape, which establishes robustness to parameters but not to the \emph{class} of landscapes --- a fair objection, since the reported effects could be artefacts of the single hand-built utility. We therefore draw a family of $M=30$ independently generated landscapes (script \texttt{landscapes.py}), varying the pairwise-interaction (epistasis) strength, the number of latent needs (modality), and the value spread, and characterize each \emph{exactly} by enumeration as in Section~\ref{sec:setup}. The family spans $1$--$6$ local optima (median $3$) and a fitness--distance correlation $\mathrm{FDC}\in[-0.84,+0.06]$, so it ranges from smooth to markedly rugged. On each landscape we re-run the loop ($K=40$ runs per condition) and ask two questions the single-instance study cannot.

\paragraph{The theorem signatures generalize.} The qualitative behaviour the theorem predicts is present across the family, not just on the tuned instance. Cumulative drift is positive in \emph{every} run on $83\%$ of landscapes, with a mean drift-positive fraction of $99.6\%$; runs settle at a \emph{local} optimum strictly below $U^*$ on $93.1\%$ of runs on average (minimum $72\%$); and the attained fraction of the global optimum averages $97.7\%$ (range $95.1$--$98.8\%$). Crucially, the \emph{residual gap} is explained by ruggedness in the way the model says it should: its correlation with $\mathrm{FDC}$ is $+0.74$ (landscapes with a weaker global gradient, $\mathrm{FDC}$ nearer $0$, leave a larger gap), while its correlation with the raw local-optimum \emph{count} is negligible ($+0.04$) --- consistent with $\mathrm{FDC}$, not the optima count, being the informative ruggedness measure (Section~\ref{sec:setup}).

\paragraph{The headline ablation ordering survives.} The distinctive prediction that \emph{belief refinement} (honest re-testing) matters more than \emph{ideation volume} is not an artefact of the one utility function either. Averaged over the family, removing belief refinement costs $+0.004$ of final utility against $+0.001$ for extra ideation, and belief refinement dominates extra ideation on $73\%$ of the landscapes. The magnitude of the gap shrinks on the easier landscapes (where any policy nearly saturates), but the \emph{ordering} is preserved, which is the claim that matters for the ``testing is the typically binding constraint'' reading (Section~\ref{sec:limitations}). Figure~\ref{fig:landscapes} shows all three findings; a per-landscape random-effect summary is written to \texttt{landscapes\_results.json}.

\begin{figure}[h!]
    \centering
    \includegraphics[width=\linewidth]{Fig15.png}
    \caption{Robustness across a family of $M=30$ independently generated landscapes ($K=40$ runs each). \textbf{(a)} Attained fraction $\E[X_\infty]/U^*$ against ruggedness (local-optimum count): every landscape converges to a local optimum strictly below $U^*$. \textbf{(b)} The residual gap $U^*-\E[X_\infty]$ rises as the landscape becomes more rugged (FDC toward $0$). \textbf{(c)} Per-landscape benefit of belief refinement (re-testing) versus extra ideation; refinement dominates on $73\%$ of landscapes. This is a check across a class of landscapes, still \emph{in silico} and within the additive desirability--viability--feasibility family; it is not empirical validation.}
    \label{fig:landscapes}
\end{figure}

\subsection{Comparison with Baseline Search and Decision Policies}
\label{sec:benchmarks}

A separate question is how the Ideate--Prototype--Test loop compares to standard search and decision baselines --- a comparison the paper would be weaker for omitting, since a convergence guarantee is cheap if the underlying procedure is a strawman. We benchmark it (script \texttt{benchmarks.py}) against four baselines: memoryless \emph{random search}, \emph{greedy hill-climbing}, an \emph{$\varepsilon$-greedy} bandit, and a \emph{UCB / best-arm} policy. All five share the same landscape, the same bounded mutation operator (Assumption~\ref{as:bounded}), the same noisy test oracle $y=U(p)+\varepsilon$, the same projected-Gaussian belief (Remark~\ref{rem:bounded}) and --- the currency of cost in the model --- an \emph{identical budget of $300$ tests} per run; each returns the true utility of the prototype it would ship (its posterior-mean incumbent). We report means with $95\%$ CIs, the Monte-Carlo standard error (MCSE), and Cohen's $d$ against the loop (Table~\ref{tab:bench}, Figure~\ref{fig:benchmarks}).

\begin{table}[h!]
\centering
\caption{Matched-budget comparison ($300$ tests per run). Main-landscape shipped utility $\E[Y]$ over $250$ runs with $95\%$ CI, Monte-Carlo standard error, and Cohen's $d$ versus the DT loop; the last column is the mean attained fraction of $U^*$ across a $10$-landscape family ($30$ runs each). A negative $d$ means the baseline underperforms the loop.}
\label{tab:bench}
\footnotesize
\begin{tabular}{@{}lccccc@{}}
\toprule
\textbf{Policy} & $\E[Y]$ (main) & $95\%$ CI & MCSE & $d$ vs.\ DT & attained \% (family) \\
\midrule
random search            & $0.6944$ & $\pm0.0033$ & $0.0017$ & $-1.38$ & $87.5\%$ \\
greedy hill-climbing     & $0.6779$ & $\pm0.0047$ & $0.0024$ & $-1.60$ & $84.5\%$ \\
$\varepsilon$-greedy     & $0.7177$ & $\pm0.0020$ & $0.0010$ & $-0.40$ & $94.2\%$ \\
UCB / best-arm           & $0.7144$ & $\pm0.0023$ & $0.0012$ & $-0.57$ & $93.6\%$ \\
\textbf{DT loop (this paper)} & $\mathbf{0.7238}$ & $\pm0.0017$ & $0.0009$ & --- & $\mathbf{95.4\%}$ \\
\bottomrule
\end{tabular}
\end{table}

\paragraph{Where the loop stands.} On the main landscape the loop attains the highest shipped utility, $\E[Y]=0.7238\pm0.0017$: ahead of the tuned bandit baselines --- $\varepsilon$-greedy ($0.7177$, $d=0.40$) and UCB ($0.7144$, $d=0.57$) --- and well ahead of random search ($0.6944$, $d=1.38$) and greedy hill-climbing ($0.6779$, $d=1.60$), whose myopia strands it at poor optima. Across the $10$-landscape family the loop attains $95.4\%$ of $U^*$ on average, above $\varepsilon$-greedy ($94.2\%$) and UCB ($93.6\%$), and is the best or tied-best policy on $8$ of $10$ landscapes. The reading is deliberately modest: the loop's \emph{retain-plus-diversify-plus-retest} structure is competitive with, and here slightly ahead of, tuned bandit policies, and clearly better than memoryless or purely greedy search --- so the convergence guarantee is a guarantee about a \emph{reasonable} decision procedure, not a degenerate one. The comparison is \emph{in silico}, matched by test count rather than wall-clock effort, and the baselines are only lightly tuned; the point is not a horse-race trophy but the exclusion of the ``your loop is a strawman'' objection.

\begin{figure}[h!]
    \centering
    \includegraphics[width=\linewidth]{Fig16.png}
    \caption{Matched-budget benchmark ($300$ tests per run; error bars are $95\%$ CIs). \textbf{(a)} Shipped true utility on the main landscape: the DT loop leads the tuned bandit baselines and dominates random and greedy search. \textbf{(b)} Mean attained fraction of $U^*$ across a $10$-landscape family; the loop is best or tied on $8/10$. This is an \emph{in-silico}, matched-test-budget comparison, not empirical evidence about design teams.}
    \label{fig:benchmarks}
\end{figure}

\section{Discussion and Implications}

\subsection{Reconciling Formalism with the `Chaos' of Creativity}

A legitimate worry is that rigour trivializes the creativity at the heart of innovation \citep{mallinson2019, menning2018}. The model avoids that by making no attempt to formalize the ``spark'', formalizing instead the \textbf{decision-making scaffolding within which creativity operates}. Divergent operators such as $\Op{Ideate}$ are deliberately \textbf{black boxes}: we formalize their \emph{effect} (a growing idea set, rising state entropy) and not the mechanism that produces a good idea, which is where intuition, framing and imagination live. Convergent operators are where a Bayesian rule supplies transparent, auditable structure for choosing among what creativity generated --- and the low-coherence ``creative leaps'' that empirically raise ideational productivity \citep{menning2018} are exactly the entropy-increasing moves our divergent operators describe. The apparatus does not generate ideas; it keeps the team honest about which ideas survive contact with evidence.

\subsection{Why, and When, Design Thinking Works}

Theorem~\ref{thm:converge} makes the operative mechanism precise: under unbiased testing the best expected utility is a bounded submartingale, so the loop has a built-in non-negative upward drift --- each honest test can only keep or improve the best-known concept in expectation. That is a precise restatement of the practitioner's intuition that disciplined iteration ``makes things better,'' and it is robust to the messiness of practice because it depends not on having the right idea first but only on testing fairly and retaining the best. The same formalism delineates the limits: on a rugged, only locally observable landscape the guarantee is \emph{local}, never global. In design-theoretic terms the model \emph{implies} that DT delivers reliable \emph{satisficing} rather than optimization, that \emph{testing quality} rather than any idea's brilliance is the \emph{typically} binding constraint (in the model, and in simulation), and that escaping a poor local optimum requires a reframe rather than further within-frame iteration. The theory's distinctive content is less the bare drift result than the \emph{decomposition} of generativity into generation, materialization and selection (Eq.~\eqref{eq:delta-factor}), which turns ``why iteration helps'' into three separately measurable factors rather than a single slogan.

\subsection{Pathways to Computational Implementation}
\label{sec:tools}

The model is a blueprint for design-support tools \citep{uebernickel2020, wuttke2024}, and this reading also resolves the cognitive-plausibility objection: designers need not compute posteriors if a tool does the bookkeeping while they retain judgement. Existing tools support \emph{individual activities} --- whiteboards capture ideas, design-space-exploration systems \citep{woodbury2006} enumerate alternatives --- but none maintains a belief over concept quality across the whole loop or reasons about when the process has converged or should reframe. What the model supplies is that \emph{process-level} logic: a belief state $X_k$, a stopping criterion (Definition~\ref{def:convergence}), a reframing criterion (Section~\ref{sec:reframe}), and an assumption-monitor for biased or impoverished testing. Three design commitments follow, each doubling as a testable claim. \emph{(i) It augments, it does not decide:} the tool surfaces $\E[U(p)\mid\filt_k]$, the current $\mathrm{EI}_k$ and the reframing criterion, and flags biased testing (incumbent-only re-tests, Section~\ref{sec:bias}), but never overrides the team --- it says where the evidence points and how strongly, leaving the leap to them. \emph{(ii) It represents disagreement rather than averaging it away:} carrying several utilities $\{U^{(s)}\}$, it displays each candidate's vector of expected utilities and the resulting Pareto front rather than a scalarized ranking --- which Proposition~\ref{prop:pareto} shows is not a compromise but the right object, since it is the front, not a point, that provably converges. It should also flag a \emph{sparse} front, the failure mode Proposition~\ref{prop:pareto}(b) predicts. \emph{(iii) It is honest about its own assumptions:} the guarantee is conditional on unbiased testing and a fixed frame, so the tool monitors those conditions and warns when they fail. Generative AI could additionally power $\Op{Ideate}$ from a formal Point of View \citep{brenner2016, enkmann2025}. Real barriers remain --- eliciting coarse utilities from busy stakeholders, integrating heterogeneous evidence, organizational resistance to making testing quality visible --- and are empirical questions the framework sharpens rather than settles.

\subsection{Scope, Limitations, and the Theory--Practice Gap}
\label{sec:limitations}

Section~\ref{sec:scope-normative} stated the model's normative status; this section states what it leaves out. Table~\ref{tab:scope} summarizes what is included, what is omitted, and how each omission bears on the conclusions; two entries deserve expansion.

\begin{table}[h!]
\centering
\caption{What the model includes, what it omits, and the effect of each omission on the conclusions.}
\label{tab:scope}
\footnotesize
\setlength{\tabcolsep}{5pt}
\begin{tabularx}{\textwidth}{@{}>{\raggedright\arraybackslash}p{0.20\textwidth} >{\raggedright\arraybackslash}X >{\raggedright\arraybackslash}X@{}}
\toprule
\textbf{Modelling choice} & \textbf{Included / Omitted} & \textbf{Effect on conclusions} \\
\midrule
Knowledge representation & Sets of insights/needs/ideas/prototypes; \emph{omits} tacit, material, aesthetic, ethical dimensions & Results are conditional on a quantified utility; the core is not all of designing \\
Utility & Bounded scalar $U$; \emph{omits} non-compensatory, fuzzy, stakeholder-specific forms & Theorem needs only boundedness, so it survives these; landscape shape and optima change \\
Utility correctness & Assumes the team optimizes the \emph{right} $U$ & If $U$ is mis-specified, convergence is to a local optimum of the wrong objective \\
Creativity ($\delta$) & Assumes improvements are generated with prob.\ $\ge\delta$; \emph{omits} the mechanism & Guarantee is conditional on generativity; $\delta$ is team/practice-specific, unmeasured here \\
Exploration & Bounded (Hamming ball); \emph{omits} unbounded creative leaps within an epoch & Makes the limit \emph{local}; leaps are modelled as reframes (new epoch) \\
Testing & Approximately unbiased, consistent; \emph{omits} unmodelled bias types & $B=0$ gives a submartingale; bounded bias degrades gracefully (Prop.~\ref{prop:bias}) \\
Framing & Fixed within an epoch; \emph{omits} continuous co-evolution & Convergence is within-epoch; reframing resets $(F,U)$ \\
Agency & Single rational decision-maker; \emph{omits} multi-stakeholder politics & Convergence holds per stakeholder; consensus is a separate social act \\
Cognition & Normative/tool logic; \emph{omits} bounded-rationality heuristics & Qualitative conclusions expected to survive; constants would change \\
\bottomrule
\end{tabularx}
\end{table}

\paragraph{What the representation leaves out.} The state $(I,N,\Pi,S_i,P_r)$ is deliberately coarse: it omits tacit knowledge, embodied and material constraints, aesthetic judgement, stakeholder power, and the ethical dimensions of design --- much of designing. The claim is only that, \emph{conditional on} a quantified utility and evidence, the convergent decision core behaves as analysed; extending the state via FBS \citep{gero2004} or C--K knowledge spaces \citep{hatchuel2009} is the natural next step. What the formalism buys beyond ``good practices help'' is three things a verbal account does not: \emph{which} practices matter and by how much (belief refinement is worth about four times as much as extra ideation, Tables~\ref{tab:sens}--\ref{tab:abl}); the \emph{shape} of failure (graceful, Proposition~\ref{prop:bias}, not catastrophic); and a computable \emph{trigger} for reframing, whose sensitivity to its own key input we measure rather than assume (Section~\ref{sec:gref-sim}).

\paragraph{Utility: mis-specification and non-compensatory forms.} Treating $U_D,U_V,U_F$ as precisely measurable and additively combined is a strong simplification, but the machinery does \emph{not} require additivity --- only that $U$ be bounded and time-invariant within the epoch. A thresholded or lexicographic utility such as $U(p)=\mathbf 1[U_F(p)\ge\theta](w_dU_D+w_vU_V)$ is still bounded, so Theorem~\ref{thm:converge} applies verbatim; only the landscape's \emph{shape} changes, non-compensatory feasibility cliffs making it more rugged and so raising the value of reframing. Convergence is likewise robust to coarse or merely \emph{ordinally correct} utilities, since the limiting incumbent depends on $U$ only through the ranking it induces near the top --- which is why ``directional'' use is the regime in which the machinery is most defensible. But a $U$ mis-specified so as to reorder near-best options still converges --- to a local optimum of the \emph{wrong} objective. Disciplined iteration will efficiently perfect the wrong thing, relocating the burden onto framing and elicitation: a formal restatement of ``you can build the thing right and still build the wrong thing.'' Stakeholder disagreement is treated separately in Section~\ref{sec:pareto}; accordingly Definition~\ref{def:convergence}(b) is a normative benchmark a support tool approximates, not a calculator a team runs unaided.

\paragraph{Cognition, sociality, and scaling.} Designers do not compute posteriors in real time: the Bayesian machinery is normative, read as a benchmark or a tool's logic, never as a description of unaided cognition, and reconciling it with bounded rationality \citep{simon1956} would leave the qualitative conclusions intact while changing constants. DT is collaborative and political, whereas we abstract the team to one decision-maker; Section~\ref{sec:pareto} handles conflicting utilities, but consensus costs and the conversational dynamics \citet{menning2018} documents remain outside the model, though the operators are agnostic to whether one agent or many enact them. On \emph{scale}, the convergence theorem is dimension-free while the rate is not, as Section~\ref{sec:rates} sets out; global optimization over the full design graph is expected to be NP-hard \citep{wikipediaNDb}, which is exactly why DT --- like Bayesian optimization --- targets good local solutions. Finally, the formal skeleton invites rigorous comparison with TRIZ/SIT \citep{altshuller1999} and the theories of Section~\ref{sec:related}.

\section{From Verification to Validation}
\label{sec:validation}\label{sec:roadmap}

Because the simulation checks internal consistency rather than empirical adequacy, the model earns its keep only if it makes predictions that could be \emph{wrong} when confronted with design data. It does, and they are distinct from what a purely descriptive account would say:
\begin{enumerate}
    \item \textbf{Drift.} Within a fixed framing, the tested quality of a team's leading concept should be non-decreasing on average over honest test--iterate cycles; runs that violate this should coincide with retention failures or biased testing, not with ``bad luck.''
    \item \textbf{Testing is the binding constraint.} Across projects, variation in \emph{testing quality} (frequency, unbiasedness) should predict final solution quality better than variation in idea novelty or quantity --- a prediction the ablation makes quantitative (Table~\ref{tab:abl}).
    \item \textbf{Reframing, not iteration, escapes plateaus.} Once marginal test gains fall below threshold, further iteration within the frame should not help; measurable improvement should follow reframes (backtracking edges), not additional cycles.
    \item \textbf{Bias degrades gracefully, not catastrophically.} Introduced confirmation bias should lower outcome quality by a \emph{bounded}, $O(B)$ amount (Proposition~\ref{prop:bias}) --- monotone, and in simulation roughly linear in its magnitude --- not cause divergence.
\end{enumerate}

\noindent Turning these predictions into a test requires three things, and the rest of this section supplies them rather than promising them: a measurement pipeline that turns design protocols into the quantities the predictions are about (Section~\ref{sec:protocol}), a pre-registration of the study that would run it (Section~\ref{sec:prereg}), and a feasibility check showing the pipeline runs end-to-end on real design narratives --- together with an honest measurement of how fragile its conclusions are (Section~\ref{sec:cases}).

\subsection{The Measurement Protocol}
\label{sec:protocol}

The pipeline runs against existing resources --- published design-protocol studies such as the DTRS corpora, the design-conversation data of \citet{menning2018}, or controlled team experiments --- and has five steps.
\begin{enumerate}
\item \textbf{Segmentation.} Segment each transcript into design moves and label each as Ideate, Prototype, or Test using a coding scheme grounded in existing protocol-analysis practice: an \emph{Ideate} move proposes a new concept or variant; a \emph{Prototype} move externalizes a concept into a testable artefact; a \emph{Test} move gathers evaluative evidence (user reaction, expert critique, feasibility check) about a specific concept. Backtracking (reframing) moves --- a return to problem definition --- are coded separately.
\item \textbf{Test-outcome coding.} For each Test move, code (a) the concept it evaluates, (b) the evidence valence on an ordinal scale (e.g.\ a $1$--$5$ desirability/feasibility rating, or a signed critique count), and (c) whether the test could induce bias (incumbent-only re-test, leading questions, non-representative users --- proxies for $b_k\neq0$).
\item \textbf{Reliability.} Two independent coders code $\ge20\%$ of segments; report Krippendorff's $\alpha$ (target $\alpha\ge0.70$ for phase labels and $\ge0.67$ for ordinal outcomes) and adjudicate disagreements before the primary coder completes the corpus. Coding manuals and a reliability sample are released with the data.
\item \textbf{Estimation.} Map each ordinal outcome to a bounded observation $y\in[0,1]$ with an explicit, pre-registered function. For a $1$--$5$ Likert rating $s$ the default is the linear map $y=(s-1)/4$ (used in the worked example below); for signed critique counts, a squashing map $y=\sigma(\alpha\,c)$ with $\alpha$ fixed in advance; for binary adopt/reject, $y\in\{0,1\}$. The mapping is a modelling choice and should be reported and, ideally, its influence checked by re-running with an alternative (e.g.\ a mildly non-linear Likert map) to confirm conclusions are not artefacts of it. Given $y$, update a conjugate ($\mathrm{Beta}$ for $[0,1]$ outcomes, truncated-Gaussian otherwise) posterior for each concept's quality, exactly as in the model; from these compute the retained-best estimate $Y_k$ and the belief $X_k$ per concept and per team. \emph{Stakeholder disagreement} is handled as in Section~\ref{sec:limitations}: code each test with the stakeholder role it represents, fit a per-role posterior, and analyse both the per-role trajectories and the Pareto set rather than forcing a single scalar quality.
\item \textbf{Tests of the predictions.} (1) \emph{Drift:} test $\E[X_{k+1}-X_k]\ge0$ within framing epochs via the cumulative-drift statistic of Section~\ref{sec:checks} (a single test per epoch, avoiding multiplicity). (2) \emph{Testing as binding constraint:} regress final concept quality on coded testing intensity/unbiasedness vs.\ ideation volume/novelty. (3) \emph{Reframing:} test whether post-plateau improvement is larger after coded reframes than after further within-frame cycles. (4) \emph{Graceful degradation:} compare drift/outcome across teams stratified by coded bias level (Proposition~\ref{prop:bias}).
\end{enumerate}

\subsubsection*{Estimating the generativity rate}

\paragraph{Three estimators, written out.} Let $k$ index I--P--T rounds within one coded epoch, $Y_k$ the retained-best coded quality, and $\mathcal{O}_k=\mathbf 1\{Y_k<Y_{\text{end}}\}$ the indicator that an improvement was still available at round $k$ --- identified \emph{post hoc} from the epoch's own final value, which is why $\delta$ is estimated retrospectively and cannot be computed online.

\emph{(E1) Coarse, from published corpora.} With only a two-event transition table, $\widehat{\pi^{\text{gen}}}=n(\text{evaluate}\to\text{generate})/n(\text{evaluate}\to\cdot)$, which in \citet{sung2019} is $39/52\approx0.75$. This bounds $\pi^{\text{gen}}$ \emph{from above} --- it counts any return to generation, not the generation of an \emph{improving} candidate --- and says nothing about the other two factors, so it bounds $\delta$ from above too. It requires no new data collection.

\emph{(E2) Direct, from coded sessions --- a retrospective estimator of the realized improvement-opportunity rate.} From the protocol of Section~\ref{sec:protocol},
\begin{equation}
\label{eq:delta-hat}
\hat\delta \;=\; \frac{\sum_k \mathbf 1\{Y_{k+1}>Y_k\}\,\mathcal{O}_k}{\sum_k \mathcal{O}_k},
\qquad
\hat\delta_k \;=\; \text{the same ratio computed in a moving window},
\end{equation}
the window form being what tests decay and hence summability. This is the estimator we apply to the simulator in Section~\ref{sec:rates-sim}, where it returns $\hat\delta=0.104$. Two biases to declare: truncation at the end of an observed epoch makes $\mathcal{O}_k$ conservative (a team that would have improved later is scored as having had no opportunity), biasing $\hat\delta$ \emph{upward}; and coding error in $Y_k$ biases it \emph{downward} by adding spurious non-improvements. Neither is negligible at realistic corpus sizes, which is why we treat $\hat\delta$ as an order-of-magnitude quantity.

\emph{(E3) Factorized, from instrumented studies.} The three factors of Eq.~\eqref{eq:delta-factor} require three different measurements, and none is recoverable from the others:
\[
\widehat{\pi^{\text{gen}}}=\frac{\#\{\text{proposals rated better than the incumbent by a blind post hoc rater}\}}{\#\{\text{proposals}\}},
\quad
\widehat{\pi^{\mu}}=\frac{\#\{\text{prototypes rated} \ge \text{their idea}\}}{\#\{\text{materializations}\}},
\quad
\widehat{\pi^{\text{sel}}}=\frac{\#\{\text{improvements retained}\}}{\#\{\text{improvements built}\}} .
\]
$\pi^{\text{gen}}$ needs a rater blind to what the team subsequently did (otherwise selection contaminates generation); $\pi^{\mu}$ needs the idea and its realization rated independently, ideally by different raters; $\pi^{\text{sel}}$ needs later-revealed quality, so it is available only retrospectively. The factorization is not identified from outcome data alone --- a low $\hat\delta$ is consistent with any of the three being the culprit --- which is precisely the diagnostic point: distinguishing them requires designing the study to do so, and a team that has not done so does not know which of ideation, prototyping or testing is its binding constraint.

\paragraph{Distinguishing within-epoch exploration from reframing.} The protocol's most delicate coding decision --- and a fair objection is that ``any creative leap could be relabelled a reframe to protect the bounded-exploration assumption'' --- is separating bounded within-epoch exploration from an epoch-ending reframe. We therefore define the boundary by \emph{observable} markers rather than by the assumption we wish to protect. A move is coded as a \emph{reframe} (new epoch) if it exhibits at least one of: (a) an explicit restatement of the problem or Point of View, or a return to Empathize/Define activity; (b) a change in the evaluation criteria or their weights (a different $U$); or (c) a candidate that shares little feature overlap with the current pool \emph{and} is justified by a changed problem understanding rather than by incremental improvement. A move is \emph{within-epoch exploration} if it proposes a variant evaluated against the \emph{same} criteria with substantial feature overlap. Crucially, these markers are coded \emph{before} and independently of any drift analysis, so the exploration/reframe split cannot be tuned to manufacture the predicted result; and the framework makes a falsifiable commitment here --- if teams routinely make large, low-overlap jumps \emph{without} any of markers (a)--(c), the bounded-exploration assumption is simply wrong, and the coding will reveal it. Such judgements are not always clean, and we expect real disagreements. The adjudication rule when markers conflict is to ask whether the \emph{evaluation criteria} changed: a move that leaves $U$ intact is a within-frame move however large it looks, and one that changes $U$ is a reframe however small. Section~\ref{sec:cases} works a real borderline case through this rule. Reporting such adjudications, with the achieved $\alpha$, is part of the protocol.

\paragraph{Power, sampling, and coding caveats.} The relevant unit is the framing epoch, not the team. Treating each epoch's cumulative drift as one observation with the effect size seen \emph{in simulation} ($\E[\Delta]\approx0.06$, s.d.\ $\approx0.05$, Cohen's $d\approx1.2$), a one-sided $t$-test at $\alpha=0.05$ reaches power $0.9$ with $\approx8$ epochs. But that simulation effect is almost certainly an over-estimate: powering for a realistic $d\approx0.3$ needs $\approx100$ epochs, so the honest requirement is several dozen sessions or a multi-study meta-analysis, not the optimistic ``$10$--$15$.'' Moreover, because epochs are \emph{nested within teams} they are not independent: the effective sample size is deflated by the design effect $1+(\bar m-1)\rho$, with $\bar m$ the mean epochs per team and $\rho$ the intraclass correlation (ICC) of the outcome within a team. At, say, $\bar m\approx4$ and a modest $\rho\approx0.1$ the design effect is $\approx1.3$, so $N\approx100$ nominal epochs behave like $\approx75$ independent ones; the ICC must be estimated and reported and the target inflated accordingly. We therefore treat $N=100$ as a floor conditional on low ICC, not a guarantee. Two further cautions: design-process coding has historically achieved lower reliability than our $\alpha$ targets, so teams should pilot the manual and coarsen the scale if it falls short; and the bias markers (incumbent-only re-tests, leading questions, non-representative users) are coarse proxies for $b_k\neq0$ --- they flag when bias is \emph{likely}, not its sign or magnitude. We state these so the study is designed for realistic expectations, not the simulation's best case.

\paragraph{A synthetic dry run.} To show the pipeline is well defined we ran steps~1--5 on a small hand-authored protocol (a coffee-cup session; \texttt{protocol\_demo.py} in the repository). This exercises the machinery, and is not evidence about real design: across the coded protocol the belief $X_k$ has positive cumulative drift ($X_T-X_0=+0.34$) and the retained incumbent never switches to a worse concept, while a single downward blip is a re-test correcting an over-estimate --- the submartingale, not monotone, behaviour of Theorem~\ref{thm:converge}, recovered from coded observations rather than assumed.

\subsection{Pre-Registration of the Validation Study}
\label{sec:prereg}

The model's conditional claims are only worth as much as the study that would test them, and a protocol that is specified only after the data are seen is not a test. We therefore state the study here in pre-registration form. It is a \emph{plan}, not a report: no data have been collected, and nothing else in this paper depends on its outcome. We intend to lodge this plan, unchanged, in a public registry before collection begins, and we state in advance the results that would count against the model.

\subsubsection*{Hypotheses}

Each hypothesis is comparative --- it relates \emph{variation in assumption-satisfaction} to \emph{variation in outcomes} across framing epochs --- because that, and not the behaviour of any single team, is what a normative model implies (Section~\ref{sec:scope-normative}).

\begin{description}[leftmargin=2.6em, style=nextline, itemsep=3pt]
\item[H1 (Drift).] Within a coded framing epoch, the tested quality of the leading concept rises on average. To avoid circularity (Section~\ref{sec:prereg-independence}) we state this on \emph{two} outcomes: \emph{(H1a, primary)} a \emph{blind-coded external quality} rating of the leading concept, coded independently of the Bayesian model, has positive cumulative drift; \emph{(H1b, secondary)} the model-derived incumbent belief satisfies $\E[X_{T}-X_{0}]>0$. \emph{Directional; H1a primary, H1b corroborating.}
\item[H2 (Testing predicts outcomes above ideation).] Across epochs, coded testing quality (frequency $\times$ unbiasedness) predicts final concept quality with a larger standardized coefficient than coded ideation volume or novelty. This is a \emph{predictive association}, not a causal claim: testing quality and general process quality are plausibly confounded (a better-organized team tests better \emph{and} ideates better \emph{and} documents better), so the contrast identifies ``testing is the binding constraint'' only descriptively; a randomized manipulation of the testing regime, not this observational comparison, is what would license the causal reading (Section~\ref{sec:cannot-settle}). \emph{Directional, primary; predictive.}
\item[H3 (Reframing, not iteration, escapes plateaus).] After a coded plateau, improvement in the $10$ moves following a coded reframe exceeds improvement in the $10$ moves following continued within-frame iteration. Reframing is \emph{not randomly assigned} --- teams reframe precisely when the current direction stalls, when new information arrives, or when membership changes --- so this comparison is subject to selection on both arms (which can bias the contrast either way) and identifies a causal effect only under strong ignorability assumptions; we read it as predictive and flag a randomized or quasi-experimental reframing manipulation as the design that would settle it. \emph{Directional, primary; predictive.}
\item[H4 (Graceful degradation).] Across epochs stratified by coded bias level, final quality declines monotonically and continuously in bias, with no discontinuity. (``Approximately linear'' is an \emph{empirical} refinement; the theory guarantees only a bounded, $O(B)$ decline, Proposition~\ref{prop:bias}.) \emph{Directional, secondary.}
\item[H5 (Generativity is non-summable).] Within epochs long enough to estimate it, the windowed $\hat\delta_k$ of Eq.~\eqref{eq:delta-hat} does not decay fast enough for $\sum_k\hat\delta_k$ to converge over the observed horizon. \emph{Exploratory} --- summability is an asymptotic property and a finite record can only fail to contradict it.
\end{description}

\subsubsection*{Why the primary outcome is model-independent}
\label{sec:prereg-independence}
The model quantity $X_k$ is reconstructed from the coded protocol \emph{through the very Bayesian update the theorem is about}, so testing $\E[X_T-X_0]>0$ risks a partial tautology: a coding-and-update pipeline engineered to have submartingale-like behaviour could produce positive drift almost by construction. We therefore make the \emph{primary} outcome an \emph{independent} one --- a blind external quality rating of the leading concept at each epoch boundary, produced by raters who never see the model's belief trajectory --- and demote the model quantity $X_k$ to a secondary, corroborating measure. The claim that then carries confirmatory weight is ``assumption-satisfaction $\rightarrow$ independent outcome,'' not ``model $\rightarrow$ model.'' Agreement between H1a and H1b is itself informative (it argues the coding-plus-update pipeline is not merely reproducing its own inputs), but only H1a is confirmatory.

\subsubsection*{Design, unit of analysis, and sampling}

\textbf{Design.} Observational, retrospective coding of design-session protocols, supplemented where available by prospective recording of student and professional design teams. \textbf{Unit of analysis: the framing epoch, not the team or the session}; a team contributes as many units as it has epochs, and epochs are nested within teams (see the analysis plan below). \textbf{Sampling frame.} Existing protocol corpora (the DTRS datasets; the design-conversation data of \citet{menning2018}) plus prospectively recorded sessions. \textbf{Target $N$.} $100$ framing epochs, from the power analysis of Section~\ref{sec:protocol}: with a realistic $d\approx0.3$ for H1, a one-sided test at $\alpha=0.05$ reaches power $0.9$ at $N\approx100$. We will not report the optimistic $N\approx8$ implied by the simulation's effect size. \textbf{Stopping rule.} Collection stops at $N=100$ epochs or at $24$ months, whichever comes first; no interim analysis of the primary outcomes will be conducted, and no data-dependent extension will be made.

\subsubsection*{Coding, blinding, and the order of operations}

Coding follows Section~\ref{sec:protocol} steps 1--3. Two commitments do the work that prevents the analysis from being tuned:

\begin{enumerate}[leftmargin=1.8em, itemsep=2pt]
\item \textbf{Segmentation precedes outcomes.} Epoch boundaries and move labels are coded, reconciled and frozen \emph{before} any outcome variable is computed. Coders see transcripts, not quality trajectories.
\item \textbf{Reliability gate.} Two independent coders code $\ge20\%$ of segments; we require Krippendorff's $\alpha\ge0.70$ for phase labels and $\ge0.67$ for ordinal outcomes. \emph{If the gate is not met we will coarsen the scale, re-pilot, and report the failure} --- we will not proceed to hypothesis testing on unreliable codes, and we will not silently drop the requirement.
\end{enumerate}

\subsubsection*{Measures}

$X_k$ and $Y_k$ from the conjugate updates of Section~\ref{sec:protocol} step~4, with the ordinal-to-$[0,1]$ map fixed in advance as $y=(s-1)/4$ for $1$--$5$ Likert ratings. Testing quality: tests per move $\times$ ($1-$ proportion of tests carrying a bias marker). Ideation volume: distinct proposals per move. Novelty: mean pairwise feature distance among proposals. $\hat\delta$ from Eq.~\eqref{eq:delta-hat}. All measures are computed by released code from the frozen codes, with no analyst discretion at this stage.

\subsubsection*{Analysis plan}

\textbf{Primary.} H1: the per-epoch cumulative drift of the \emph{independent} outcome (H1a) is tested with a cluster-robust procedure --- a mixed model with random intercepts for team and corpus, cross-checked against a one-sided \emph{cluster (team-level) bootstrap} and a \emph{permutation} test --- rather than a plain one-sample $t$-test, because the drift statistic is heteroskedastic, epoch-length-dependent, scale-bounded and clustered within teams. The model quantity (H1b) is tested the same way as a secondary outcome. A single statistic per epoch avoids the multiplicity of per-iteration tests. H2: mixed-effects regression of final quality on standardized testing quality and ideation volume, with random intercepts for team and corpus; the hypothesis is a contrast between the two coefficients, tested once. H3: paired comparison within epochs that contain both a post-plateau reframe and a post-plateau continuation. \textbf{Multiplicity.} Three primary hypotheses, Holm-corrected at family-wise $\alpha=0.05$. Secondary and exploratory results are reported as such and are not corrected, and will not be described as confirmatory. \textbf{Exclusions, fixed in advance.} Epochs with fewer than $5$ coded test events (too short to estimate drift); sessions where recording failure loses more than $20\%$ of the protocol; corpora that fail the reliability gate. No exclusion may be decided after seeing an epoch's outcome. \textbf{Sensitivity analyses, pre-specified.} (a) Re-run under an alternative segmentation in which every ambiguous boundary is resolved the other way, per Section~\ref{sec:framing}; (b) re-run under a mildly non-linear Likert map; (c) re-run with the coder-perturbation procedure of Section~\ref{sec:cases}, reporting the proportion of perturbed codings that preserve each primary conclusion. Conclusions that do not survive (a)--(c) will be reported as fragile.

\subsubsection*{What would count against the model}

We state these in advance so that they cannot be reinterpreted afterwards.

\begin{itemize}[leftmargin=1.6em, itemsep=2pt]
\item \textbf{H1 fails} (drift not positive) \emph{while} retention, bias and generativity measure as satisfying the assumptions: this contradicts the theorem's conditional and is the cleanest available disconfirmation.
\item \textbf{H2 reverses} --- ideation volume predicts outcomes better than testing quality: this contradicts the model's most distinctive prediction, since the submartingale argument makes testing, not idea supply, the binding constraint.
\item \textbf{H3 fails} --- post-plateau iteration improves outcomes as much as reframing: this would undercut the claim that the guarantee is local and that reframing is the only escape.
\item \textbf{H4 shows a discontinuity} rather than linear decline: this would contradict Proposition~\ref{prop:bias}'s graceful degradation.
\item \textbf{Bounded exploration is violated wholesale} --- teams routinely make large, low-overlap jumps without any of the reframe markers (a)--(c) of Section~\ref{sec:protocol}: this would falsify Assumption~\ref{as:bounded} directly and, with it, the local-optimum characterization.
\end{itemize}

\noindent A null result on any primary hypothesis will be reported as such. We commit to publishing the outcome regardless of direction, and to releasing codes, coding manual, reliability sample and analysis scripts with it.

\subsubsection*{What this study cannot settle}\label{sec:cannot-settle}

For symmetry, the limits. It cannot establish that teams \emph{do} maximize expected utility, which the model does not claim. It cannot test summability, an asymptotic property, on finite records. It cannot separate the three factors of Eq.~\eqref{eq:delta-factor} without the instrumented ratings of estimator (E3), which retrospective corpora do not contain. And observational data cannot license causal claims about testing quality; a randomized manipulation of the testing regime, which we regard as the natural follow-on, would be required for that. The same limitation applies to \emph{bias}: the coded bias markers (incumbent-only re-tests, leading questions, non-representative users) flag when bias is \emph{likely} but do not recover its sign or magnitude ($\text{marker}\neq b_k$), so H4 is a stratified association, and an \emph{instrumental manipulation} of testing bias --- not stratification of observational data --- is what a causal test of Proposition~\ref{prop:bias} would require.

\subsection{A Feasibility Check on Published Cases}
\label{sec:cases}

The last thing the protocol needs is a demonstration that it runs at all, so we applied it to two \emph{published} design-thinking studies as secondary data. \textbf{The purpose of that exercise is procedural, not evidential}, and the reasoning is worth stating plainly rather than hedging. Both cases are reported successes, so the sample cannot exhibit the failure modes (retention loss, biased testing) on which the model's most distinctive predictions bear; a design that cannot produce a disconfirming observation cannot produce a confirming one either. Add $N=2$, a single coder (ourselves), and ordinal outcome codes read off authors' prose, and the appropriate conclusion is that \emph{nothing here is evidence for the model}. What the exercise does establish is that the coding-and-estimation steps run end-to-end on real design narratives and produce the quantities the predictions are about --- which is a precondition for the confirmatory study, not a substitute for it. We also quantify how fragile the resulting numbers are to coder disagreement, and the answer is: quite. The coding sheets and code (\texttt{empirical.py}) are in the repository.

\paragraph{Case A: a longitudinal mHealth project (SAM).} \citet{vanasselt2025} document seven years of co-creative development of the Stress Autism Mate app through Empathize--Define--Ideate--Prototype--Test cycles with autistic panel members. We coded the reported sequence into $14$ test events with bounded outcomes (adopt / confirm / mixed / reject / plateau; step~4 of the protocol) and computed the incumbent belief $X_k$ (Figure~\ref{fig:empirical}a).

We begin by \emph{withdrawing} a claim made in an earlier version of this work. It is tempting to report that ``the retained-best design never regresses across all $14$ events'', and it is true --- but vacuously, because the retained best is computed as a running maximum and is therefore non-decreasing \emph{by construction}. It tests nothing. The non-trivial quantity is the incumbent \emph{belief} $X_k$, which a submartingale permits to fall: across the coded sequence $X_k$ decreases at $2$ of $13$ steps with a cumulative drift of $X_T-X_0=+0.015$ and a mean increment of $+0.0012$. That is the qualitative signature the theory predicts --- upward drift, not monotonicity --- but on $14$ events it is far too small a sample to be more than a demonstration that the quantity is computable.

Two further observations are worth recording as \emph{illustrations of the mechanism}. The one candidate testing scored below the incumbent --- a Fitbit-integration add-on that ``showed no effect \ldots\ discrepancies between different stress indicators could lead to confusion'' --- was correctly not retained, which is the mechanism the submartingale argument rests on. And the two large moves that changed both the feature set and the criteria (the initial shift from a generic ``Self-mate'' concept to stress-focused monitoring, and the 2025 v2.0 overhaul that followed in-depth interviews with \emph{dissatisfied} users reporting ``lack of feature depth'') fall where Definition~\ref{def:reframe} says reframes should fall: after within-frame gains had thinned.

\paragraph{How much would a second coder change this?} The single-coder limitation cannot be repaired retrospectively, but its \emph{consequence} can be bounded, and doing so is more informative than restating the caveat. We simulate a disagreeing coder by shifting each coded outcome one step along the ordinal scale with probability $q=0.30$ and jittering the ordinal-to-$[0,1]$ map, then re-running the pipeline on $20{,}000$ such re-codings. The results are sobering and we report them as such: the rejected candidate still scores below the incumbent in $86.8\%$ of re-codings, but \textbf{the positive cumulative drift of $X_k$ survives in only $61.3\%$} (drift mean $+0.020$, $5$th--$95$th percentile $[-0.056,+0.148]$). In other words, the central quantity of Case~A is \emph{not robust} to plausible coder disagreement. This is precisely why step~3 of the protocol (two coders, reported Krippendorff's $\alpha$) is not a formality, and why we decline to describe this case as supporting the model.

\paragraph{Case B: a design-protocol analysis (Sung \& Kelley).} \citet{sung2019} report a sequential (lag) analysis of think-aloud design protocols, coding cognition into states including \emph{Designing} (generating ideas) and \emph{Predicting} (anticipating a design's consequences --- i.e.\ evaluating). Their published transition statistics let us test whether the generate$\leftrightarrow$evaluate iteration our model assumes is empirically present: the transitions Designing$\to$Predicting ($z=3.66$, $p<0.001$) and Predicting$\to$Designing ($z=5.17$, $p<0.001$) are both significantly \emph{more} frequent than chance (Figure~\ref{fig:empirical}b), i.e.\ a statistically significant, bi-directional Ideate$\leftrightarrow$Test loop --- the empirical counterpart of our iterated cycle. From the same table we can bound a \emph{generativity} rate relevant to Assumption~\ref{as:creative}: after an evaluation (\emph{Predicting}), designers returned to generating (\emph{Designing}) in $39/52\approx0.75$ of cases --- i.e.\ evaluation is followed by renewed idea generation about three-quarters of the time. This is not the model's $\delta$ (the probability an \emph{improving} neighbour is generated \emph{and} selected), but it is an upper-bounding, empirically grounded first estimate of the raw generativity on which $\delta$ depends, and it shows the quantity is measurable from existing protocol data.

\begin{figure}[h!]
    \centering
    \includegraphics[width=\linewidth]{Fig12.png}
    \caption{Feasibility run of the coding pipeline on published cases ($N=2$, single coder, retrospective); this is a feasibility demonstration, not evidence --- see the text and the robustness analysis. \textbf{(a)} Case~A \citep{vanasselt2025}: the incumbent belief $X_k$ across $14$ coded test events, with two down-steps and a cumulative drift of $+0.015$; the running maximum (dashed) is monotone by construction and tests nothing; the one option testing scored below the incumbent (a Fitbit add-on) was rejected; vertical lines mark reframes. \textbf{(b)} Case~B \citep{sung2019}: the Ideate$\leftrightarrow$Test (Designing$\leftrightarrow$Predicting) iteration is statistically significant ($z>1.96$), with a raw generativity rate of $0.75$ bounding $\pi^{\text{gen}}$.}
    \label{fig:empirical}
\end{figure}

\paragraph{What this does and does not show.} It shows that the pipeline runs end-to-end on real design narratives; that its outputs are concrete enough to be checked, argued with, and --- as the robustness analysis demonstrates --- found wanting; and that at least one assumption-relevant quantity ($\pi^{\text{gen}}$, via Case~B) is already recoverable from published data. It does not show that the model is right. Two success cases cannot test the failure predictions, one coder cannot establish that the codes mean anything, and a drift statistic that survives only $61\%$ of plausible re-codings is not a result. The honest summary is that Section~\ref{sec:protocol} has been converted from a promise into a demonstrated procedure with a measured fragility, and that the confirmatory study it describes remains to be done.

\section{Conclusion}

We asked under what sufficient, empirically meaningful conditions the iterative core of Design Thinking provably improves an outcome. The answer is that four structural conditions suffice within a framing epoch --- bounded exploration, retention, non-summable generativity, and approximately unbiased testing --- and that, while we do not establish their necessity, each is \emph{operationally meaningful}: relaxing it degrades the guarantee in a specific, measurable way: bias costs $O(B)$ in the limit (Proposition~\ref{prop:bias}), slow criterion drift and feature accretion cost $O(D)$ (Corollaries~\ref{cor:drift}--\ref{cor:accretion}), and stakeholder disagreement converts the limit from a point to a locally non-dominated set without destroying convergence (Proposition~\ref{prop:pareto}). The theorem itself is a short argument from Doob's theorem; the contribution is the identification of what must be true of a design practice for it to hold.

We have tried throughout to keep three kinds of claim apart. The mathematics \emph{proves} conditional convergence. The simulation \emph{demonstrates} internal consistency and maps the parameter space --- including two findings that cut against the model's own presentation: the process is an order of magnitude slower than its stopping rule suggests (Section~\ref{sec:rates-sim}), and the reframing criterion's advantage over a fixed schedule is a matter of robustness to reframing cost rather than of per-configuration superiority (Section~\ref{sec:reframe-sim}). Whether real teams satisfy the assumptions \emph{remains untested}. Section~\ref{sec:protocol} says how one would find out, Section~\ref{sec:prereg} pre-registers the study and commits in advance to the results that would count against the model, and Section~\ref{sec:cases} shows on published cases both that the pipeline runs and how fragile its conclusions are to plausible coder disagreement. We are explicit that this is a \emph{normative} model throughout (Section~\ref{sec:scope-normative}): it says what a team would have to do, not what teams do, and the evidence it calls for is comparative --- that closer approximation to the four conditions predicts better outcomes. What the formalization offers design research is therefore not a validated theory but a falsifiable vocabulary: a way of saying precisely what ``disciplined iteration works'' means, what would show it false, and where --- at the framing level, which remains genuinely wicked --- it does not reach.

\section*{Data and code availability}
All code required to reproduce the experiments, figures and tables in this paper --- the simulation of the Ideate--Prototype--Test loop (\texttt{simulation.py}), the sensitivity and ablation analysis (\texttt{analysis.py}), the landscape-geometry and rate analysis (\texttt{rates.py}), the reframing-policy comparison (\texttt{reframing.py}), the landscape-family robustness study (\texttt{landscapes.py}), the matched-budget baseline benchmarks (\texttt{benchmarks.py}), the coded-protocol demonstrations (\texttt{protocol\_demo.py}, \texttt{empirical.py}), the figure-generating scripts, and the raw numerical outputs --- is openly available in the repository at \url{https://github.com/pawelkuras/design-thinking-formalization} and will be deposited in a permanent archive (Zenodo) with a citable DOI upon acceptance. The experiments use fixed random seeds and require only \texttt{numpy} and \texttt{matplotlib}. A single driver script (\texttt{run\_all.sh}) re-runs the entire pipeline from scratch in a clean environment and regenerates every figure, table and reported statistic, so the results can be reproduced end-to-end rather than merely inspected.

\section*{CRediT authorship contribution statement}
\textbf{Pawe\l{} Kura\'s:} Conceptualization, Methodology, Formal analysis, Software, Writing -- original draft. \textbf{Adrian Michalski:} Formal analysis, Validation, Writing -- review \& editing. \textbf{Alicja Gerka:} Conceptualization, Investigation, Writing -- review \& editing.

\section*{Declaration of competing interest}
The authors declare that they have no known competing financial interests or personal relationships that could have appeared to influence the work reported in this paper.

\section*{Declaration of generative AI use}
In the interest of full transparency we describe the use of generative AI specifically. A large-language-model assistant was used to (i) draft and revise prose from author-provided outlines, arguments, and bullet points, and to improve clarity and concision; (ii) help write and debug the accompanying Python code that produces the figures, simulation, and analysis; and (iii) assist with reference formatting. It was \emph{not} used to originate the theoretical model, the assumptions, the proofs, the choice of experiments, or the interpretation of results, and it was not used to generate or alter any data. The authors conceived the contribution, verified every mathematical statement and proof line by line, independently checked all numerical results against the released code, confirmed each cited reference against its original source, and take full responsibility for the entire content of the paper. Where AI assistance touched the writing, the authors reviewed and edited all text so that it reflects their own understanding and claims. The accompanying repository preserves the full commit history of the code, so the record of which implementation changes were AI-assisted --- as distinct from the conceptual, mathematical and experimental contributions, which were not --- remains auditable.














\begin{thebibliography}{Wagenmakers et~al.(2018)}

\bibitem[Altshuller(1999)]{altshuller1999} Altshuller, G. (1999). \textit{The Innovation Algorithm: TRIZ, Systematic Innovation and Technical Creativity} (L. Shulyak \& S. Rodman, Trans.). Technical Innovation Center.

\bibitem[Arora(2017)]{wikipediaNDb} Arora, J. S. (2017). \textit{Introduction to Optimum Design} (4th ed.). Academic Press.

\bibitem[Bagaria(2021)]{bagaria2021} Bagaria, J. (2021). Set theory. In E. N. Zalta (Ed.), \textit{The Stanford Encyclopedia of Philosophy} (Winter 2021 ed.). Metaphysics Research Lab, Stanford University. Retrieved from \url{https://plato.stanford.edu/archives/win2021/entries/set-theory/}

\bibitem[Berger(1985)]{theodoridisND} Berger, J. O. (1985). \textit{Statistical Decision Theory and Bayesian Analysis} (2nd ed.). Springer. \doi{10.1007/978-1-4757-4286-2}

\bibitem[Bhamra(2020)]{majka2024} Bhamra, T. S. (2020). The Double Diamond: A universally accepted model of the design process. In C. Cautela, M. Mortati, \& F. Rizzo (Eds.), \textit{Design Management in the Era of Disruption}. Springer. \doi{10.1007/978-981-15-4704-2_4}

\bibitem[Bonchev \& Rouvray(1991)]{smith2024} Bonchev, D., \& Rouvray, D. H. (Eds.). (1991). \textit{Chemical Graph Theory: Introduction and Fundamentals}. Gordon and Breach.

\bibitem[Durrett(2019)]{durrett2019} Durrett, R. (2019). \textit{Probability: Theory and Examples} (5th ed.). Cambridge University Press. \doi{10.1017/9781108591034}

\bibitem[Even-Dar et~al.(2006)]{evendar2006} Even-Dar, E., Mannor, S., \& Mansour, Y. (2006). Action elimination and stopping conditions for the multi-armed bandit and reinforcement learning problems. \textit{Journal of Machine Learning Research, 7}, 1079--1105.

\bibitem[Frazier et~al.(2008)]{frazier2008} Frazier, P. I., Powell, W. B., \& Dayanik, S. (2008). A knowledge-gradient policy for sequential information collection. \textit{SIAM Journal on Control and Optimization, 47}(5), 2410--2439. \doi{10.1137/070693424}

\bibitem[Laumanns et~al.(2002)]{laumanns2002} Laumanns, M., Thiele, L., Deb, K., \& Zitzler, E. (2002). Combining convergence and diversity in evolutionary multiobjective optimization. \textit{Evolutionary Computation, 10}(3), 263--282. \doi{10.1162/106365602760234108}

\bibitem[Brenner \& Uebernickel(2016)]{brenner2016} Brenner, W., \& Uebernickel, F. (Eds.). (2016). \textit{Design Thinking for Innovation: Research and Practice}. Springer. \doi{10.1007/978-3-319-26100-3}

\bibitem[Brenner et~al.(2016)]{abrell2016} Brenner, W., Uebernickel, F., \& Abrell, T. (2016). Design thinking as mindset, process, and toolbox. In W. Brenner \& F. Uebernickel (Eds.), \textit{Design Thinking for Innovation: Research and Practice} (pp. 3--21). Springer. \doi{10.1007/978-3-319-26100-3_1}

\bibitem[Broberg et~al.(2023)]{broberg2023} Broberg, O., Andersen, V., \& Seim, R. (2023). Applying design thinking in occupational health and safety management. In V. Andersen, O. Broberg, A. Karltun, \& N. M{\o}ller (Eds.), \textit{Human-Centered and Error-Resilient Systems Development} (Lecture Notes in Networks and Systems, Vol.\ 785). Springer. \doi{10.1007/978-3-031-41712-8_41}

\bibitem[Brown(2008)]{brown2008} Brown, T. (2008). Design thinking. \textit{Harvard Business Review, 86}(6), 84--92.

\bibitem[Bubeck \& Cesa-Bianchi(2012)]{bubeck2012} Bubeck, S., \& Cesa-Bianchi, N. (2012). Regret analysis of stochastic and nonstochastic multi-armed bandit problems. \textit{Foundations and Trends in Machine Learning, 5}(1), 1--122. \doi{10.1561/2200000024}

\bibitem[Chen(2021)]{chen2021} Chen, C. H. (2021). Case study on the diffusion of digital dental technology innovation: Double Diamond model perspective. \textit{Journal of Healthcare Engineering, 2021}, Article 9924555. \doi{10.1155/2021/9924555}

\bibitem[DeGroot(2004)]{elgamalND} DeGroot, M. H. (2004). \textit{Optimal Statistical Decisions}. Wiley-Interscience. (Original work published 1970)

\bibitem[Design Council(2005)]{designcouncilNDa} Design Council. (2005). \textit{The Double Diamond: A universally accepted depiction of the design process}. Design Council.

\bibitem[Design Council(2007)]{designcouncilNDb} Design Council. (2007). \textit{Eleven Lessons: Managing Design in Eleven Global Brands. A Study of the Design Process}. Design Council.

\bibitem[Diestel(2017)]{mateosND} Diestel, R. (2017). \textit{Graph Theory} (5th ed.). Springer. \doi{10.1007/978-3-662-53622-3}

\bibitem[Dorst \& Cross(2001)]{dorstcross2001} Dorst, K., \& Cross, N. (2001). Creativity in the design process: Co-evolution of problem--solution. \textit{Design Studies, 22}(5), 425--437. \doi{10.1016/S0142-694X(01)00009-6}

\bibitem[Dorst \& Dijkhuis(1995)]{dorst1995} Dorst, K., \& Dijkhuis, J. (1995). Comparing paradigms for describing design activity. \textit{Design Studies, 16}(2), 261--274. \doi{10.1016/0142-694X(94)00012-3}

\bibitem[Drew(2019)]{shabanND} Drew, J. (2019). The Double Diamond: A practical and scalable design process. \textit{Journal of Design, Business \& Society, 5}(2), 143--162. \doi{10.1386/dbs_00003_1}

\bibitem[Enkmann et~al.(2025)]{enkmann2025} Enkmann, J. M., Beermann, V., Mohr, P. N. C., \& Uebernickel, F. (2025). The mechanisms of attention and decision-making: A neuroscience-informed heuristic framework for the design of digital choice architectures. In F. D. Davis et al. (Eds.), \textit{Information Systems and Neuroscience (NeuroIS Retreat 2024)} (Lecture Notes in Information Systems and Organisation, Vol.\ 66). Springer. \doi{10.1007/978-3-031-71385-9_27}

\bibitem[Fillingim et~al.(2022)]{fillingim2022} Fillingim, J., Johnson, J. K., \& Bogan, C. (2022). To improve quality, leverage design. \textit{BMJ Quality \& Safety, 31}(1), 70--74. \doi{10.1136/bmjqs-2021-013665}

\bibitem[Gelman et~al.(2013)]{johnson2019} Gelman, A., Carlin, J. B., Stern, H. S., Dunson, D. B., Vehtari, A., \& Rubin, D. B. (2013). \textit{Bayesian Data Analysis} (3rd ed.). CRC Press. \doi{10.1201/b16018}

\bibitem[Gero \& Kannengiesser(2004)]{gero2004} Gero, J. S., \& Kannengiesser, U. (2004). The situated function--behaviour--structure framework. \textit{Design Studies, 25}(4), 373--391. \doi{10.1016/j.destud.2003.10.010}

\bibitem[Halmos(1960)]{numberanalyticsNDa} Halmos, P. R. (1960). \textit{Naive Set Theory}. Van Nostrand.

\bibitem[Hatchuel \& Weil(2009)]{hatchuel2009} Hatchuel, A., \& Weil, B. (2009). C--K design theory: An advanced formulation. \textit{Research in Engineering Design, 19}(4), 181--192. \doi{10.1007/s00163-008-0043-4}

\bibitem[Hermsen et~al.(2016)]{hermsen2015} Hermsen, S., Mulder, S., Renes, R. J., \& van der Lugt, R. (2016). Integrating behavioural theory into design: The Double Diamond of persuasive design. In \textit{Proceedings of DRS 2016, Design Research Society 50th Anniversary Conference}. Design Research Society. \doi{10.21606/drs.2016.326}

\bibitem[Howard(1966)]{howard1966} Howard, R. A. (1966). Decision analysis: Applied decision theory. In D. B. Hertz \& J. Melese (Eds.), \textit{Proceedings of the Fourth International Conference on Operational Research} (pp. 55--71). Wiley-Interscience.

\bibitem[K{\"o}rding(2014)]{wikipediaNDd} K{\"o}rding, K. P. (2014). A tutorial on Bayesian inference for learning and decision-making. \textit{Journal of Mathematical Psychology, 58}, 83--93. \doi{10.1016/j.jmp.2013.11.006}

\bibitem[Keeney \& Raiffa(1976)]{keeney1976} Keeney, R. L., \& Raiffa, H. (1976). \textit{Decisions with Multiple Objectives: Preferences and Value Tradeoffs}. Wiley.

\bibitem[Levinthal \& March(1993)]{levinthal1993} Levinthal, D. A., \& March, J. G. (1993). The myopia of learning. \textit{Strategic Management Journal, 14}(S2), 95--112. \doi{10.1002/smj.4250141009}

\bibitem[Levinthal(1997)]{levinthal1997} Levinthal, D. A. (1997). Adaptation on rugged landscapes. \textit{Management Science, 43}(7), 934--950. \doi{10.1287/mnsc.43.7.934}

\bibitem[Liedtka(2015)]{liedtka2018} Liedtka, J. (2015). Perspective: Linking design thinking with innovation outcomes through cognitive bias reduction. \textit{Journal of Product Innovation Management, 32}(6), 925--938. \doi{10.1111/jpim.12163}

\bibitem[March(1991)]{march1991} March, J. G. (1991). Exploration and exploitation in organizational learning. \textit{Organization Science, 2}(1), 71--87. \doi{10.1287/orsc.2.1.71}

\bibitem[Martin(2009)]{mallinson2019} Martin, R. L. (2009). \textit{The Design of Business: Why Design Thinking Is the Next Competitive Advantage}. Harvard Business Press.

\bibitem[Menning et~al.(2018)]{menning2018} Menning, A., Ewald, B., Nicolai, C., \& Weinberg, U. (2018). ``$\ldots$ and not building on that'': The relation of low coherence and creativity in design conversations. In H. Plattner, C. Meinel, \& L. Leifer (Eds.), \textit{Design Thinking Research: Making Distinctions: Collaboration Versus Cooperation} (pp. 195--213). Springer. \doi{10.1007/978-3-319-60967-6_10}

\bibitem[Moustakas \& Kyriakou(2023)]{moustakas2023} Moustakas, A., \& Kyriakou, P. (2023). Computational design methodologies for branded product forms. \textit{Machines, 11}(2), Article 232. \doi{10.3390/machines11020232}

\bibitem[Norman(2013)]{norman2013} Norman, D. A. (2013). \textit{The Design of Everyday Things} (Rev.\ and expanded ed.). Basic Books.

\bibitem[Pierri et~al.(2023)]{chen2022} Pierri, C., Biggemann, S., \& El-Kafafi, S. (2023). A systematic literature review on integrating sustainability in design thinking. \textit{Journal of Cleaner Production, 419}, Article 137885. \doi{10.1016/j.jclepro.2023.137885}

\bibitem[Plattner et~al.(2011)]{plattner2011} Plattner, H., Meinel, C., \& Leifer, L. (Eds.). (2011). \textit{Design Thinking: Understand -- Improve -- Apply} (Understanding Innovation). Springer. \doi{10.1007/978-3-642-13757-0}

\bibitem[Rittel \& Webber(1973)]{rittel1973} Rittel, H. W. J., \& Webber, M. M. (1973). Dilemmas in a general theory of planning. \textit{Policy Sciences, 4}(2), 155--169. \doi{10.1007/BF01405730}

\bibitem[Shahriari et~al.(2016)]{shahriari2016} Shahriari, B., Swersky, K., Wang, Z., Adams, R. P., \& de Freitas, N. (2016). Taking the human out of the loop: A review of Bayesian optimization. \textit{Proceedings of the IEEE, 104}(1), 148--175. \doi{10.1109/JPROC.2015.2494218}

\bibitem[Simon(1956)]{simon1956} Simon, H. A. (1956). Rational choice and the structure of the environment. \textit{Psychological Review, 63}(2), 129--138. \doi{10.1037/h0042769}

\bibitem[Simon(1996)]{simon1996} Simon, H. A. (1996). \textit{The Sciences of the Artificial} (3rd ed.). MIT Press.

\bibitem[Sohaib et~al.(2021)]{wikipediaNDc} Sohaib, M. M. A. W., Waseem, M., Lashari, S. A., \& Bham, S. U. (2021). Framework of designers' design process: A comparison between the Double Diamond and other models. \textit{Journal of Engineering, Design and Technology, 19}(5), 1177--1200. \doi{10.1108/JEDT-05-2020-0193}

\bibitem[Srinivas et~al.(2012)]{srinivas2012} Srinivas, N., Krause, A., Kakade, S. M., \& Seeger, M. W. (2012). Information-theoretic regret bounds for Gaussian process optimization in the bandit setting. \textit{IEEE Transactions on Information Theory, 58}(5), 3250--3265. \doi{10.1109/TIT.2011.2182033}

\bibitem[Stiny \& Gips(1972)]{stiny1972} Stiny, G., \& Gips, J. (1972). Shape grammars and the generative specification of painting and sculpture. In C. V. Freiman (Ed.), \textit{Information Processing 71} (pp. 1460--1465). North-Holland.

\bibitem[Suh(1990)]{suh1990} Suh, N. P. (1990). \textit{The Principles of Design}. Oxford University Press.

\bibitem[Sung \& Kelley(2019)]{sung2019} Sung, E., \& Kelley, T. R. (2019). Identifying design process patterns: A sequential analysis study of design thinking. \textit{International Journal of Technology and Design Education, 29}(2), 283--302. \doi{10.1007/s10798-018-9448-1}

\bibitem[Thomas(2021)]{thomas2021} Thomas, M. (2021). \textit{Developing a fit-for-purpose assessment framework for collaboration in the digital emerging environment} (Doctoral dissertation, Victoria University, Melbourne, Australia). Retrieved from \url{https://vuir.vu.edu.au/43467/}

\bibitem[Uebernickel et~al.(2020)]{uebernickel2020} Uebernickel, F., Jiang, L., Brenner, W., Pukall, B., Naef, T., \& Schindlholzer, B. (2020). \textit{Design Thinking: The Handbook}. World Scientific. \doi{10.1142/11329}

\bibitem[van Asselt \& Roke(2025)]{vanasselt2025} van Asselt, A., \& Roke, Y. (2025). Co-creating the Stress Autism Mate (SAM): Development of an mHealth application for autistic adults using a design thinking approach. \textit{Health Psychology and Behavioral Medicine, 13}(1), Article 2597607. \doi{10.1080/21642850.2025.2597607}

\bibitem[Wagenmakers et~al.(2018)]{laskey2023} Wagenmakers, E.-J., Marsman, M., Jamil, T., Ly, A., Verhagen, J., Love, J., \ldots\ Morey, R. D. (2018). Bayesian inference for psychology. Part I: Theoretical advantages and practical ramifications. \textit{Psychonomic Bulletin \& Review, 25}(1), 35--57. \doi{10.3758/s13423-017-1262-3}

\bibitem[Wang et~al.(2020)]{smith2025} Wang, N., Wang, D., \& Yang, Y. (2020). A comprehensive review on design optimization of structures and materials. \textit{Archives of Computational Methods in Engineering, 27}(2), 317--340. \doi{10.1007/s11831-018-9297-6}

\bibitem[Wiltse \& Stolterman(2016)]{wiltse2016} Wiltse, H., \& Stolterman, E. (2016). Design togetherness, pluralism and convergence. In \textit{Proceedings of DRS 2016, Design Research Society 50th Anniversary Conference}. Design Research Society. \doi{10.21606/drs.2016.367}

\bibitem[Woodbury \& Burrow(2006)]{woodbury2006} Woodbury, R. F., \& Burrow, A. L. (2006). Whither design space? \textit{Artificial Intelligence for Engineering Design, Analysis and Manufacturing, 20}(2), 63--82. \doi{10.1017/S0890060406060127}

\bibitem[Wuttke et~al.(2024)]{wuttke2024} Wuttke, T., Haskamp, T., Perscheid, M., \& Uebernickel, F. (2024). Building the processes behind the product: How digital ventures create business processes that support their growth. \textit{Business \& Information Systems Engineering, 66}(5), 565--583. \doi{10.1007/s12599-024-00884-5}

\bibitem[Zadeh(1965)]{wikipediaNDa} Zadeh, L. A. (1965). Fuzzy sets. \textit{Information and Control, 8}(3), 338--353. \doi{10.1016/S0019-9958(65)90241-X}

\end{thebibliography}

\end{document}

